\documentclass[11pt]{article}
\usepackage[T1]{fontenc}
\usepackage{lmodern}
\usepackage{amsmath,amssymb,amsthm,mathtools,mathrsfs}
\usepackage[margin=1in]{geometry}
\usepackage{booktabs,array,longtable,enumitem}
\usepackage{microtype}
\usepackage[colorlinks=true,linkcolor=blue,citecolor=blue,urlcolor=blue]{hyperref}
\hypersetup{pdftitle={Theta correlations, Laguerre signs, and hypercomplex jets of the theta heat flow},
 pdfauthor={Drafted for Edward Baker},pdfsubject={Complete theta correlations, reflected monogenic lifts, and stationary energy-relative signs}}
\allowdisplaybreaks
\setlength{\emergencystretch}{2em}
\numberwithin{equation}{section}
\newtheorem{theorem}{Theorem}[section]
\newtheorem{proposition}[theorem]{Proposition}
\newtheorem{lemma}[theorem]{Lemma}
\newtheorem{corollary}[theorem]{Corollary}
\newtheorem{classicalinput}[theorem]{Classical input}
\theoremstyle{definition}
\newtheorem{definition}[theorem]{Definition}
\theoremstyle{remark}
\newtheorem{remark}[theorem]{Remark}
\newcommand{\R}{\mathbb R}
\newcommand{\C}{\mathbb C}
\newcommand{\HH}{\mathbb H}
\newcommand{\OO}{\mathbb O}
\newcommand{\ii}{\mathbf i}
\newcommand{\jj}{\mathbf j}
\newcommand{\kk}{\mathbf k}
\newcommand{\dd}{\,d}
\newcommand{\Lag}{\mathcal L}
\newcommand{\RH}{\mathrm{RH}}
\newcommand{\Ephys}{E}
\DeclareMathOperator{\Rea}{Re}
\DeclareMathOperator{\Ima}{Im}
\title{Theta correlations, Laguerre signs, and\\hypercomplex jets of the theta heat flow}
\author{Drafted for Edward Baker}
\date{15 October 2026}
\begin{document}
\maketitle
\begin{abstract}
The physical first jet of the theta heat flow has energy
$E=H^2+H_x^2$. Its evolution at two stationary sets is governed by
$\Lag_0=H_x^2-HH_{xx}$ and $\Lag_1=H_{xx}^2-H_xH_{xxx}$.
We express both forms and $E$ through one complete-source correlation
family, retaining the full arithmetic interactions and the original
cutoff errors. The octonionic observations reduce to planar dot products
and determinants. A reflected monogenic extension preserves the physical
four-jet without the auxiliary correction required by a stronger
off-slice prescription. Exact calibration and logarithmic-norm identities
clarify which information these geometric descriptions retain.
We further prove a signed boundary-defect formula, sharp spatial and disk
integral bounds, centered-average positivity, and an exact finite
$E_8$-character reconstruction. A positive auxiliary theta transform
reconstructs $H$ through a signed differential combination with a modular
boundary constant. Complete Poisson localization and its heat-integrability
obstructions, a finite rational-closure obstruction, and a smooth
strongly log-concave collision control delimit these approaches.
For the theta source, both full Laguerre forms are strictly positive at
sufficiently large height when $0<t\le1/2$ and
$t\log(|x|/(4\pi))\ge61/24$, including the closed boundary; the
starting height is unevaluated. Complete-source certificates give both
stationary signs and $\sqrt E>10^{-17}$ on $0\le t\le1/2$, $|x|\le112$.
Collision expansions and analytic rigidity turn prescribed all-cutoff
lower bounds into sufficient logarithmic and subpower loss criteria.
The reviewed octonionic, Fourier-kernel, and exceptional-theta theories
identify possible tools, but do not supply the missing theta-specific
signed estimate. The inner cusp and the finite transition remain open.
\end{abstract}
\makeatletter
\renewcommand*\l@subsection{\@dottedtocline{2}{1.5em}{3em}}
\makeatother
\tableofcontents

\section{The theta heat flow and the main result}\label{sec:introduction}

The four real derivatives
\begin{equation}
 h_j(t,x)=\partial_x^jH(t,x),\qquad j=0,1,2,3,
 \label{eq:four-jets}
\end{equation}
form the smallest consecutive jet containing the collision detector
$(H,H_x)$ and both quadratic expressions
\begin{equation}
 \Lag_0=h_1^2-h_0h_2,\qquad
 \Lag_1=h_2^2-h_1h_3.
 \label{eq:laguerre-definitions}
\end{equation}
Here ``four jets'' means these four consecutive derivative components.
Their common scalar origin is essential: they are not four independent
fields. The first two give the quaternionic observation
$F=h_0+\jj h_1$ and its energy $E=h_0^2+h_1^2$. Since $H_t=-H_{xx}$,
\begin{equation}
 \frac12E_t=-h_0h_2-h_1h_3
 =\begin{cases}\Lag_0,&h_1=0,\\ \Lag_1,&h_2=0.\end{cases}
 \label{eq:stationary-energy-intro}
\end{equation}
The stationary sets in this identity are those of the original scalar
function and its first derivative. Their signs will remain the same
objects throughout the analytic and hypercomplex constructions.

\subsection{The fixed source and normalization}
The family under consideration is
\begin{align}
 \Phi(u)&=\sum_{n\ge1}\phi_n(u),\qquad
 \phi_n(u)=(2\pi^2n^4e^{9u}-3\pi n^2e^{5u})e^{-\pi n^2e^{4u}},
 \label{eq:theta-source}\\
 H(t,z)&=\int_0^\infty e^{tu^2}\Phi(u)\cos(zu)\dd u,
 \qquad H(0,z)=\frac18\xi\left(\frac12+\frac{iz}{2}\right).
 \label{eq:theta-family}
\end{align}
We use the convention
$\xi(s)=s(s-1)\pi^{-s/2}\Gamma(s/2)\zeta(s)/2$, agreeing with
\cite{polymath}. In particular the critical line corresponds to real
$z$, and the heat equation has the backward sign
\begin{equation}
 \partial_tH=-\partial_z^2H.
 \label{eq:backward-heat}
\end{equation}

For completeness, the normalization follows directly from the theta
Mellin identity. Put $A(u)=e^u\sum_{n\ge1}e^{-\pi n^2e^{4u}}$.
Poisson summation makes $P(u)=e^u+2A(u)$ even, and differentiation gives
$\Phi=(P''-P)/16=(A''-A)/8$. Splitting the theta Mellin integral at one
and substituting $y=e^{4u}$ gives
\[
 \xi(1/2+iz/2)=\frac12-(1+z^2)
                    \int_0^\infty A(u)\cos(zu)\dd u.
\]
Evenness of $P$ gives $A'(0)=-1/2$. Two integrations by parts now prove
\eqref{eq:theta-family}; all positive-infinity boundary terms vanish.
This also proves that the complete source $\Phi$ is even. For $u\ge0$
its terms are positive, and their successive majorant ratios yield
\begin{equation}
 0<\Phi(u)<4\pi^2e^{9u-\pi e^{4u}}.
 \label{eq:source-majorant}
\end{equation}
Indeed the ratio of successive terms of $n^4e^{-\pi n^2e^{4u}}$ is at
most $16e^{-3\pi}<1/2$. The bound justifies differentiation of every
order on compact sets of complex $(t,z)$. Thus $H$ is jointly entire,
even in $z$, and real on real $(t,z)$; \eqref{eq:backward-heat} follows
under the integral. For each fixed real time it is a real entire
function of order at most one, as is also immediate from the same
source bound by splitting the integral at $u=O(\log(2+|z|))$.

\subsection{Uniform positivity outside the inner region}
The leading result concerns the dimensionless parameter $t\log|x|$,
which remains nontrivial as time tends to zero while height increases.
\begin{theorem}[Uniform four-jet asymptotic]\label{thm:outer}
There is a finite, presently unevaluated $X_*>4\pi$ with the
following property. Uniformly for
\begin{equation}
 x\ge X_*,\qquad 0<t\le\tfrac12,\qquad
 \theta:=tL\ge61/24,\qquad L:=\log\frac{x}{4\pi},
 \label{eq:outer-domain}
\end{equation}
there is an explicitly defined nonzero complex amplitude $Q=Q_t(x)$
such that, with $\omega=L/4$ and $\phi=\arg Q$,
\begin{equation}
 \frac{1}{2|Q|}\left(h_0,\frac{h_1}{\omega},
                   \frac{h_2}{\omega^2},\frac{h_3}{\omega^3}\right)
 = (\cos\phi,-\sin\phi,-\cos\phi,\sin\phi)
       +O(L^{-1}).
 \label{eq:main-four-jet}
\end{equation}
Consequently
\begin{align}
 h_0^2+(h_1/\omega)^2&=4|Q|^2(1+O(L^{-1})),
 \label{eq:main-balanced}\\
 \Lag_0&=4|Q|^2\omega^2(1+O(L^{-1})),\qquad
 \Lag_1=4|Q|^2\omega^4(1+O(L^{-1})).
 \label{eq:main-signs}
\end{align}
After increasing $X_*$, both Laguerre forms are strictly
positive everywhere in \eqref{eq:outer-domain}, and $F(t,x)\ne0$ there.
Evenness gives the reflected statements at negative heights.
\end{theorem}

The theorem combines the effective approximation in
\cite[Theorem 1.3]{polymath} with five exponent-pair estimates, uniformly
in all four derivatives. Strict rational margins absorb the positive
arithmetic losses and the complex-disc crossing, so the lower boundary
$61/24$ is included exactly. The starting height remains unevaluated;
no optimality claim is made. Lemma~\ref{lem:middle-block-transfer} gives
a conditional continuation to $5/2$, while
Proposition~\ref{pg:weighted-transfer} records a weaker sufficient
weighted-remainder hypothesis. Neither hypothesis is proved here.
A different, more conservative region has a fully explicit height.
\begin{proposition}[An explicit outer region]\label{prop:explicit-outer}
Set $X_{\mathrm{exp}}=4\pi\exp(10^9)$. For $x\ge X_{\mathrm{exp}}$,
$0<t\le1/2$, and $tL\ge5$, every coordinate error in
\eqref{eq:main-four-jet} is at most $108000000/L\le0.108<1/8$.
Consequently
\begin{equation}
 \frac{h_0^2+(h_1/\omega)^2}{4|Q|^2}\ge\frac12,
 \qquad
 \frac{\Lag_j}{4|Q|^2\omega^{2j+2}}\ge\frac{31}{64},\quad j=0,1.
 \label{eq:explicit-signs-intro}
\end{equation}
Reflection gives the negative-height statements.
\end{proposition}
The complete explicit estimates are in Appendix~\ref{app:explicit-height}.
This enormous cutoff does not provide a practical bridge from height 112,
and it is not an evaluated cutoff for the sharper $61/24$ theorem.

The geometric interpretation is that a common nonzero amplitude controls
all four jets, including at phases where individual leading components
vanish. In quaternionic and octonionic coordinates, the same quadratic
forms become signed geometric observations. Those constructions do not
supply positivity by themselves; they identify precisely which part of
the theta analysis is needed. A second, complementary result is the
validated bounded cover
\begin{equation}
 |F(t,x)|>10^{-17}\quad(0\le t\le1/2,\ |x|\le112),
 \label{eq:intro-bounded}
\end{equation}
together with both strict stationary signs there. Its complete-source
proof and certificate structure are given in
Section~\ref{sec:fc-certificate}, including Theorem~\ref{fc:height112}.

Real multiple zeros at positive times are the relevant obstruction:
the classical heat-flow inputs imply
\begin{equation}
 \RH\quad\Longleftrightarrow\quad
 F(t,x)\ne0\quad\hbox{for every }t>0\hbox{ and }x\in\R.
 \label{eq:intro-rh}
\end{equation}
Section~\ref{sec:fc-collision} gives the reduction and the sufficient
stationary bounded-loss criterion. Theorem~\ref{thm:outer} and
\eqref{eq:intro-bounded} leave an inner small-time region and a finite
transition strip. The remaining problem admits several exact descriptions.
Section~\ref{sec:hypercomplex} reduces the oriented octonionic observations
to two-dimensional determinants and dot products, and gives a simpler
reflected lift of the physical trace. Section~\ref{sec:complete-correlations}
expresses both signs and the energy in one complete theta correlation
family. Its hybrid version retains the full-amplitude cutoff used in
Proposition~\ref{pg:J-transfer}. These reductions simplify the objects
to be estimated; they add no unconditional sign region.
Sections~\ref{sf:section} and~\ref{pp:section} add signed-form identities,
integrated bounds, and a positive primitive with exact reconstruction.
Section~\ref{upd:section} records the Poisson integrability obstructions,
sharper weighted budgets, and a broader conditional loss transfer.
These results preserve the distinction between auxiliary positivity and
the pointwise stationary estimate. The applicability of further octonionic
analysis is assessed in Section~\ref{subsec:octonionic-analysis-literature}.

% Prepared for Edward Baker; gpt-6-astra, effort ultra; 14 October 2026.
\section{Proof of the uniform outer-region theorem}\label{sec:outer-proof}

\subsection{A holomorphic comparison amplitude}
Use the principal logarithm and set
\begin{align}
 M_0(s)&=\frac{s(s-1)}{16}\pi^{-s/2}\sqrt{2\pi}
 \exp\left\{\left(\frac s2-\frac12\right)\operatorname{Log}\frac s2
                         -\frac s2\right\},\label{eq:M0}\\
 \alpha(s)&=\frac{M_0'(s)}{M_0(s)}
 =\frac{1}{2s}+\frac{1}{s-1}+\frac12\operatorname{Log}\frac{s}{2\pi},
 \qquad M_t(s)=e^{t\alpha(s)^2/4}M_0(s),\label{eq:alpha-Mt}\\
 s_\pm(z)&=\frac{1\pm iz}{2},\qquad
 w_\pm(z)=s_\pm(z)+\frac t2\alpha(s_\pm(z)),\label{eq:wpm}\\
 Q_\pm(z)&=M_t(s_\pm(z))\zeta(w_\pm(z)),\qquad
 \mathcal P_t(z)=Q_+(z)+Q_-(z).\label{eq:analytic-model}
\end{align}
For real $x$ we take $Q=Q_+(x)$; the other branch is its conjugate.
The functions in these definitions are analytic on the small discs
used below. Their arguments stay at imaginary parts approximately
$\pm x/2$, away from all logarithmic cuts, while $\Re w_\pm>1$ keeps
the zeta arguments away from the pole.

Write $z=u+iv$, $L_u=\log(u/(4\pi))$, and
\[
 N=\left\lfloor\sqrt{u/(4\pi)+t/16}\right\rfloor,
 \qquad S_{t,N}(w)=\sum_{n\le N}e^{t\log^2n/4}n^{-w}.
\]
Polymath's estimate applies when $u\ge200$, $0\le v\le1$, and
$0<t\le1/2$. In the present notation it reads
\begin{equation}
 H(t,z)=M_t(s_-(z))S_{t,N}(w_-(z))
       +M_t(s_+(z))S_{t,N}(w_+(z))+\mathcal E(t,z),
 \label{eq:RS-input}
\end{equation}
where $|\mathcal E|\le |M_t(s_-)|(e_A+e_B+e_{C,0})$.
This is the exact two-branch expression of the cited theorem: its first
exponent $s_*$ equals $w_-$ and its second exponent, after incorporating
the factor $n^v$, is $\overline{s_*}+\kappa-v=w_+$. To retain the full
imported error information, define
\[
 \gamma=\frac{M_t(s_+)}{M_t(s_-)},\qquad
 \kappa=\frac t2\left(\alpha\left(\frac{1-v+iu}{2}\right)
                -\alpha\left(\frac{1+v+iu}{2}\right)\right).
\]
The bounds needed from \cite[eqs. (20)--(24)]{polymath} are
\begin{align}
 |\gamma|&\le e^{0.02v}(u/(4\pi))^{-v/2},\qquad
 |\kappa|\le\frac{tv}{2(u-6)},\label{eq:RS-ratios}\\
 e_A+e_B&\le\sum_{n\le N}(1+|\gamma|N^{|\kappa|}n^v)
 e^{t\log^2n/4}n^{-\Re w_-}
 \left[\exp\left\{\frac{(t^2/16)\log^2(u/(4\pi n^2))+0.626}
                             {u-6.66}\right\}-1\right],
 \label{eq:RS-AB}\\
 e_{C,0}&\le (u/(4\pi))^{-(1+v)/4}
 \exp\left\{-\frac{tL_u^2}{16}
 +\frac{1.24(3^v+3^{-v})}{N-0.125}
 +\frac{3|L_u+i\pi/2|+10.44}{u-12}\right\}.
 \label{eq:RS-C}
\end{align}
No derivatives of this finite approximation will be taken. Its cutoff
can jump with $u$; the analytic model \eqref{eq:analytic-model} will
remove that difficulty before differentiation.

\subsection{A maximal estimate for logarithmic-phase blocks}
The gain over absolute summation comes from retaining the oscillation in
$n^{-iT}$ before estimating each arithmetic block.
\begin{lemma}[Logarithmic-phase blocks]\label{lem:logarithmic-block}
For $T\ge1$, $M\ge1$, and every interval $I$ of integers contained in
$(M,2M]$,
\begin{equation}
 \left|\sum_{n\in I}n^{-iT}\right|
 \ll T^{1/6}M^{1/2}+T^{-1/6}M.
 \label{eq:logarithmic-block}
\end{equation}
The implied constant is absolute. The same estimate holds with $T$
replaced by $-T$, and is uniform over all shortened terminal blocks.
\end{lemma}
\begin{proof}
Apply the third derivative test, in the form of
\cite[Theorem 2 and its preceding proof]{robert}, to
$f(y)=-(T/(2\pi))\log y$. On $(M,2M]$ its third derivative has
magnitude comparable to $T/M^3$, with an absolute comparison ratio.
For an interval of length $J\le M$, the test gives
\[
 J(T/M^3)^{1/6}+J^{1/2}(T/M^3)^{-1/6}
 \le T^{1/6}M^{1/2}+T^{-1/6}M.
\]
Empty intervals contribute zero; one-point intervals are absorbed by
increasing the absolute constant. Conjugation handles the opposite phase.
\end{proof}

For a continuously differentiable, possibly complex weight $W$, Abel
summation on a shortened block is the exact identity
\begin{equation}
 \sum_{M<n\le R}W(n)n^{-iT}
 =W(R)S_M(R)-\int_M^R W'(y)S_M(y)\dd y,
 \quad S_M(y)=\sum_{M<n\le y}n^{-iT},
 \quad M\le R\le2M.
 \label{eq:abel-logarithmic-block}
\end{equation}
Thus \eqref{eq:logarithmic-block} is multiplied by
$|W(R)|+\int_M^R|W'(y)|\dd y$. For a nonnegative decreasing weight,
this total-variation factor is exactly $W(M)$. The maximal form of the
lemma is needed because the last block ends at the actual cutoff $N$.

\subsection{Five exponent pairs and a closed outer boundary}
We strengthen the elementary third-derivative estimate above using five
limiting exponent pairs. Write
\begin{equation}
 \begin{aligned}
 P&=(13/84,55/84),& AP&=(13/194,76/97),\\
 A^2P&=(13/414,359/414),& D&=(18/199,593/796),\\
 R&=(4742/38463,35731/51284).
 \end{aligned}
 \label{eq:closed-pairs}
\end{equation}
Here $D$ is the name of the displayed pair, not an assertion that an
arbitrary transformation called $D$ preserves exponent pairs.
Bourgain's theorem supplies $P$ with arbitrarily small positive losses
\cite[Theorem 6]{bourgain14}; Trudgian and Yang supply $D,R$ and the
$A$-process
\cite[Section 1.1, Lemma 1.1, and eq.~(1.10)]{trudgianyang23},
\[
 A(k,l)=\left(\frac{k}{2k+2},\frac{k+l+1}{2k+2}\right).
\]
Their exponent-pair definition is uniform over arbitrary subintervals.
For $f(y)=(T/(2\pi))\log y$, its first derivative is exactly the model
monomial with exponent parameter one, so all derivative-approximation
errors are zero. Conjugation handles the opposite phase. Consequently,
for each fixed $\eta>0$, every pair $(k,l)$ in
\eqref{eq:closed-pairs} gives the complete maximal estimate
\begin{equation}
 \sup_{I\subset(M,2M]}
 \left|\sum_{n\in I}n^{-iT}\right|
 \ll_\eta T^{k+\eta}M^{l-k},
 \qquad M\ge1,\quad T\ge2\pi M.
 \label{eq:closed-maximal-block}
\end{equation}
The supremum is over integer intervals, including shortened terminal
blocks. Bourgain's proper-subinterval completion may cost a logarithm;
choose a smaller initial positive loss to absorb it into the displayed
$\eta$. No zero-loss endpoint theorem is used.

For clarity, the descendants of $P$ also have admissibly small losses.
Under equal positive increments $\epsilon$ in both coordinates, the
increments in the coordinates of $A(k,l)$ are
\[
 \frac{\epsilon}{2(1+k)(1+k+\epsilon)},\qquad
 \frac{\epsilon(1+k-l)}{2(1+k)(1+k+\epsilon)}.
\]
They are nonnegative and at most $\epsilon/2$ for the pairs used here;
a second application remains continuous. If the final coordinate
increments are at most $\eta$, their contribution to the logarithm of
the bound is at most
$\eta(\log T-\log M)+\eta\log M=\eta\log T$.
Thus \eqref{eq:closed-maximal-block} includes all losses and is a full
bound in the indicated range.

Set
\[
 \theta_*:=\frac{61}{24},\qquad \eta=\rho:=\frac1{100000}.
\]
For $(t,x)$ in \eqref{eq:outer-domain}, put
\begin{equation}
 D_x=\{z:|z-x|\le L^{-1}\},\qquad
 L_u=\log(u/(4\pi)),\qquad \theta_u=tL_u,
 \quad z=u+iv.
 \label{eq:shrinking-disc}
\end{equation}
In this proof write $N_{\rm RS}$ for the integer $N$ in
\eqref{eq:RS-input}; it will not be confused with the source cutoff
introduced in Section~\ref{sec:residuals}. Uniformly on $D_x$, the
explicit formula \eqref{eq:alpha-Mt} gives, for either branch
$w=\sigma+i\tau=w_\pm(z)$,
\begin{equation}
 \begin{gathered}
 L_u=L+O((xL)^{-1}),\quad |\tau|\asymp u,\quad
 \log N_{\rm RS}\le L_u/2+O(u^{-1}),\\
 \theta_u\ge\theta_*-\rho,\qquad
 \sigma\ge\tfrac12+\tfrac{\theta_u}{4}-\delta_x,\qquad
 0\le\delta_x=O(L^{-1})\le\rho.
 \end{gathered}
 \label{eq:exponent-bounds}
\end{equation}
These estimates hold after one fixed increase of $X_*$ and are uniform
for $0<t\le1/2$, including growing $\theta_u\le L_u/2$. In particular
$\sigma\ge9/8$ there. The real disc is allowed to cross slightly below
$\theta_*$; the following strict margins pay that displacement.

For $\beta=(\log M)/L_u$ define the unperturbed weighted exponent
\begin{equation}
 F_{k,l}(\beta;\theta)
 =k+(l-k-\tfrac12)\beta-\frac{\theta}{4}\beta(1-\beta).
 \label{eq:closed-weighted-exponent}
\end{equation}
At $\theta=\theta_*$ it is convex, with second derivative $61/48$.
The following endpoint values are exact.
\begin{center}
\small
\begin{tabular}{@{}llll@{}}
\toprule
Pair & Interval for $\beta$ & Left value & Right value\\
\midrule
$A^2P$ & $[1/5,3/10]$ & $-43/13800$ & $-869/662400$\\
$AP$ & $[3/10,3/8]$ & $-459/310400$ & $-145/198656$\\
$D$ & $[3/8,81/200]$ & $-215/407552$ & $-21807/254720000$\\
$R$ & $[81/200,3/7]$ & $-4368359/49232640000$ & $-1/1176$\\
$P$ & $[3/7,1/2]$ & $-1/1176$ & $-11/2688$\\
\bottomrule
\end{tabular}
\end{center}
The intervals cover $[1/5,1/2]$ without gaps. Convexity therefore gives,
with the corresponding pair in each interval,
\begin{equation}
 F_{k,l}(\beta;\theta_*)\le-\mu_*,\qquad
 \mu_*:=\frac{21807}{254720000}.
 \label{eq:closed-margin}
\end{equation}
The positive exponent-pair loss costs at most $\eta$, the real-part
error at most $\delta_x/2$, and the decrease of $\theta_u$ below
$\theta_*$ at most $\rho/16$. For the tiny possible extension
$\beta\le1/2+O((uL_u)^{-1})$, continue to use $P$. Its principal change
beyond $1/2$ is $\theta_u(\beta-1/2)^2/4$, which tends to zero uniformly
even when $\theta_u$ grows; the other perturbations do also. Increase
$X_*$ until this extension costs at most $\rho$. The surviving margin is
\begin{equation}
 \mu_*-\rho(1+\tfrac12+\tfrac1{16}+1)
 =\frac{76399}{1273600000}>\frac1{20000}.
 \label{eq:closed-paid-margin}
\end{equation}
Replacing $T=|\tau|$ by a constant multiple of $u$ changes bounded
factors only. All these losses are paid at the actual lower real
boundary $\theta=61/24$.

\subsection{Uniform replacement of the weighted sums at the closed boundary}
\begin{lemma}[Weighted logarithmic moments]\label{lem:weighted-moments}
Under \eqref{eq:exponent-bounds}, for every fixed nonnegative integer $r$,
\begin{equation}
 \sum_{n\le N_{\rm RS}}(\log n)^r e^{t\log^2n/4}n^{-w}
       -(-1)^r\zeta^{(r)}(w)=O_r(L^{-1}),\qquad z\in D_x.
 \label{eq:sum-replacement}
\end{equation}
For $r=0$ the factor at $n=1$ is interpreted as one. The implied
constant is uniform throughout \eqref{eq:outer-domain}.
\end{lemma}
\begin{proof}
Write $\ell=L_u$ and split at $Y=e^{\ell/5}$. For $n\le Y$, use
$e^q-1\le qe^q$ and
$e^{t\log^2n/4}\le n^{\theta_u/20}$. The small-index discrepancy is
at most
\begin{equation}
 \frac t4\sum_{n\ge2}(\log n)^{r+2}
                    n^{-(1/2+\theta_u/5-\delta_x)}.
 \label{eq:weighted-zeta-difference}
\end{equation}
Put $\theta_u=\theta_*-\rho+s$ with $s\ge0$. Since
\[
 \frac12+\frac{\theta_*}{5}=\frac{121}{120},\qquad
 \frac1{120}-\frac\rho5-\rho>\frac1{240},
\]
the exponent in \eqref{eq:weighted-zeta-difference} is at least
$1+1/240+s/5$. Hence, with the fixed finite constant
$D_r=\sum_{n\ge2}(\log n)^{r+2}n^{-1-1/240}$, that discrepancy is
bounded by
\begin{equation}
 \frac{D_r}{4\ell}(\theta_*-\rho+s)2^{-s/5}=O_r(\ell^{-1}).
 \label{eq:uniform-growing-theta}
\end{equation}
The displayed numerator is bounded on the whole half-line $s\ge0$;
this proves uniformity for growing $\theta_u$, rather than merely for
time parameters in a fixed compact substrip. Since $\sigma\ge9/8$,
the omitted ordinary zeta tail satisfies
\begin{equation}
 \sum_{n>Y}(\log n)^r n^{-\sigma}
       =O_r(\ell^r e^{-\ell/40}).
 \label{eq:ordinary-zeta-tail-three}
\end{equation}

Partition $Y<n\le N_{\rm RS}$ into the exact dyadic intervals
$(M,\min(2M,N_{\rm RS})]$, starting at $M=Y$. The weight
\[
 W_r(y)=(\log y)^r e^{t\log^2y/4}y^{-\sigma}
\]
is nonnegative and, for sufficiently large $\ell$, decreasing, because
\begin{equation}
 y(\log W_r)'=\frac r{\log y}+\frac t2\log y-\sigma
 \le\frac{5r}{\ell}-\frac12+\delta_x+O(u^{-1})<0.
 \label{eq:weighted-block-monotonicity}
\end{equation}
Apply \eqref{eq:abel-logarithmic-block} and
\eqref{eq:closed-maximal-block}, choosing the pair by the left endpoint
$\beta=(\log M)/\ell$. The weighted block is at most a constant times
$W_r(M)T^{k+\eta}M^{l-k}$. A block may cross a boundary in the table:
the estimate is still valid, because \eqref{eq:closed-maximal-block}
already controls every shortened interval at its given scale $M$.
Equations \eqref{eq:closed-margin}--\eqref{eq:closed-paid-margin}
therefore yield
\begin{equation}
 \sum_{Y<n\le N_{\rm RS}}(\log n)^r e^{t\log^2n/4}n^{-w}
       =O_r(\ell^re^{-\ell/20000}).
 \label{eq:high-index-decay}
\end{equation}
The convex dyadic summation in Proposition~\ref{upd:convex-dyadic}
pays their total with a bounded geometric multiplier, including the
shortened final block. Combining
\eqref{eq:weighted-zeta-difference},
\eqref{eq:ordinary-zeta-tail-three}, and \eqref{eq:high-index-decay}
proves the lemma.
\end{proof}

\subsection{Complete value errors and derivative control}
The positive sums in the imported errors require no further arithmetic
cancellation. For $n\le N_{\rm RS}$,
\[
 e^{t\log^2n/4}n^{-\theta_u/4}
 =\exp\{-\tfrac t4\log n(\ell-\log n)\}\ll1,
 \qquad n^{\delta_x}\ll1.
\]
On the upper half-disc \eqref{eq:RS-ratios} bounds $|\gamma|$ and
$N_{\rm RS}^{|\kappa|}$ uniformly. The factor $n^v$ replaces the first
branch's real part $(1+v)/2$ by $(1-v)/2$, already covered in
\eqref{eq:exponent-bounds}. Thus the complete positive weight sum in
\eqref{eq:RS-AB} is $O(\sum_{n\le N_{\rm RS}}n^{-1/2})=O(u^{1/4})$.
Its per-index error multiplier is $O(L^2/u)$, uniformly for $t\le1/2$.
The correction exponent in \eqref{eq:RS-C} tends uniformly to zero.
Consequently all three imported errors obey
\begin{equation}
 e_A+e_B=O(L^2u^{-3/4}),\qquad
 e_{C,0}\ll e^{-\ell/4-\theta_u\ell/16}=O(u^{-2/5}),
 \label{eq:RS-small-errors}
\end{equation}
where $1/4+(\theta_*-\rho)/16>2/5$.

Set $\Gamma=|M_t(s_+(x))|$. The logarithmic $z$-derivative of either
$M_t(s_\pm(z))$ is $O(L)$, so integration over a segment of length
at most $1/L$ gives
\begin{equation}
 |M_t(s_\pm(z))|\asymp\Gamma\qquad(z\in D_x).
 \label{eq:gamma-comparison}
\end{equation}
Combining \eqref{eq:RS-input}, Lemma~\ref{lem:weighted-moments} at
$r=0$, and \eqref{eq:RS-small-errors}--\eqref{eq:gamma-comparison}
gives
\begin{equation}
 H(t,z)-\mathcal P_t(z)=O(\Gamma/L),\qquad z\in D_x.
 \label{eq:analytic-value-error}
\end{equation}
Initially the imported approximation applies to the upper half-disc;
Schwarz symmetry supplies the lower half.

Absolute convergence of the Dirichlet series for $\zeta$ and
$1/\zeta$ in $\Re w\ge9/8$ gives
\begin{equation}
 \zeta(9/8)^{-1}\le|\zeta(w)|\le\zeta(9/8).
 \label{eq:zeta-nonzero}
\end{equation}
Thus $Q\ne0$ and $|Q|\asymp\Gamma$. The complete difference in
\eqref{eq:analytic-value-error} is holomorphic on a neighborhood of
the closed disc. Cauchy's formula, with time fixed, yields
\begin{equation}
 \partial_x^k(H(t,\cdot)-\mathcal P_t)(x)
       =O_k(|Q|L^{k-1}),\qquad k=0,1,2,3.
 \label{eq:analytic-jet-error}
\end{equation}
At $k=0$ this means $O(|Q|/L)$. The moving Riemann--Siegel floor has
been eliminated before differentiation.

\subsection{The common phase and the two signs}
At the real center the explicit carrier satisfies
\begin{equation}
 \alpha(s_+(x))=L/2+i\pi/4+O(x^{-1}),\qquad
 \alpha^{(m)}(s_+(x))=O_m(x^{-m}),\quad m\ge1.
 \label{eq:alpha-asymptotics}
\end{equation}
Fixed-order logarithmic derivatives of $\zeta$ are uniformly bounded
in $\Re w\ge9/8$. Therefore
\begin{align}
 (\log Q_+)'&=\frac i2\left[\alpha+\frac t2\alpha\alpha'
  +\left(1+\frac t2\alpha'\right)
                        \frac{\zeta'(w_+)}{\zeta(w_+)}\right]
             =i\omega+O(1),\label{eq:phase-derivative}\\
 (\log Q_+)''&=O(1),\qquad
 (\log Q_+)'''=O(1).\label{eq:higher-log-derivatives}
\end{align}
Here all $\alpha$ terms are evaluated at $s_+(x)$ and time is fixed.
The estimates remain uniform as $\theta=tL$ grows because $t\le1/2$.
The exponential derivative formulas give
\[
 Q_+^{(k)}(x)=Q\{(i\omega)^k+O(L^{k-1})\},\qquad k=1,2,3;
\]
the real-axis derivatives of $Q_-$ are their conjugates. Equation
\eqref{eq:analytic-jet-error} proves \eqref{eq:main-four-jet} with
$O(L^{-1})$, throughout the closed domain \eqref{eq:outer-domain}.

Write $P_k=h_k/(2|Q|\omega^k)$ and
$P=(c,-s,-c,s)+e$, where $c^2+s^2=1$. If $|e_k|\le d$, direct
expansion gives
\begin{equation}
 P_0^2+P_1^2\ge1-4d,\qquad
 P_1^2-P_0P_2,\ P_2^2-P_1P_3\ge1-4d-d^2.
 \label{eq:explicit-perturbation}
\end{equation}
For example the first signed form equals
$1-2se_1+e_1^2-ce_2+ce_0-e_0e_2$; retain its nonnegative square
and bound the other terms. Increasing $X_*$ until $\omega\ge1$ and
$d\le1/8$ proves Theorem~\ref{thm:outer} and the source-defined envelopes
\begin{equation}
 \begin{gathered}
 E\ge2|Q|^2\ge\frac{2\Gamma^2}{\zeta(9/8)^2},\\
 \Lag_0\ge\frac{31}{16}|Q|^2\omega^2,\qquad
 \Lag_1\ge\frac{31}{16}|Q|^2\omega^4.
 \end{gathered}
 \label{eq:closed-sign-envelopes}
\end{equation}
The endpoint $\theta=61/24$ is included because the arithmetic margins
were strictly positive before the disc and exponent-pair losses were
paid. This strengthens the earlier $8/3+\varepsilon$ result, and in
particular covers its formerly excluded boundary $8/3$. The height
$X_*$ remains unevaluated. Appendix~\ref{app:explicit-height} retains
the independent explicit height under the stronger condition $tL\ge5$.

\begin{corollary}[Asymptotic collision localization]
\label{cor:threshold-closed}\label{cor:threshold-three}\label{cor:threshold-four}
If $H(T_m,a_m)=H_x(T_m,a_m)=0$, $0<T_m\le1/2$, and
$|a_m|\to\infty$, then every sufficiently large $m$ satisfies
\[
 T_m\log\frac{|a_m|}{4\pi}<\frac{61}{24}.
\]
In particular the limsup of these products is at most $61/24$.
\end{corollary}
\begin{proof}
For $|a_m|\ge X_*$, a product at least $61/24$ contradicts the energy
lower envelope \eqref{eq:closed-sign-envelopes}, using reflection when
$a_m<0$.
\end{proof}
This localizes any hypothetical collision sequence; it does not assert
that collisions exist or impose this threshold below the unevaluated
height cutoff.

\subsection{A conditional continuation to five halves}
The largest blocks alone do not control the remaining interval:
\[
 F_P(1/2;5/2)=-\frac1{672},\qquad
 F_D(2/5;5/2)=\frac9{3980}>0.
\]
At $\beta=2/5$, $D$ is the best of the five displayed pairs. The positive
upper exponent prevents these particular estimates from proving a small
remainder there. It does not prove growth of the actual sum or failure
of a theta sign. Absolute treatment of low indices at $\theta=5/2$
requires a split $Y=e^{\alpha\ell}$ with $\alpha<1/5$, so merely
moving that split beyond the failing middle block is not a repair.

\begin{lemma}[Conditional middle-block transfer]\label{lem:middle-block-transfer}
Suppose there are fixed constants $C,\kappa>0$ and $T_0$ such that
for $T\ge T_0$, $\beta=\log M/\log T\in[7/25,47/100]$, and every
integer interval $I\subset(M,2M]$,
\begin{equation}
 \left|\sum_{n\in I}n^{-iT}\right|
 \le C T^{\beta/2+(5/8)\beta(1-\beta)-\kappa}.
 \tag{MB}\label{eq:middle-block-hypothesis}
\end{equation}
Then, for sufficiently large $x$, uniformly for $0<t\le1/2$ and the
entire closed strip $5/2\le tL\le8/3$, the four-jet asymptotic
\eqref{eq:main-four-jet} holds with $O(L^{-1})$, and
$E,\Lag_0,\Lag_1$ are strictly positive. The two cutoff conclusions of
Corollary~\ref{cor:signed-staircase-cutoffs} also hold there, with
$\zeta(11/10)$ replacing $\zeta(9/8)$ in the relative coefficient.
The hypothesis \eqref{eq:middle-block-hypothesis} is not established here.
\end{lemma}
\begin{proof}
Use the same shrinking disc and put $Y=e^{19\ell/100}$. At the real
lower boundary $\theta=5/2$, the exponent in the absolute small-index
discrepancy is
\[
 \frac12+\frac{(5/2)(1-19/100)}4=\frac{161}{160}>1.
\]
For sufficiently large height the disc perturbations leave it at least
$1+1/320$. Since $\theta$ is bounded on this strip, the discrepancy is
$O_r(L^{-1})$. Also $\Re w\ge11/10$, so the ordinary zeta tail is
$O_r(\ell^r e^{-19\ell/1000})$.

For the high blocks outside the assumed middle range, use $A^2P$ on
$[19/100,2/7]$ and $P$ on $[23/50,1/2]$. Their exact endpoint
values at $\theta=5/2$ are
\[
 \begin{array}{c|cc}
 A^2P&-3293/3312000&-1/4508\\
 P&-41/84000&-1/672.
 \end{array}
\]
Convexity gives strict savings on both full intervals, preserved by
sufficiently small positive exponent-pair losses. The middle range in
\eqref{eq:middle-block-hypothesis} overlaps these intervals by $1/175$
and $1/100$, respectively. Those fixed overlaps absorb the difference
between $\log T=\ell+O(1)$ and the scale $\ell$ used on the complex disc.

For a remaining block, write $b=\log M/\ell$, while the hypothesis uses
$\beta=\log M/\log T$. Multiplication by the actual decreasing weight
and Abel summation give, apart from fixed powers of $\ell$, a logarithm
bounded by
\[
 -\kappa\ell-\frac{\theta_u-5/2}{4}b(1-b)\ell+O(1).
\]
The $O(1)$ includes the logarithmic-coordinate change and the real-part
error, since $\delta_x\ell=O(1)$. On this fixed strip the disc's
possible displacement below $\theta=5/2$ tends to zero. Thus the bound
is at most $-\kappa\ell/2$ after increasing the starting height.
Summing the $O(\ell)$ blocks proves Lemma~\ref{lem:weighted-moments}
on the entire closed strip, conditionally on
\eqref{eq:middle-block-hypothesis}.

All imported errors remain negligible; in particular
$e_{C,0}=O(u^{-2/5})$ since $1/4+(5/2)/16=13/32>2/5$, with room
for the vanishing disc perturbation. The analytic comparison, Cauchy
estimates, and phase-uniform algebra above apply with the zero-free
half-plane $\Re w\ge11/10$. They give
$E\ge2\Gamma^2/\zeta(11/10)^2$ and the same two positive Laguerre
envelopes. The proof of Corollary~\ref{cor:signed-staircase-cutoffs}
then gives both cutoff conclusions.
\end{proof}
Together with Theorem~\ref{thm:outer}, the hypothesis would cover all
$tL\ge5/2$ in the allowed time range at sufficiently large height.
It is a sufficient arithmetic condition, not a necessary condition or
an RH-equivalent reformulation. A direct estimate for the complete
weighted or signed physical observation could bypass it. Neither this
conditional transfer nor the unconditional theorem settles the inner
cusp or evaluates the finite transition to the asymptotic region.

\section{Why positive-time collisions are the obstruction}
\label{sec:fc-collision}

The preceding asymptotic result is useful because a failure of RH forces
a finite collision of real zeros at a positive time.  We record the
precise reduction, including the compactness input, before formulating
the signed estimate still needed outside the proved regions.

We use the normalization
\[
 H(0,z)=\frac18\xi\!\left(\frac12+\frac{iz}{2}\right),
 \qquad H_t=-H_{zz}.
\]
The positive, rapidly decreasing theta source makes $H$ jointly entire
in $(t,z)$, even in $z$, and real on the real $(t,z)$ plane.  Moreover,
$H(t,iy)>0$ for real $t,y$, so no time slice is identically zero and
every zero has finite multiplicity.

The classical de Bruijn--Newman theory gives a finite real threshold
$\Lambda$ such that all zeros of $H(t,\cdot)$ are real exactly when
$t\ge\Lambda$; the endpoint is included
\cite{debruijn,newman}.  In this normalization $\Lambda\le1/2$ and
RH is equivalent to $\Lambda\le0$.  De Bruijn's strip contraction
keeps all zeros in $|\operatorname{Im}z|\le1$ for $t\ge0$.
We also use the uniform confinement theorem of
\cite[Theorem 1.5(i)]{polymath}: there is an absolute constant $C>0$
such that, for $0<t\le1/2$, a nonreal zero satisfies
\begin{equation}
 |\operatorname{Re}z|<e^{C/t}.
 \label{fc:uniform-confinement}
\end{equation}
Uniformity in time, rather than a separate uncontrolled radius for
each slice, is what prevents a first collision from escaping to infinity.

\begin{lemma}[Local splitting at a multiple zero]
\label{fc:hermite}
Let $W$ be real analytic near $(T,a)$, solve $W_t=-W_{xx}$, and have
an exact spatial zero of multiplicity $m\ge2$ there.  Its nearby zeros
have leading locations
\[
 a+\sqrt{2(t-T)}\,\lambda_k+O(|t-T|),\qquad k=1,\ldots,m,
\]
where the $\lambda_k$ are the distinct real roots of the probabilists'
Hermite polynomial $\operatorname{He}_m$.  In particular, $W$ has a
nonreal zero at times immediately before $T$.
\end{lemma}
\begin{proof}
Put $W_q=\partial_x^qW(T,a)$.  The convergent heat Taylor series is
\[
 W(T+\delta,a+w)=
 \sum_{p,q\ge0}\frac{(-\delta)^p w^q}{p!q!}W_{2p+q}.
\]
With $\delta=\eta^2$ and $w=\sqrt2\eta\lambda$, division by $\eta^m$
extends analytically to $\eta=0$, where its value is
$2^{m/2}W_m\operatorname{He}_m(\lambda)/m!$.
The roots of $\operatorname{He}_m$ are real and simple: its Rodrigues
formula gives orthogonality to all lower-degree polynomials for the
positive Gaussian weight, and the polynomial formed from its real
sign-changing factors would contradict that orthogonality unless there
were $m$ such roots. The simple roots therefore continue by the implicit function
theorem.  For $t<T$, take $\eta=i\sqrt{T-t}$; at least one Hermite root
is nonzero, so its leading imaginary displacement cannot be cancelled
by the $O(T-t)$ remainder.
\end{proof}

\begin{proposition}[Finite collision criterion for RH]
\label{fc:rh-collision}
RH holds if and only if
\begin{equation}
 H(t,x)=H_x(t,x)=0
 \quad\hbox{has no solution with }0<t\le1/2,\ x\in\R.
 \label{fc:collision-criterion}
\end{equation}
Equivalently, $E(t,x)=h_0(t,x)^2+h_1(t,x)^2$ is positive throughout
this region.  No uniform positive lower bound for $E$ is required. Equivalently,
there is no such joint zero at any positive time.
\end{proposition}
\begin{proof}
If RH holds, all zeros are real at all nonnegative times.  A multiple
zero at a positive time would contradict Lemma~\ref{fc:hermite}.

Conversely, if RH fails then $0<\Lambda\le1/2$.  Choose
$t_k\uparrow\Lambda$ with $t_k\ge\Lambda/2$, and a nonreal zero
$z_k$ at each $t_k$.  The strip bound and
\eqref{fc:uniform-confinement} put these zeros in the fixed compact set
$|\operatorname{Im}z|\le1$, $|\operatorname{Re}z|\le e^{2C/\Lambda}$.
A subsequence converges to a zero $a$ of $H(\Lambda,\cdot)$, which is
real by the threshold theorem.  If that zero were simple, the local
implicit zero branch would be unique; conjugation symmetry would make
it real for real times on both sides of $\Lambda$.  This contradicts
$z_k\notin\R$.  Thus $H(\Lambda,a)=H_x(\Lambda,a)=0$.
The source positivity gives $H(\Lambda,0)>0$, so $a\ne0$.
\end{proof}

\subsection{A sufficient estimate only on two stationary sets}

Recall
\[
 \mathcal L_j=h_{j+1}^2-h_jh_{j+2},\qquad j=0,1.
\]
On a smooth stationary curve $h_{j+1}(t,\chi(t))=0$, the heat equation
gives
\[
 \frac12\frac{d}{dt}h_j(t,\chi(t))^2=\mathcal L_j(t,\chi(t)).
\]
The following result includes singular stationary sets and all finite
multiplicities.  Its losses may grow with height or as $t\downarrow0$;
only local boundedness through each possible positive-time collision
is required.

\begin{lemma}[A stationary predecessor of an even zero]
\label{fc:predecessor}
Suppose $W_t=-W_{xx}$ and $W$ has exact spatial multiplicity $2r$ at
$(T,a)$, where $r\ge1$.  With $W_q=\partial_x^qW(T,a)$, there is a
real analytic function $\chi(h)$ for small $h\ge0$ such that
\begin{align*}
 W_x(T-h,\chi(h))&=0,&
 \chi(h)&=a-\frac{W_{2r+1}}{rW_{2r}}h+O(h^2),\\
 W(T-h,\chi(h))&=\frac{W_{2r}}{r!}h^r+O(h^{r+1}),&
 W_{xx}(T-h,\chi(h))&=\frac{W_{2r}}{(r-1)!}h^{r-1}+O(h^r).
\end{align*}
\end{lemma}
\begin{proof}
Use the heat Taylor series above with $\delta=-h$ and $w=hy$.
The function $h^{-r}W_x(T-h,a+hy)$ extends analytically to $h=0$,
where it equals
$W_{2r}y/(r-1)!+W_{2r+1}/r!$.
Its $y$ derivative is nonzero.  The implicit function theorem gives
$y(h)$, and substitution in the same series gives the two displayed
asymptotics.
\end{proof}

\begin{theorem}[Stationary bounded-loss criterion]
\label{fc:bounded-loss}
Suppose there are nonnegative functions $K_0,K_1$, locally bounded
on $(0,1/2]\times\R$, for which
\begin{equation}
 \begin{split}
 h_1=0&\quad\Longrightarrow\quad\mathcal L_0\ge-K_0E,\\
 h_2=0&\quad\Longrightarrow\quad\mathcal L_1\ge-K_1E.
 \end{split}
 \label{fc:stationary-loss}
\end{equation}
Then RH holds.  Local boundedness here includes a neighborhood of
every candidate joint zero, using a left neighborhood in time at
$t=1/2$.
\end{theorem}
\begin{proof}
Suppose $H$ has a real multiple zero at $(T,a)$ in this time range,
of exact multiplicity $m\ge2$.
If $m=2r$, apply Lemma~\ref{fc:predecessor} to $W=H$.
On its predecessor branch $h_1=0$, $E=H^2>0$, and
\[
 -\frac{\mathcal L_0}{E}=\frac{H_{xx}}H
 =\frac rh+O(1),\qquad h=T-t\downarrow0.
\]
This contradicts the first inequality with locally bounded $K_0$.
If $m=2r+1$, apply the lemma to $W=H_x$.
Then $h_2=0$, $H_x=c h^r+O(h^{r+1})$, and
$H_{xxx}=rc h^{r-1}+O(h^r)$, with $c\ne0$.
The heat Taylor series also gives $H=O(h^{r+1})$ on this branch,
because its displacement from $a$ is $O(h)$.  Consequently
\[
 E=H_x^2(1+O(h^2)),\qquad
 -\frac{\mathcal L_1}{E}=\frac rh+O(1),
\]
contradicting the second inequality.  All divisions occur on proved
nonzero predecessor branches, not at the supposed collision.
Proposition~\ref{fc:rh-collision} now applies.
\end{proof}

Both levels matter: the heat polynomial
$H(t,x)=x^3-6(t-T)x$ has
$\mathcal L_0=3x^4+36(t-T)^2\ge0$ everywhere, but has a triple zero
at $(T,0)$.  Its second stationary loss diverges on approaching that
collision from earlier times.  Thus positivity of the first observation
alone is insufficient for arbitrary multiplicity.



% Prepared for Edward Baker, 14 October 2026.
% Substantial LLM assistance: OpenAI gpt-6-astra, effort ultra,
% as verified and supplied by the parent session for this turn.
% Standalone section fragment. Required macros: \R,\C,\HH,\OO,\ii,\jj,\kk.
% Required citation keys: sudbery, baez, kauhanenorelma.

\subsection{The finite part of the collision pole and delayed cancellation}
\label{subsec:collision-finite-part}

The preceding collision argument uses only the leading pole.  Its next
coefficient also identifies cancellations that a delayed source estimate
must preserve.  All statements in this subsection are conditional on a
finite collision; they do not assert that the theta family has one.

\begin{proposition}[Finite part on the stationary predecessor]
\label{prop:collision-finite-part}
Suppose $H(T,\cdot)$ has exact multiplicity $m=2r+j$ at $x_*$, where
$T>0$, $r\ge1$, and $j\in\{0,1\}$. Put
\[
 \begin{gathered}
 A=h_m(T,x_*)\ne0,\qquad B=h_{m+1}(T,x_*),\qquad
 D=h_{m+2}(T,x_*),\\
 k=A/r!,\qquad \alpha_* =\frac{D}{(r+1)A}-\frac{B^2}{2rA^2}.
 \end{gathered}
\]
On the predecessor $h_{j+1}(T-s,\chi(s))=0$ of
Lemma~\ref{fc:predecessor}, as $s\downarrow0$,
\begin{align}
 \chi(s)&=x_*-\frac{B}{rA}s+O(s^2),\notag\\
 h_j(T-s,\chi(s))&=k s^r(1+\alpha_*s+O(s^2)),\notag\\
 \frac{\Lag_j}{E}&=-\frac rs-\alpha_*+O(s),
 \qquad E\sim k^2s^{2r},\qquad
 \Lag_j\sim-rk^2s^{2r-1}.
 \label{eq:collision-finite-part}
\end{align}
For odd multiplicity, the more precise comparison is
\begin{equation}
 \frac{H}{H_x}=-\frac{B}{r(r+1)A}s+O(s^2),\qquad
 E=h_1^2\left(1+\frac{B^2s^2}{r^2(r+1)^2A^2}+O(s^3)\right).
 \label{eq:odd-collision-energy}
\end{equation}
\end{proposition}
\begin{proof}
For $W=h_j$, the heat Taylor expansion at the collision is
\[
 W(T-s,x_*+y)=\sum_{p,q\ge0}
 \frac{s^py^q}{p!q!}\,h_{j+2p+q}(T,x_*).
\]
The implicit-function construction in Lemma~\ref{fc:predecessor}
gives $y=-Bs/(rA)+O(s^2)$.  The terms of order $s^{r+1}$ in $W$
are $Ds^{r+1}/(r+1)!$, $Bs^ry/r!$, and
$As^{r-1}y^2/[2(r-1)!]$.  Their sum, divided by $ks^r$,
is $\alpha_*s$, proving the second expansion.  Along the branch,
$dW(T-s,\chi(s))/ds=W_{xx}$ exactly, so
$W_{xx}/W=r/s+\alpha_*+O(s)$.  For $j=0$, $E=W^2$ and the result
follows.  For $j=1$, the same series for $H$ gives
\[
 H(T-s,\chi(s))=-\frac{B}{r(r+1)r!}s^{r+1}+O(s^{r+2}),
\]
which proves \eqref{eq:odd-collision-energy}.  Since
$\Lag_1=-W W_{xx}$ on this branch, the relative $O(s^2)$ energy
correction changes the displayed quotient only by $O(s)$.
\end{proof}

The moving energy derivative must be distinguished from the fixed-position
one.  With $\xi(t)=\chi(T-t)$,
\[
 \frac d{dt}E(t,\xi(t))=
 \begin{cases}2\Lag_0,&h_1=0,\\
 2\Lag_1+2\xi' h_0h_1,&h_2=0.
 \end{cases}
\]
The extra term divided by $E$ is $O(s)$ on the odd predecessor.
The exact branch identity without that term is
$\frac12\frac d{dt}h_1^2=\Lag_1$.
For odd multiplicity, the surviving real zero has displacement
$-Bs/((r+1)A)+O(s^2)$ from $x_*$: substituting $y=sz$ in the heat
series for $H$ and applying the implicit function theorem gives this
coefficient.  Its displacement from the stationary predecessor is
$Bs/[r(r+1)A]+O(s^2)$, consistently with
\eqref{eq:odd-collision-energy}.

Fix $a=\pi/16$, $b=1/4$, and $\mu=a^2+b^2$, and retain
$C_j=-h_{j+2}-2ah_{j+1}-\mu h_j$.  For any declared delay put
\[
 G_j=C_j(t-\delta,x),\qquad P_j=h_jG_j,\qquad
 Q_j=h_j(C_j-G_j/2).
\]
At current stationarity the exact combination is
\begin{equation}
 \frac{P_j/2+Q_j}{E}
 =\frac{\Lag_j}{E}-\mu\frac{h_j^2}{E}
 =-\frac rs-(\alpha_*+\mu)+O(s).
 \label{eq:delay-collision-combination}
\end{equation}
Thus changing the delay changes the two summands but not the observation
whose adverse part must be controlled.

\begin{proposition}[Stronger poles introduced by a positive delay]
\label{prop:positive-delay-cancellation}
Suppose the prescribed delay is analytic near the candidate event,
$\delta_*:=\delta(T,x_*)>0$, and put
$g=C_j(T-\delta_*,x_*)$.  If $g\ne0$, then on the predecessor
\[
 \frac{P_j}{E}=\frac gk s^{-r}(1+O(s)),\qquad
 \frac{Q_j}{E}=-\frac{P_j}{2E}-\frac rs-(\alpha_*+\mu)+O(s).
\]
For $r>1$, the two products therefore have opposite leading poles of
order $s^{-r}$.  Their intervening negative powers also cancel in
\eqref{eq:delay-collision-combination}.  If
$A_{\min}=(-P_j/E)_+$ and $B_{\min}=(-Q_j/E)_+$ are used only as
diagnostic quantities on this nonzero branch, then for $r>1$ and $g\ne0$
\[
 A_{\min}/2+B_{\min}\sim\frac{|g/k|}{2}s^{-r}.
\]
Both split ratios are $O(s^{-1})$ if and only if
$G_j=O(s^{r-1})$.  For $r=1$, writing $\gamma=g/k$ (including
$g=0$),
\[
 P_j/E=\gamma/s+O(1),\qquad
 Q_j/E=-(1+\gamma/2)/s+O(1).
\]
The necessary leading conditions for $P_j\ge0$ and
$Q_j\ge-KE$ with locally bounded $K$ are respectively
$\gamma\ge0$ and $\gamma\le-2$, and are incompatible.
\end{proposition}
\begin{proof}
Analyticity gives $G_j=g+O(s)$, while
$h_j=ks^r(1+O(s))$ and $E=h_j^2(1+O(s^2))$; the even case has
the exact equality $E=h_j^2$.  These facts prove the first expansion,
and the second follows from \eqref{eq:delay-collision-combination}.
Taking the negative part according to the sign of $g/k$ proves the
diagnostic formula when $r>1$.  Finally
$P_j/E=G_j/[ks^r](1+O(s))$ proves the equivalence; the combined
identity gives the same bound for $Q_j/E$.  The case $r=1$ follows
without assuming $g\ne0$.
\end{proof}

For comparison, a diagnostic delay $\delta=\rho s+O(s^2)$ with fixed
$\rho\ge0$ gives, from the leading heat polynomial,
\[
 \frac{P_j}{E}=-\frac{r(1+\rho)^{r-1}}s+O(1),\qquad
 \frac{Q_j}{E}=
 \frac{r\{(1+\rho)^{r-1}/2-1\}}s+O(1).
\]
If only $\delta/s\to\rho$ is known, replace $O(1)$ by $o(s^{-1})$.
These rules refer to an unknown collision time and are stress tests,
not proposed source protocols.  A split can even create an adverse pole
on a favorable branch: for
$H=y^4-12uy^2+12u^2$, $u=t-T>0$, at $y=0$,
\[
 E=144u^4,\quad\Lag_0=288u^3>0,\quad
 \frac{P_0}{E}=-\frac{2\delta+\mu\delta^2}{u^2}
                 +\frac{2+2\mu\delta}{u}-\mu
\]
for a fixed $0<\delta<T$, whereas
$P_0/(2E)+Q_0/E=2/u-\mu$.  This is a polynomial control outside
the theta class.

\subsection{Logarithmic losses and local analytic rigidity}
\label{subsec:logarithmic-loss}

For prescribed $M>0$ and finite $q\ge0$, define
\[
 \mathfrak f_q(e)=e[1+\log_+(M/e)]^q\quad(e>0),\qquad
 \mathfrak f_q(0)=0,\qquad\log_+v=\max\{0,\log v\}.
\]
\begin{proposition}[A logarithmic stationary-loss criterion]
\label{prop:logarithmic-collision-criterion}
Suppose on the two stationary sets that
$\Lag_j\ge-K_j\mathfrak f_q(E)$, with nonnegative $K_j$ locally
bounded through every candidate positive-time collision.  Then no such
collision occurs in a domain containing the required predecessor
neighborhoods.  If the domain covers $(0,1/2]\times\R$, RH follows.
\end{proposition}
\begin{proof}
On the predecessor of a zero of multiplicity $2r+j$,
$E\sim k^2s^{2r}$ and $-\Lag_j/E=r/s+O(1)$.  Hence
\[
 \frac{-\Lag_j}{\mathfrak f_q(E)}
 =\frac{r(1+o(1))}{s[2r\log(1/s)]^q}\longrightarrow\infty,
\]
contradicting the local bound for $K_j$.  Nontriviality of the theta
family ensures finite spatial multiplicity.  All divisions are on the
nonzero predecessor, and Proposition~\ref{fc:rh-collision} supplies
the final RH implication.
\end{proof}

A loss by a fixed smaller energy power does not have the same
all-multiplicity implication: $-\Lag_j/E^p$ has order
$s^{2r(1-p)-1}$, which diverges only if $p>1-1/(2r)$.
At equality its limit is $r|k|^{1/r}$, and below equality it tends
to zero.  Also, a loss coefficient integrable in time along the
predecessor would suffice to exclude the event: any coefficient allowing
it must dominate $r/s+O(1)$, whose integral diverges.  This is a
branchwise requirement; two-dimensional local integrability alone gives
no such control on the stationary curve.
Analyticity is essential for the logarithmic statement when $q>1$.
With $\beta=1/(q-1)$, the smooth function
$y(t)=\exp[-(T-t)^{-\beta}]$, continued by zero for $t\ge T$,
satisfies $y'=-\beta y[\log(1/y)]^q$ and reaches zero.  It is not
analytic at $T$ and is not a theta or backward-heat example.

The logarithmic allowance is useful quantitatively, but on analytic
stationary germs it does not enlarge the qualitative class of locally
bounded losses.
\begin{lemma}[Analytic rigidity of logarithmic losses]
\label{lem:analytic-log-rigidity}
Let $S$ be semianalytic near a finite point $p$, and let $E\ge0$ and
$L$ be real analytic in an ambient neighborhood.  Suppose, on $S$ near
$p$,
\[
 L\ge-K E[1+\log_+(M/E)]^q,
\]
where $K\ge0$ is bounded, $M>0$ and $q\ge0$ are fixed and finite,
and the right side is zero at $E=0$.  On a smaller neighborhood in $S$
there is a finite $C\ge0$ such that $L\ge-CE$.  Conversely this
linear bound implies the logarithmic bound with $K=C$.
\end{lemma}
\begin{proof}
If no such $C$ exists, there are $z_n\in S$ tending to $p$ with
$E(z_n)>0$, $-L(z_n)>0$, and $E(z_n)/[-L(z_n)]\to0$.
The semianalytic set
\[
 Z=\{(z,u):z\in S,\ E(z)>0,\ -L(z)>0,\ u>0,
                     \ E(z)=u[-L(z)]\}
\]
therefore has $(p,0)$ in its closure.  Real analytic curve selection
\cite[(7.2), p.~40]{ruiz} supplies an analytic arc
$(z(\rho),u(\rho))$ into $Z$ for $\rho>0$.  Boundedness of $L$
and $E=u(-L)$ imply $E(p)=0$.  Since $u(0)=0$ and
$E(z(0))=0$, but both functions are positive for $\rho>0$, they have
finite orders
$u(\rho)=a\rho^k(1+o(1))$ and
$E(z(\rho))=b\rho^m(1+o(1))$, with $a,b>0$ and integers
$k,m\ge1$.  The assumed estimate would force
\[
 u(\rho)^{-1}\le K[1+\log_+(M/E(z(\rho)))]^q
                  =O((1+|\log\rho|)^q),
\]
contradicting its order $\rho^{-k}$.  At $E=0$ the assumed estimate
already gives $L\ge0$.  The converse is immediate.
\end{proof}

The theta stationary sets are analytic, including their singular points.
Restrictions by $t\le1/2$ or by a closed $t\log(x/(4\pi))$ strip
are locally semianalytic for $x>4\pi$, so the lemma applies there as
well.  It yields an existence constant on a germ, not a source formula
for that constant, nor uniform control at infinity or as $t\downarrow0$.
In particular it does not authorize defining a proof coefficient from
an observed quotient $-\Lag_j/E$.
The exact record
\path{numerics/COLLISION_BEHAVIOR_CHECKS_20261014.json} contains
1072 named assertions over 60 rational jet cases of multiplicities
2 through 13, including 30 delay comparisons.  Its driver
\path{numerics/check_collision_behavior.py} independently forms the
physical numerators and denominators and checks the finite-part
coefficients.  These finite polynomial checks support the algebra;
they do not prove analytic branch existence, the logarithmic limit, or
a theta-specific sign estimate.  Review 20 records the internal audit.

% Prepared for Edward Baker, 15 October 2026.
% Substantial LLM assistance: recorded model gpt-6-astra, effort ultra,
% as verified and supplied by the parent turn metadata.
\section{The four jets as quaternionic and octonionic fields}
\label{sec:hypercomplex}

The physical geometry is planar.  The vectors
\(v_j=(h_j,h_{j+1})\) satisfy
\[
 E=|v_0|^2,\qquad \Lag_j=-\det(v_j,v_{j+1}),\qquad
 \tfrac12E_t=-v_0\cdot v_2.
\]
Quaternionic fields extend these overlapping pairs away from the real
slice; octonionic associators preserve their signed dot products and
determinants.  We first give this scalar dictionary, then construct two
monogenic extensions with different off-slice traces.  The constructions
organize the observations but impose no additional theta inequality.

\subsection{A quaternionic heat field with independent external time}

Use Hamilton units \(\ii\jj=\kk\),
\(\ii^2=\jj^2=\kk^2=-1\), and identify
\(\C\) with \(\C_{\ii}=\R+\ii\R\).
Complex coefficients are placed on the right of \(\jj\). Thus
\begin{equation}
 (a+\jj b)(c+\jj d)
 =ac-\overline b d+\jj(bc+\overline a d),
 \qquad \overline{a+\jj b}=\overline a-\jj b,
 \label{eq:hc-quaternion-product}
\end{equation}
for \(a,b,c,d\in\C_{\ii}\), and
\(|a+\jj b|^2=|a|^2+|b|^2\).

Write the independent quaternion coordinates as
\[
 q=x_0+\ii x_1+\jj x_2+\kk x_3=z_0+\jj z_1,
 \qquad z_0=x_0+\ii x_1,\quad z_1=x_2-\ii x_3.
\]
The unnormalized left Fueter operator and its conjugate are
\[
 D_q=\partial_{x_0}+\ii\partial_{x_1}
                  +\jj\partial_{x_2}+\kk\partial_{x_3},
 \qquad
 \overline D_q=\partial_{x_0}-\ii\partial_{x_1}
                  -\jj\partial_{x_2}-\kk\partial_{x_3}.
\]
A differentiable field is \emph{left Fueter regular} if \(D_qF=0\).
The convention using the factor \(1/2\) has the same kernel; all units
here multiply on the left. For the classical theory see \cite{sudbery}.

Let \(H(\tau,z)\) be the jointly entire theta heat family, satisfying
\(H_\tau=-H_{zz}\). With \(t\in\R\) an independent external parameter,
set
\[
 \tau=t+\overline{z_1}=t+x_2+\ii x_3.
\]
\begin{equation}
 G_m(t,q)=\partial_z^mH(\tau,z_0)
                +\jj\partial_z^{m+1}H(\tau,z_0),
 \qquad F=G_0.
 \label{eq:hc-lift}
\end{equation}
The coordinate \(z_1\) supplies a complex heat-time displacement; both
entries of \(G_m\) have the same spectral argument.

\begin{proposition}[Regular lift and jet identities]
\label{prop:hc-fueter}
Each \(G_m\) is globally real analytic and left Fueter regular in \(q\).
For every \(m\ge0\),
\begin{align}
 D_qG_m&=0,&\overline D_qG_m&=2G_{m+1},&
 \partial_tG_m&=-G_{m+2}.\label{eq:hc-jet-operators}
\end{align}
\begin{equation}
 \begin{aligned}
 \partial_{x_0}G_m&=G_{m+1},&\qquad
 \partial_{x_1}G_m&=G_{m+1}\ii,\\
 \partial_{x_2}G_m&=-G_{m+2},&
 \partial_{x_3}G_m&=-G_{m+2}\ii.
 \end{aligned}
 \label{eq:hc-coordinate-jets}
\end{equation}
\end{proposition}
\begin{proof}
For complex-valued \(U,V\), direct use of
\eqref{eq:hc-quaternion-product} gives
\[
 \tfrac12D_q(U+\jj V)
 =U_{\overline z_0}-V_{z_1}
              +\jj\bigl(V_{z_0}+U_{\overline z_1}\bigr).
\]
The derivatives are componentwise Wirtinger derivatives. Taking
\(U=\partial_z^mH(\tau,z_0)\) and
\(V=\partial_z^{m+1}H(\tau,z_0)\), the first two terms vanish by
holomorphy and the final pair cancel by the heat equation. Joint
entireness proves global real analyticity. Differentiating the declared
arguments proves \eqref{eq:hc-coordinate-jets} and the external time
identity. Finally \(D_q+\overline D_q=2\partial_{x_0}\) proves the
remaining assertion.
\end{proof}

\subsection{The planar observations and their exact identities}

On the physical slice \(q=x\in\R\), write
\(h_m=h_m(t,x)=\partial_x^mH(t,x)\). The jets are real and
\begin{equation}
 F=h_0+\jj h_1,\qquad F_x=G_1=h_1+\jj h_2,
 \qquad F_{xx}=G_2=h_2+\jj h_3=-F_t.
 \label{eq:hc-physical-jets}
\end{equation}
Thus the four physical jets occur in three consecutive overlapping
quaternion pairs. The basic observations are
\begin{equation}
 E=|F|^2=h_0^2+h_1^2,\qquad
 \mathcal L_j=h_{j+1}^2-h_jh_{j+2}\quad(j=0,1).
 \label{eq:hc-observations}
\end{equation}
In particular \(F(t,x)=0\) means exactly that \(H(t,\cdot)\) has a multiple
real zero at \(x\). The scalar positive-time collision criterion therefore
translates to
\[
 \mathrm{RH}\quad\Longleftrightarrow\quad
 F(t,x)\ne0\quad(t>0,\ x\in\R).
\]
This uses the scalar collision theorem for the theta family. It makes
no assertion about zeros at nonreal quaternion arguments.

The signed products and the heat equation give, on this slice,
\begin{align}
 \overline{G_j}G_{j+1}
   &=h_jh_{j+1}+h_{j+1}h_{j+2}-\jj\mathcal L_j,
   \label{eq:hc-signed-product}\\
 E_t&=-2(h_0h_2+h_1h_3),&
 (\partial_t+\partial_x^2)E&=2(h_1^2+h_2^2).
 \label{eq:hc-energy-evolution}
\end{align}
Consequently
\begin{equation}
 h_1=0\ \Longrightarrow\ E_t=2\mathcal L_0,
 \qquad
 h_2=0\ \Longrightarrow\ E_t=2\mathcal L_1.
 \label{eq:hc-stationary-energy}
\end{equation}
These formulas preserve the two different stationary conditions. The
positive source term in the backward-heat identity
\eqref{eq:hc-energy-evolution} does not give a forward minimum principle.

For \(a=A+\jj B\) and \(b=C+\jj D\) in this physical plane, put
\[
 a\cdot b=AC+BD,\qquad a\wedge b=AD-BC.
\]
Then \(\bar a b=a\cdot b+\jj(a\wedge b)\), and the planar Gram
identity gives
\begin{equation}
 |a|^2|b|^2=(a\cdot b)^2+(a\wedge b)^2.
 \label{eq:hc-planar-gram}
\end{equation}
In particular, for any physical residual \(R\in\R+\jj\R\),
\begin{equation}
 E|F_x|^2=\tfrac14E_x^2+\Lag_0^2,\qquad
 E|R|^2=(F\cdot R)^2+(F\wedge R)^2.
 \label{eq:hc-radial-transverse}
\end{equation}
A lower bound for \(F\cdot R\) need not control the transverse
component \(F\wedge R\).  This distinction will matter when a full
residual-norm bound is replaced by a one-sided signed estimate.

There is also a quotient-free relation among the consecutive pairs.
Writing \(\mathfrak d=h_0h_3-h_1h_2\), the two-by-three determinant
identity gives
\begin{equation}
 \Lag_0F_t=\Lag_1F+\mathfrak d F_x.
 \label{eq:hc-plucker}
\end{equation}
Indeed, its real and \(\jj\) components are respectively
\(-\Lag_0h_2=\Lag_1h_0+\mathfrak d h_1\) and
\(-\Lag_0h_3=\Lag_1h_1+\mathfrak d h_2\).  These polynomial
identities remain valid when a jet or a Laguerre form vanishes.

On any interval where \(h_j+ih_{j+1}\ne0\), the same orientation has
the scalar phase description
\begin{equation}
 \partial_x\arg(h_j+ih_{j+1})
 =-\frac{\Lag_j}{h_j^2+h_{j+1}^2}.
 \label{eq:hc-phase}
\end{equation}
This formula describes phase away from a common zero; the polynomial
observations above remain the appropriate quantities at a candidate
collision.

The quadratic form itself has only three real coordinates, with
\begin{equation}
 \Lag_j=h_{j+1}^2+\tfrac14(h_j-h_{j+2})^2
                         -\tfrac14(h_j+h_{j+2})^2.
 \label{eq:hc-inertia}
\end{equation}
Its inertia is \((2,1)\).  Each frequency kernel
\((y+z)^2\) and \((y-z)^2\), considered separately, uses only the
features \(1,y,y^2\), since \((y\pm z)^2=y^2+z^2\pm2yz\).
Each corresponding finite frequency matrix therefore has rank at most
three; no such rank bound is asserted here for the combined weighted
Hermitian and bilinear expression.
Shared moments can therefore be evaluated without constructing every
pair separately.  An invertible change of coordinates preserves the
negative direction; low rank does not give positivity.

\subsection{Two octonionic extensions and their different traces}

For a reader unfamiliar with octonions, it suffices here to regard
\(\OO\) as two copies of \(\HH\), with a specified multiplication.
Write \((a,b)=a+\ell b\), where \(\ell=(0,1)\), and adopt the
Cayley--Dickson convention
\begin{equation}
 (a,b)(c,d)=(ac-d\overline b,\overline a d+cb),
 \qquad \overline{(a,b)}=(\overline a,-b).
 \label{eq:hc-octonion-product}
\end{equation}
Then \(|(a,b)|^2=|a|^2+|b|^2\). This is the convention in
\cite[Section 2.2]{baez}. Octonionic multiplication is nonassociative;
we keep parentheses whenever three octonionic factors occur. All
component products on the right of \eqref{eq:hc-octonion-product} are
associative quaternion products.

Let \(r=y_0+\ii y_1+\jj y_2+\kk y_3\), so that
\(p=q+\ell r\) is the eight-dimensional argument. Put
\(e_0=1,e_1=\ii,e_2=\jj,e_3=\kk\), and define the left operator
\begin{equation}
 \mathcal D_8=\sum_{\mu=0}^3e_\mu\partial_{x_\mu}
                    +\sum_{\mu=0}^3(\ell e_\mu)\partial_{y_\mu}.
 \label{eq:hc-D8}
\end{equation}
We call a field \emph{left monogenic} if \(\mathcal D_8\mathcal O=0\).
The pair convention gives the component formula
\begin{equation}
 \mathcal D_8(A,B)=
 \left(D_qA-\sum_{\mu=0}^3B_{y_\mu}\overline e_\mu,
       \overline D_qB+\sum_{\mu=0}^3A_{y_\mu}e_\mu\right).
 \label{eq:hc-D8-components}
\end{equation}
This follows by applying \eqref{eq:hc-octonion-product} to each basis
term of \eqref{eq:hc-D8}.  General octonionic Cauchy--Riemann systems
are discussed in \cite{kauhanenorelma}; their regularity classes must
be distinguished from the distinct Clifford and slice formulations
\cite[Proposition 3.5 and Section 4]{kauhanenorelma19}.

\begin{proposition}[Reflected extension of the physical trace]
\label{prop:hc-reflected-extension}
For every \(m\ge0\), the field
\begin{equation}
 \widetilde{\mathcal O}_m(t,q,r)
 =\bigl(G_m(t,q),G_{m+2}(t,\bar q)\bigr)
 \label{eq:hc-reflected-lift}
\end{equation}
is globally real analytic and left monogenic.  It satisfies
\begin{equation}
 \partial_{x_0}\widetilde{\mathcal O}_m
       =\widetilde{\mathcal O}_{m+1},\qquad
 \partial_t\widetilde{\mathcal O}_m
       =-\widetilde{\mathcal O}_{m+2},\qquad
 \widetilde{\mathcal O}_m(t,x,0)=(G_m(t,x),G_{m+2}(t,x)).
 \label{eq:hc-reflected-identities}
\end{equation}
\end{proposition}
\begin{proof}
Reflection reverses all three imaginary quaternion coordinates, hence
\[
 \bar D_q[B(\bar q)]=(D_qB)(\bar q).
\]
Both \(r\)-derivatives in \eqref{eq:hc-D8-components} vanish, and
both remaining components vanish by \(D_qG_k=0\).  Reflection
commutes with \(\partial_{x_0}\) and \(\partial_t\), proving the
derivative statements.  Real analyticity follows from that of the
\(G_k\), and \(\bar x=x\) gives the physical trace.
\end{proof}

This extension is constant in the four \(r\) coordinates. Its second
component uses \(z_0=x_0-\ii x_1\) and
\(\tau=t-x_2-\ii x_3\), whereas the first uses the original
arguments. A nonzero such field is not in whole-space
\(L^2(\R^8)\); its local square integrability on bounded domains
is unaffected. For
\(m=0\), its values use derivatives of \(H\) only through order
three, evaluated at the original and reflected arguments.  Its
derivatives still generate higher jets.  If the full quaternionic
trace \((G_m(t,q),G_{m+2}(t,q))\) at \(r=0\) is required, the
following extension retains that stronger prescription.

\begin{proposition}[Extension with the original quaternionic trace]
\label{prop:hc-octonion-extension}
For each \(m\ge0\), the field
\begin{equation}
 \mathcal O_m(t,q,r)
 =\left(G_m(t,q)-\tfrac12G_{m+3}(t,q)\overline r,
                        G_{m+2}(t,q)\right)
 \label{eq:hc-octonion-lift}
\end{equation}
is globally real analytic and satisfies
\begin{equation}
 \mathcal D_8\mathcal O_m=0,\qquad
 \partial_{x_0}\mathcal O_m=\mathcal O_{m+1},\qquad
 \partial_t\mathcal O_m=-\mathcal O_{m+2},\qquad
 \mathcal O_m(t,q,0)=(G_m,G_{m+2}).
 \label{eq:hc-octonion-identities}
\end{equation}
\end{proposition}
\begin{proof}
For \eqref{eq:hc-octonion-lift}, \(B\) is independent of \(r\).
Moreover \(D_q(G_{m+3}\overline r)=(D_qG_{m+3})\overline r=0\),
since the component algebra is associative and \(r\) is held fixed.
The first component of \eqref{eq:hc-D8-components} is therefore zero.
The second is
\[
 2G_{m+3}-\tfrac12\sum_{\mu=0}^3
                         G_{m+3}\overline e_\mu e_\mu
 =2G_{m+3}-2G_{m+3}=0.
\]
The derivative identities follow directly from
\eqref{eq:hc-jet-operators}, and the explicit formula proves global
real analyticity.
\end{proof}

The uncorrected pair with the same \(q\) argument, held constant in
\(r\), has defect \((0,2G_{m+3})\).  This does not apply to the
reflected pair.  More generally, two left
Fueter-regular fields \(A(q),B(q)\) admit the extension
\((A-\tfrac12B_{x_0}\overline r,B)\), by the same calculation.
Thus the construction imposes no new condition on the original pair.
For \(m=0\), its value away from \(r=0\) involves \(G_3\), hence
\(h_4\); it is not a closed evolution of only four jets.

Both extensions have the same physical trace \((q,r)=(x,0)\):
\[
 \mathcal O_0=(h_0+\jj h_1)+\ell(h_2+\jj h_3)
 \in\operatorname{span}_{\R}\{1,\jj,\ell,\ell\jj\}.
\]
This span is an associative quaternion subalgebra, as can be checked
from \eqref{eq:hc-octonion-product}. Its total norm is
\begin{equation}
 |\mathcal O_0|^2=E+h_2^2+h_3^2.
 \label{eq:hc-total-norm}
\end{equation}
At a nondegenerate double zero, \(E=0\) but \(h_2\ne0\); the total
octonionic norm is then positive. A collision criterion must retain
the distinguished first pair.

\subsection{Oriented associators preserve the signed observations}

The \emph{associator} measures the failure of associativity:
\[
 [A,B,C]=(AB)C-A(BC),\qquad \mathbf u=\ell\jj.
\]
For embedded quaternions \(a,b\in\HH\), the pair multiplication gives
\begin{equation}
 [a,b,\ell]=\ell(\overline{ab}-\overline{ba})
             =-\ell(ab-ba).
 \label{eq:hc-commutator-associator}
\end{equation}
The last equality holds because a quaternion commutator has zero real
part. The untwisted physical jets lie in the commutative plane
\(\R+\jj\R\), so their associator with \(\ell\) vanishes.
Fixed right multiplication by \(\ii\) and \(\kk\) produces the needed
signed quantities while preserving quaternion left regularity.

\begin{lemma}[Dot products and determinants as associators]
\label{lem:hc-twisted-associators}
For \(a=A+\jj B\), \(b=C+\jj D\) with real \(A,B,C,D\),
\begin{equation}
 [a\ii,b\ii,\ell]=-2\mathbf u(AD-BC),
 \qquad
 [a\ii,b\kk,\ell]=2\mathbf u(AC+BD).
 \label{eq:hc-twisted-associators}
\end{equation}
\end{lemma}
\begin{proof}
One has \(a\ii=A\ii-B\kk\), \(b\ii=C\ii-D\kk\), and
\(b\kk=D\ii+C\kk\). Their quaternion commutators are, respectively,
\(2\jj(AD-BC)\) and \(-2\jj(AC+BD)\). Substitute these into
\eqref{eq:hc-commutator-associator}.
\end{proof}

In particular define, on the physical slice,
\begin{align}
 \mathcal A_E&=[F\ii,F\kk,\ell]=2\mathbf u E,
       \label{eq:hc-associator-energy}\\
 \mathcal A_j&=[G_j\ii,G_{j+1}\ii,\ell]
                    =2\mathbf u\mathcal L_j\quad(j=0,1),
       \label{eq:hc-associator-signs}\\
 \mathcal A_t&=[F\ii,F_t\kk,\ell]=\mathbf u E_t.
       \label{eq:hc-associator-time}
\end{align}
These are immediate from Lemma~\ref{lem:hc-twisted-associators} and
\eqref{eq:hc-energy-evolution}. In contrast to
\eqref{eq:hc-total-norm}, \(\mathcal A_E=0\) is equivalent to \(F=0\).
Equation~\eqref{eq:hc-stationary-energy} becomes
\[
 h_1=0\ \Longrightarrow\ \mathcal A_t=\mathcal A_0,
 \qquad
 h_2=0\ \Longrightarrow\ \mathcal A_t=\mathcal A_1.
\]

There is no order on \(\OO\) being invoked. Specify the oriented ray
\(\mathcal C_{\mathbf u}=\{s\mathbf u:s\ge0\}\).
The two stationary sign conditions are exactly
\begin{equation}
 \mathcal A_t\in\mathcal C_{\mathbf u}
          \quad\hbox{on}\quad\{h_1h_2=0\}.
 \label{eq:hc-oriented-sign}
\end{equation}
Likewise the locally bounded loss condition is exactly
\begin{equation}
 \mathcal A_t+K\mathcal A_E\in\mathcal C_{\mathbf u}
          \quad\hbox{on}\quad\{h_1h_2=0\},
 \qquad K\ge0,
 \label{eq:hc-oriented-loss}
\end{equation}
with \(K\) locally bounded through every candidate collision. Indeed
the coefficient on this ray is \(E_t+2KE\). Thus the associators
preserve both the collision detector and the required orientation.
Their norms alone would replace the signed quantities by absolute values.

\subsection{What calibration, symmetry, and logarithmic norms retain}

For imaginary octonions, the Harvey--Lawson identity is
\begin{equation}
 \varphi(a,b,c)^2+\tfrac14|[a,b,c]|^2=|a\wedge b\wedge c|^2,
 \qquad \varphi(a,b,c)=\langle a,bc\rangle.
 \label{eq:hc-calibration}
\end{equation}
Here the wedge denotes exterior volume, and
\(\langle a,b\rangle=\Rea(\bar a b)\).  See
\cite[Chapter IV, Theorem 1.6]{harveylawson} and
\cite[Corollary 3.52]{karigiannis}.  For the twisted triples
in Lemma~\ref{lem:hc-twisted-associators}, \(\varphi=0\), since the
first two arguments lie in \(\HH\) and the third is \(\ell\).
Thus \eqref{eq:hc-calibration} reduces to the planar identity
\eqref{eq:hc-planar-gram}.  In particular \(F\ii\) and \(F\kk\)
are orthogonal with norm \(|F|\), and
\[
 |[F\ii,F\kk,\ell]|=2E.
\]
Whenever \(F\ne0\), this triple spans the fixed three-plane
\(\operatorname{span}_{\R}\{\ii,\kk,\ell\}\).  Its volume is a
useful consistency check, but its norm does not retain a signed
coefficient.  No calibrated-submanifold condition on a varying tangent
plane has been established by this construction.

An octonion automorphism \(g\in G_2\) satisfies
\([gA,gB,gC]=g[A,B,C]\).  Transforming both the observation and the
oriented ray therefore leaves \eqref{eq:hc-oriented-loss} unchanged.
The larger triality symmetry acts through three related
representations, rather than a single algebra automorphism
\cite[Section 2.4]{baez}\cite{schraymanogue}.
Any use of it here must also transport the distinguished first-pair
projection and the stationary constraints.  Mixing that pair with the
second pair without transporting the projection changes the collision
detector, as \eqref{eq:hc-total-norm} already shows.

The special holomorphic structure does give an exact logarithmic-norm
identity.  Put
\[
 \widehat E(t,q)=|H(\tau,z_0)|^2+|H_z(\tau,z_0)|^2,
 \qquad \Delta_q=\sum_{\mu=0}^3\partial_{x_\mu}^2.
\]
At a physical point with \(E>0\),
\begin{equation}
 \left.\Delta_q\log\widehat E\right|_{q=x}
 =\frac{4\{\Lag_0^2+(h_0h_3-h_1h_2)^2\}}{E^2}.
 \label{eq:hc-log-energy}
\end{equation}
To verify it, write \(v=(H,H_z)\).  In either complex coordinate
\(\zeta=z_0,\tau\), holomorphy gives
\[
 \partial_\zeta\partial_{\bar\zeta}\log|v|^2
 =\frac{|v|^2|v_\zeta|^2-|\langle v,v_\zeta\rangle_{\C}|^2}
        {|v|^4}
 =\frac{|\det(v,v_\zeta)|^2}{|v|^4}.
\]
The complex inner product is Hermitian. The same formula in any complex
direction proves that \(\log\widehat E\) is plurisubharmonic away
from the common zero set. On the physical slice the two
determinants are \(-\Lag_0\) and
\(-(h_0h_3-h_1h_2)\), while
\(\Delta_q=4(\partial_{z_0}\partial_{\bar z_0}
                   +\partial_\tau\partial_{\bar\tau})\).
This proves \eqref{eq:hc-log-energy}.  The nonnegative right side has
squared away the required orientation and is defined only where
\(E>0\); it supplies no sign for the external derivative \(E_t\).

This logarithmic positivity comes from the holomorphic pair, not
generic octonionic monogenicity.  The monogenic Cauchy kernel
\(\bar p/|p|^8\), away from the origin in eight dimensions, has norm
\(|p|^{-7}\), and
\[
 \Delta_8\log|p|^{-7}=-42/|p|^2<0.
\]
Likewise a frame chosen using \(F/|F|\) is singular at a potential
collision.  A source-prescribed frame can be used, with its derivative
terms retained, but cannot assume the desired first-pair nonvanishing.

\subsection{Why a source-specific estimate remains necessary}

The constructions above apply to every jointly entire backward-heat
family real on the physical slice. The following control shows that
even strict positivity and analytic decay of its cosine-transform source
do not supply the missing sign.

\begin{proposition}[A positive entire source with a positive-time collision]
\label{prop:hc-positive-source-control}
There exist strictly positive, even, entire functions \(\Phi_*\),
decaying superexponentially on the real axis, whose heat cosine
transforms are jointly entire and have a nondegenerate real double zero
at a positive time. Their quaternionic and octonionic lifts satisfy all
the preceding identities.
\end{proposition}
\begin{proof}
Set \(\rho(v)=2e^{-v^4}/\Gamma(1/4)\), so that
\(\int_{\R}\rho(v)\,dv=1\), and
\(\rho_\epsilon(v)=\epsilon^{-1}\rho(v/\epsilon)\) for real
\(\epsilon>0\). Use
\[
 \Phi_\epsilon(u)=
 2\{\rho_\epsilon(u-1)+\rho_\epsilon(u+1)\}
   +\rho_\epsilon(u-2)+\rho_\epsilon(u+2).
\]
This source is strictly positive, even, and entire, with quartic decay
on the real axis. Its half-line heat cosine transform is
\begin{align*}
 \widetilde H_\epsilon(t,z)
 &=2\int_{\R}\rho(v)e^{t(1+\epsilon v)^2}
                         \cos\bigl(z(1+\epsilon v)\bigr)\,dv\\
 &\quad+\int_{\R}\rho(v)e^{t(2+\epsilon v)^2}
                         \cos\bigl(z(2+\epsilon v)\bigr)\,dv.
\end{align*}
Quartic decay supplies an integrable majorant on every compact set of
complex \((t,z,\epsilon)\). Hence this formula extends jointly
entirely through \(\epsilon=0\), obeys backward heat, and there becomes
\[
 \widetilde H_0(t,z)=2e^t\cos z+e^{4t}\cos(2z).
\]
At \(T_0=(\log2)/3\), \(x_0=\pi\), it satisfies
\(\widetilde H_0=(\widetilde H_0)_x=0\) and
\((\widetilde H_0)_{xx}=-6e^{T_0}\ne0\).
By the heat equation the Jacobian of
\((\widetilde H_\epsilon,(\widetilde H_\epsilon)_x)\)
with respect to \((t,x)\), at this point and \(\epsilon=0\), has
determinant \(-((\widetilde H_0)_{xx})^2\ne0\).
The real analytic implicit-function
theorem therefore gives a nearby nondegenerate real double zero at a
positive time for every sufficiently small \(\epsilon>0\).
\end{proof}

The failure is visible in the signed observation. At such a collision
\((T,a)\), put \(c=H_{xx}(T,a)\ne0\). The implicit-function theorem
gives a real stationary branch \(H_x(t,\chi(t))=0\) through it.
For \(\delta=T-t\downarrow0\), backward heat gives
\[
 H(t,\chi(t))=c\delta+O(\delta^2),\qquad
 E(t,\chi(t))=c^2\delta^2+O(\delta^3),\qquad
 \mathcal L_0(t,\chi(t))=-c^2\delta+O(\delta^2).
\]
Thus \eqref{eq:hc-oriented-loss} would require
\(K\ge\delta^{-1}+O(1)\), contradicting local boundedness at the
collision. The lift and associator identities alone cannot provide that
bound. For the theta family, its particular arithmetic structure must
enter a further signed estimate.

For the genuine theta family, Theorem~\ref{thm:outer} supplies the
additional sign information in its declared outer region:
$\mathcal A_0$ and $\mathcal A_1$ lie strictly on
$\mathcal C_{\mathbf u}$ there. At genuine stationary points the same
is true of $\mathcal A_t$. This is the geometric consequence of the
common arithmetic amplitude, not a consequence of monogenicity alone.



\section{Centered residuals and the common theta carrier}\label{sec:residuals}

The outer theorem uses oscillatory cancellation to reach a nonzero zeta
comparison amplitude. To study the remaining region, we return to the
exact theta source.
The aim is to retain the same four derivatives and their signed
correlations before estimating any error.

\subsection{A contour identity for the physical derivatives}
\begin{proposition}[Complete-source contour shift]\label{prop:complete-shift}
For real $t\ge0$, $x\ge0$, $0<a<\pi/8$, and every integer $k\ge0$,
\begin{equation}
 h_k(t,x)=e^{-ax}\Re\int_0^\infty
 (i(u+ia))^k e^{t(u+ia)^2}\Phi(u+ia)e^{ixu}\dd u.
 \label{eq:complete-shift}
\end{equation}
\end{proposition}
\begin{proof}
The theta series is normally convergent in the strip
$|\Im z|<\pi/8$ because $\Re(e^{4z})>0$ there. Analytic continuation
of the real reflection identity makes $\Phi$ even throughout the
strip; Schwarz symmetry then makes $\Phi(iy)$ real. On the left edge
of the contour shift, every factor of
$(iz)^ke^{tz^2}\Phi(z)e^{ixz}$ is real at $z=iy$, while $\dd z=i\dd y$.
Its contribution has zero real part. On the right edge,
$\cos(4y)\ge\cos(4a)>0$ gives uniform double-exponential decay.
Cauchy's theorem proves the formula. The evenness argument applies
to the complete source; it must precede any separation into individual
theta summands.
\end{proof}

The all-index and integration bounds in Appendix~\ref{app:payments}
make truncation of \eqref{eq:complete-shift} rigorous. They control
absolute errors at the gamma exponential scale. Obtaining a sign still
requires information about the retained, correlated jets.

\subsection{A residual adapted to genuine stationarity}
Freeze real $a,b$ with $0<a<\pi/8$ and $b\ge0$, and put $\mu=a^2+b^2$.
For every integer $j\ge0$, define
\begin{equation}
 C_j=e^{-ax}\Re\int_0^\infty (u^2-b^2)(i(u+ia))^j
             e^{t(u+ia)^2}\Phi(u+ia)e^{ixu}\dd u.
 \label{eq:C-integral}
\end{equation}
The required stationary forms use $j=0,1$; higher residuals will enter
exact heat-transport identities. The polynomial identity
$u^2-b^2=-(i(u+ia))^2-2a\,i(u+ia)-\mu$ gives
\begin{equation}
 C_j=-h_{j+2}-2ah_{j+1}-\mu h_j,
 \qquad h_{j+1}=0\ \Longrightarrow\
 \Lag_j=\mu h_j^2+h_jC_j.
 \label{eq:C-stationary}
\end{equation}
Thus the contour drift disappears at exactly the original stationary
constraint. No derivative of an exponentially normalized observation
has replaced that constraint.

\begin{proposition}[A sufficient residual cone]\label{prop:residual-cone}
If $\eta>0$ and
\begin{equation}
 \Delta:=\eta^2E-C_0^2-C_1^2\ge0,
 \label{eq:cone-slack}
\end{equation}
then the respective stationary sets satisfy
\begin{equation}
 \Lag_0\ge(\mu-\eta)E\quad(h_1=0),\qquad
 \Lag_1\ge-\frac{\eta^2}{4\mu}E\quad(h_2=0).
 \label{eq:cone-losses}
\end{equation}
A sufficient common negative-loss coefficient is $K=\eta^2/(4\mu)$.
It is locally bounded wherever the prescribed parameters are locally
bounded and $\mu$ stays positive.
\end{proposition}
\begin{proof}
At $h_1=0$, $E=h_0^2$ and $|C_0|\le\eta\sqrt E$, so the first
inequality follows from \eqref{eq:C-stationary}. At $h_2=0$, complete
the square in $\mu h_1^2-\eta|h_1|\sqrt E$ to obtain the second.
Finally $\eta-\mu\le\eta^2/(4\mu)$ because
$(\eta-2\mu)^2\ge0$. No division by $E$ occurs.
\end{proof}

For use with exact algebra or interval certificates, the same argument
has the following identities before stationarity is imposed:
\begin{align}
 2\eta[\Lag_0+(\eta-\mu)E]
 &= (C_0+\eta h_0)^2+C_1^2+\Delta\notag\\
 &\quad+2\eta h_1[(1-\mu+\eta/2)h_1+2ah_0],
 \label{eq:cone-SOS0}\\
 4\mu\Lag_1+\eta^2E
 &= (2\mu h_1+C_1)^2+C_0^2+\Delta
             +4\mu h_2(h_2+2ah_1).
 \label{eq:cone-SOS1}
\end{align}
Each final term is a multiple of the genuine stationary derivative.
The cone is sufficient, but a one-sided bound for $h_jC_j$ can be
substantially weaker and therefore more useful.

For example, suppose paid normalized approximations give
$Z_j=e^{\pi x/8}h_j=p+\delta$ and
$e^{\pi x/8}C_j=q+\epsilon'$, where
$|\delta|\le\varepsilon_0$, $|\epsilon'|\le\zeta_0$.
At $h_{j+1}=0$, expansion of \eqref{eq:C-stationary}, retaining
$\mu\delta^2\ge0$, gives
\begin{equation}
 e^{\pi x/4}\Lag_j\ge\mu p^2+pq
 -|2\mu p+q|\varepsilon_0-|p|\zeta_0-\varepsilon_0\zeta_0.
 \label{eq:paid-signed-residual}
\end{equation}
Every truncation, quadrature, and parameter error must be included in
the two bounds. The retained product $pq$ preserves its sign.

\subsection{One carrier for all four derivatives}
There is an exact arithmetic reduction behind the centered residual.
Put $d_n=\tfrac12\log n$. The source summands obey
\begin{equation}
 \phi_n(z)=n^{-1/2}\phi_1(z+d_n).
 \label{eq:theta-translation}
\end{equation}
For $v\ge d_n$, define
\begin{equation}
 A_n(t,x;v)=n^{-1/2}
       \exp\{td_n^2-2td_n(v+ia)-ixd_n\}.
 \label{eq:An}
\end{equation}
Changing variables $v=u+d_n$ in the complete shifted source yields
\begin{equation}
 h_j=e^{-ax}\Re\int_0^\infty
 e^{t(v+ia)^2}\phi_1(v+ia)e^{ixv}
 \sum_{n\le e^{2v}}A_n(t,x;v)[i(v+ia-d_n)]^j\dd v.
 \label{eq:shared-carrier}
\end{equation}
In particular $j=0,1,2,3$ uses only the first four logarithmic moments
of the same arithmetic weights. The staircase cutoff $n\le e^{2v}$
is part of the identity. It is not an independently chosen truncation.

To justify the rearrangement, let $0\le t\le\tau$. Since
$0\le d_n\le v$, we have $|A_n|\le n^{-1/2}$ and
$\sum_{n\le e^{2v}}n^{-1/2}\le2e^v$. Therefore the absolute integrand
in \eqref{eq:shared-carrier}, before $e^{-ax}$, is bounded by
\begin{equation}
 2(2\pi^2+3\pi)(v+a)^j
           e^{\tau v^2+10v-\pi\cos(4a)e^{4v}}.
 \label{eq:shared-majorant}
\end{equation}
This is integrable. It also gives a full integration tail, stated in
Appendix~\ref{app:payments}; no unretained index is omitted from that
bound.

The common oscillatory exponential in the original $n$-th channel has
phase $(x+2ta)u-\pi n^2\sin(4a)e^{4u}$. Its critical location is
\begin{equation}
 b_n=\frac14\log\frac{x+2ta}{4\pi n^2\sin(4a)}
     =b_1-d_n.
 \label{eq:channel-centers}
\end{equation}
Only nonnegative locations belong to the integration half-line.
This observation supplies a centering choice; it does not assert
first-channel dominance or a full stationary-phase approximation.
Taking $b=b_1$, the multiplier defining $C_j$ becomes
\begin{equation}
 (v-d_n)^2-b^2
 =\underbrace{-b\log n+\tfrac14\log^2n}_{\delta_n}
       +2(b-d_n)(v-b)+(v-b)^2.
 \label{eq:logarithmic-residual}
\end{equation}
Thus the remaining signed residual involves logarithmic weights of the
same arithmetic sum, rather than unrelated errors for each derivative.
At $v=b$ those weights have the transparent form
\begin{equation}
 A_n=n^{-1/2-tb-i(x/2+ta)}e^{t\log^2n/4}.
 \label{eq:saddle-weights}
\end{equation}
The factors $A_n$ are complex and their products retain arithmetic
phases. A positive-weight covariance inequality does not apply to them.

For perspective, if $b>0$ the frozen-parameter balanced quaternion
$W=h_0+\jj(h_1+ah_0)/b$ obeys
\begin{equation}
 W_t=(b^2-a^2-2ab\jj)W+R,\qquad
 R=C_0+\jj(C_1-aC_0)/b.
 \label{eq:balanced-residual-evolution}
\end{equation}
The skew term rotates $W$ without changing its norm. Equation
\eqref{eq:balanced-residual-evolution} separates the expected oscillatory
motion from the residual, whereas \eqref{eq:C-stationary} keeps the
precise scalar conditions needed for a sign. If $a,b$ vary with time,
the additional term is
$\jj(a_t\Re W/b-b_t\Im_{\jj}W/b)$; all cell computations below freeze
them before differentiating.


\subsection{Staircase derivatives and exact entry terms}
Put $T=x/2+ta$ and, for integers $r\ge0$,
\[
 B_r(v)=\sum_{n\le e^{2v}}(\log n)^r A_n(t,x;v),
 \qquad
 D_j(v)=\sum_{n\le e^{2v}}A_n(t,x;v)[i(v+ia-d_n)]^j.
\]
For $r=0$ the logarithmic weight is one, including at $n=1$.
On $v>0$, distributional differentiation gives
\begin{equation}
 B_r'=-tB_{r+1}+
 \sum_{n\ge2}(\log n)^r A_n(t,x;d_n)\delta(v-d_n).
 \label{eq:staircase-recurrence}
\end{equation}
Between entry points,
\begin{equation}
 D_j'=-2t(v+ia)D_j-2itD_{j+1}+ijD_{j-1},
 \label{eq:staircase-jet-recurrence}
\end{equation}
where the last term is omitted for $j=0$. The jump at $d_n$ is
$A_n(t,x;d_n)(-a)^j$. If
$K(v)=e^{t(v+ia)^2}\phi_1(v+ia)e^{ixv}$, its full contribution is
\begin{equation}
 K(d_n)A_n(t,x;d_n)(-a)^j
   =e^{-ta^2}\phi_n(ia)(-a)^j.
 \label{eq:staircase-cancellation}
\end{equation}
Indeed $\partial_v A_n=-2td_nA_n$ and
$d_n=v+ia+i[i(v+ia-d_n)]$ give the smooth recurrence; direct
substitution and \eqref{eq:theta-translation} give the jump formula.
At $v=0$ the index $n=1$ is already present and belongs to the lower
endpoint of any integration by parts. Every entering-index phase in
\eqref{eq:staircase-cancellation} cancels in $x$. These terms cannot
be assigned an artificial oscillatory gain.

The oscillatory block estimate in Section~\ref{sec:outer-proof} can also
be applied to this exact carrier. On a block $M<n\le R(v)$, where
$R(v)=\min(2M,e^{2v})$, the logarithmic-moment weight is
\[
 w_r(y)=(\log y)^r e^{t\log^2y/4}y^{-1/2-tv}.
\]
For $M\ge e^{2r}$ it is decreasing, since
$y(\log w_r)'\le r/\log y-1/2\le0$; for $r=0$ this holds for
all $M\ge1$. Abel summation therefore gives, when $|T|\ge1$,
\[
 \left|\sum_{M<n\le R(v)}w_r(n)n^{-iT}\right|
 \ll w_r(M)\bigl(|T|^{1/6}M^{1/2}+|T|^{-1/6}M\bigr).
\]
For other blocks use the full total variation of $w_r$, and for
$|T|<1$ use the trivial bound. The actual shortened endpoint remains
$R(v)$; these complex block estimates do not imply a signed product
bound.

\subsection{The complete paired observations}
Define the staircase measure space and its coordinates by
\begin{equation}
 \begin{split}
 \Omega&=\{(v,n):n\ge1,\ v\ge d_n\},\qquad
 d\nu=dv\times\text{counting measure},\\
 y_\alpha&=v-d_n,\qquad q_\alpha=-a+iy_\alpha,\\
 W_\alpha&=e^{t(v+ia)^2}\phi_1(v+ia)e^{ixv}A_n(t,x;v),
 \qquad \alpha=(v,n).
 \end{split}
 \label{eq:pair-domain}
\end{equation}
Then $h_k=e^{-ax}\Re\int_\Omega W_\alpha q_\alpha^k\,d\nu(\alpha)$.
The two factors in the following integrals have independent coordinates
and independent indices.
\begin{proposition}[Paired Laguerre and residual forms]
For $j=0,1$,
\begin{align}
 \Lag_j=\frac{e^{-2ax}}4\Re\iint_{\Omega^2}\bigl[&
 (y_\alpha+y_\beta)^2 W_\alpha\overline{W_\beta}
       (q_\alpha\overline{q_\beta})^j\notag\\
 &+(y_\alpha-y_\beta)^2 W_\alpha W_\beta
       (q_\alpha q_\beta)^j\bigr]\,d\nu(\alpha)d\nu(\beta).
 \label{eq:paired-laguerre}
\end{align}
Writing
\begin{equation}
 J_j(\alpha,\beta)=W_\alpha\overline{W_\beta}
       (q_\alpha\overline{q_\beta})^j
       +W_\alpha W_\beta(q_\alpha q_\beta)^j,
 \label{eq:paired-J}
\end{equation}
the centered product is
\begin{equation}
 h_jC_j=\frac{e^{-2ax}}4\Re\iint_{\Omega^2}
 (y_\alpha^2+y_\beta^2-2b^2)J_j(\alpha,\beta)\,
 d\nu(\alpha)d\nu(\beta).
 \label{eq:paired-residual}
\end{equation}
Before imposing stationarity, the exact difference is
\begin{equation}
 \Lag_j-\mu h_j^2-h_jC_j
       =h_{j+1}(h_{j+1}+2ah_j).
 \label{eq:stationary-ideal}
\end{equation}
\end{proposition}
\begin{proof}
Use $\Re A\Re B=\tfrac12\Re(A\overline B+AB)$ and symmetrize
in $\alpha,\beta$. The Hermitian difference in the Laguerre form is
\[
 q_\alpha\overline{q_\beta}
 -\tfrac12(q_\alpha^2+\overline{q_\beta}^{2})
 =-\tfrac12(q_\alpha-\overline{q_\beta})^2
 =\tfrac12(y_\alpha+y_\beta)^2;
\]
the bilinear difference is
$-\tfrac12(q_\alpha-q_\beta)^2
=\tfrac12(y_\alpha-y_\beta)^2$.
Conjugate-swapping the Hermitian term preserves its real integral.
The same polarization applied to the real multiplier $y^2-b^2$
gives \eqref{eq:paired-residual}, and \eqref{eq:C-stationary} gives
\eqref{eq:stationary-ideal}. The majorant
\eqref{eq:shared-majorant} has every polynomial moment integrable.
Together with $(y\pm z)^2\le2(y^2+z^2)$, this proves absolute
convergence of each displayed integral and justifies Fubini.
\end{proof}

The squared frequencies in \eqref{eq:paired-laguerre} do not give the
phase of either amplitude product a sign. As a finite algebra control,
two real signed atoms at frequencies $1,2$, with weights $1,-1/2$
and $a=0$, give $(h_0,h_1,h_2)=(1/2,0,1)$ and $\Lag_0=-1/2$.
The diagonal contributions sum to $2$ and the ordered cross terms to
$-5/2$. This control lies outside the theta class and positive-contour
hypotheses; it explains why retaining only diagonal pairs loses the
observation.

\subsection{A complete omitted-pair payment}
For a measurable retained set $A\subset\Omega$, denote its jets by
$h_k^A$ and set $\Lag_j^A=(h_{j+1}^A)^2-h_j^Ah_{j+2}^A$. Define
absolute masses before the exterior exponential normalization:
\[
 M_k^A=\int_A|W_\alpha||q_\alpha|^k\,d\nu,
 \qquad T_k^A=\int_{\Omega\setminus A}|W_\alpha||q_\alpha|^k\,d\nu.
\]
Then
\begin{equation}
 |\Lag_j-\Lag_j^A|\le e^{-2ax}
 \bigl(M_j^AT_{j+2}^A+M_{j+2}^AT_j^A+T_j^AT_{j+2}^A\bigr).
 \label{eq:pair-tail-payment}
\end{equation}
Indeed, including the factor $1/4$ in \eqref{eq:paired-laguerre},
the absolute sum of its two kernels is at most
\[
 \tfrac12\bigl(|W_\alpha||q_\alpha|^{j+2}
              |W_\beta||q_\beta|^j
       +|W_\alpha||q_\alpha|^j
              |W_\beta||q_\beta|^{j+2}\bigr),
\]
using $(y+z)^2+(y-z)^2=2(y^2+z^2)$ and $y^2\le|q|^2$.
Integrating over $A\times A^c$, $A^c\times A$, and $A^c\times A^c$
gives the claim, including all omitted mixed and omitted--omitted pairs.
No bound for $T_{j+1}^A$ is required.

For $A=\{(v,n)\in\Omega:v\le V,\ n\le N\}$, use the complete
index payment and the shared-coordinate integration payment in
Appendix~\ref{app:payments}, removing their $e^{-ax}$ normalization
before using them for $T_k^A$. Overlap may be overpaid. Retained
quadrature and parameter variation require their own errors. The
absolute payment \eqref{eq:pair-tail-payment} is not a relative bound
through an unknown zero of $E$.

\subsection{Signed staircase cutoffs on the closed outer domain}
Fix $a=\pi/16$ and retain exactly
\begin{equation}
 \Omega_N=\{(v,n)\in\Omega:0\le v\le V_N\},\qquad
 V_N=\tfrac12\log N,
 \label{eq:signed-staircase-domain}
\end{equation}
where the intrinsic restriction $n\le e^{2v}$ remains part of
$\Omega$. There is no additional rectangular approximation to that
staircase. Write $h_k^{(N)}$ for the restricted common-source moments
and $\Lag_j^{(N)}=(h_{j+1}^{(N)})^2-h_j^{(N)}h_{j+2}^{(N)}$.
All cross terms in these products are retained.

To avoid confusing source constants with the centered residual $C_0$,
use the notation
\begin{equation}
 \begin{gathered}
 \beta_0=\pi\cos(4a)=\frac\pi{\sqrt2},\qquad
 \mathsf C_0=2(2\pi^2+3\pi),\qquad
 c=\beta_0/2,\qquad \mathsf D_0=\frac{\mathsf C_0}{2\beta_0},\\
 \mathsf M_k=\mathsf C_0\int_0^\infty(v+a)^k
   \exp\{v^2/2+10v-\beta_0e^{4v}\}\dd v,\qquad 0\le k\le3,\\
 B_j=(\mathsf M_j+\mathsf M_{j+2})\mathsf D_0+\mathsf D_0^2,
 \qquad j=0,1.
 \end{gathered}
 \label{eq:signed-cutoff-constants}
\end{equation}
These finite constants depend only on the fixed complete source and
$a$, not on $t,x,N$.

\begin{corollary}[Signed staircase cutoff estimates]
\label{cor:signed-staircase-cutoffs}
For all integers $N\ge8$, $j=0,1$, and $(t,x)$ in
\eqref{eq:outer-domain},
\begin{equation}
 \Lag_j^{(N)}\ge-A_j^*(t,x)E,
 \qquad
 \Lag_j^{(N)}\ge-e^{-2ax}B_j e^{-cN^2},
 \label{eq:signed-cutoff-conclusions}
\end{equation}
where
\begin{equation}
 A_j^*(t,x)=
 \frac{e^{-2ax}\mathsf M_j\mathsf M_{j+2}\zeta(9/8)^2}
      {\Gamma(t,x)^2},\qquad
 \Gamma(t,x)=|M_t((1+ix)/2)|.
 \label{eq:signed-cutoff-relative-coefficient}
\end{equation}
The coefficient is prescribed independently of $N$ and of the measured
physical energy; it is locally bounded but is not asserted uniformly
bounded as $x\to\infty$. The second estimate has zero energy loss and
height-uniform Gaussian tail constants. The observations and stationary
sets remain those of the full function.
\end{corollary}
\begin{proof}
For $0\le t\le1/2$, \eqref{eq:shared-majorant} bounds every retained
absolute $k$-th moment by $\mathsf M_k$. Summing the absolute values of
the Hermitian and bilinear kernels in \eqref{eq:paired-laguerre}, using
$(y+z)^2+(y-z)^2=2(y^2+z^2)$ and its factor $1/4$, gives
\begin{equation}
 |\Lag_j^{(N)}|\le
 e^{-2ax}\mathsf M_j\mathsf M_{j+2}.
 \label{eq:signed-cutoff-absolute}
\end{equation}
The energy envelope \eqref{eq:closed-sign-envelopes} and
\eqref{eq:signed-cutoff-relative-coefficient} therefore prove the first
inequality in \eqref{eq:signed-cutoff-conclusions}, with a spare factor
two. Its denominator is the explicitly nonzero analytic carrier, not
$E$; this proves the asserted independence and local boundedness.

For the second inequality, the entire omitted source region is
$v>V_N$. For $v\ge1$,
\[
 \log(v+a)\le v,\qquad
 v^2/2+13v\le(\beta_0/2)e^{4v}.
\]
The latter follows because $e^{-4v}(v^2/2+13v)$ is decreasing there
and its value at one is less than $\beta_0/2$. Since $V_N\ge1$ for
$N\ge8$, every omitted moment of order $k\le3$ satisfies
\begin{equation}
 T_k^{(N)}\le
 \frac{\mathsf C_0}{2\beta_0N^2}e^{-cN^2}
 \le\mathsf D_0e^{-cN^2}.
 \label{eq:signed-cutoff-tail-moments}
\end{equation}
Indeed substitute $q=e^{4v}$ in the integral of $e^{-ce^{4v}}$ and
use $q^{-1}\le N^{-2}$. Apply the complete payment
\eqref{eq:pair-tail-payment}, including both mixed sectors and the
omitted--omitted sector. With the constants in
\eqref{eq:signed-cutoff-constants} it gives
\begin{equation}
 |\Lag_j-\Lag_j^{(N)}|
 \le e^{-2ax}B_j e^{-cN^2}.
 \label{eq:signed-cutoff-complete-tail}
\end{equation}
The independently proved full sign $\Lag_j>0$ establishes the second
inequality in \eqref{eq:signed-cutoff-conclusions}.
\end{proof}

The absolute and tail estimates in the proof hold for every $x\ge0$
and $0\le t\le1/2$; the signed conclusions use the independent outer
positivity theorem and are asserted only in its domain. They do not
assert positivity of every truncated form. Also $h_{j+1}=0$ does not
imply $h_{j+1}^{(N)}=0$; no such substitution is made. At $E=0$ on the corresponding full stationary set, an unpaid inequality $\Lag_j^{(N)}\ge-AE$ would require
$\Lag_j^{(N)}\ge0$ for every cutoff, which does not follow merely from
$\Lag_j=0$. The Gaussian additive allowance avoids imposing that extra
cutoff condition.

One exact alternative retaining the complete amplitude is
\[
 \mathcal J_j^{(N)}=-h_jh_{j+2}^{(N)}
 \quad\hbox{on }h_{j+1}=0,\qquad j=0,1.
\]
It obeys
\begin{equation}
 \mathcal J_j^{(N)}-\Lag_j=h_j(h_{j+2}-h_{j+2}^{(N)}),\qquad
 |\mathcal J_j^{(N)}-\Lag_j|
       \le\sqrt E\,e^{-ax}T_{j+2}^{(N)}.
 \label{eq:full-amplitude-cutoff-error}
\end{equation}
Thus $\mathcal J_j^{(N)}$ vanishes at $E=0$ without imposing positivity
of a raw truncated quadratic form. This is an exact error identity,
not a new signed estimate in the remaining cusp.

% Proposed insertion: after eq:full-amplitude-cutoff-error and its
% concluding paragraph, before sec:fc-certificate.
% Drafted for Edward Baker, 15 October 2026.
% Substantial LLM assistance; parent session supplies recorded model/effort.
% New bibliography keys: csordas, dimitrovxu, elkiesgross,
% kim93, pollackg2, pollackspin8, shan26, viazovska, romik, lee24,
% lee26polynomial, lee26positive.

\section{Complete-source correlation kernels}
\label{sec:complete-correlations}

The paired observations can also be organized on the complete real
source.  This gives a scalar correlation family containing both
Laguerre forms and the physical energy.  The construction is related
to the associated kernels of \cite[Theorem~3.7]{csordas} and the
Wronskian identities of \cite[Theorem~1.3]{dimitrovxu}.
We derive the formulas with the present normalization rather than
importing additional log-concavity hypotheses for the deformed source.
They preserve the signed problem and do not establish a new sign region.

\subsection{The correlation identities and their heat hierarchy}

Use the even extension of the complete source and put
\begin{equation}
\begin{gathered}
 w_t(u)=e^{tu^2}\Phi(u),\qquad
 H(t,x)=\frac12\int_{\R}w_t(u)e^{ixu}\dd u,\\
 \mathscr K_m(t,s)=
 \int_{\R}r^{2m}w_t(s+r)w_t(s-r)\dd r,
 \qquad m\ge0.
\end{gathered}
 \label{eq:corr-definitions}
\end{equation}
Here \(t,x,s\) are real and \(m\) is an integer.  The following
complete-source identities hold at every real time; their domination
is uniform when time is restricted to a bounded interval.

\begin{proposition}[Complete scalar correlations]
\label{prop:complete-correlations}
Each \(\mathscr K_m(t,\cdot)\) is real, even, positive, and integrable
with every polynomial weight.  The two Laguerre observations are
\begin{align}
 \Lag_0(t,x)
 &=\int_{\R}\mathscr K_1(t,s)\cos(2xs)\dd s,
 \label{eq:corr-L0}\\
 \Lag_1(t,x)
 &=\int_{\R}
   \{\mathscr K_2(t,s)-s^2\mathscr K_1(t,s)\}
                     \cos(2xs)\dd s.
 \label{eq:corr-L1}
\end{align}
The exact time identities are
\begin{align}
 \mathscr K_m(t,s)
 &=e^{2ts^2}\int_{\R}r^{2m}e^{2tr^2}
                    \Phi(s+r)\Phi(s-r)\dd r,
 \label{eq:corr-heat-factor}\\
 \partial_t^\ell\mathscr K_m(t,s)
 &=2^\ell\sum_{k=0}^{\ell}\binom{\ell}{k}
                   s^{2(\ell-k)}\mathscr K_{m+k}(t,s)
 \qquad(\ell\ge0).
 \label{eq:corr-heat-hierarchy}
\end{align}
In particular,
\begin{equation}
 (\partial_t+\partial_x^2)\Lag_0=2\Lag_1.
 \label{eq:corr-L-heat}
\end{equation}
\end{proposition}
\begin{proof}
By evenness, \eqref{eq:source-majorant} gives a
double-exponential majorant in \(|u|\).  On \(|t|\le T\), multiplying
that majorant by \(e^{Tu^2}\) still gives double-exponential decay.
Thus the source's
absolute moments of every order are finite.  Under
\[
 u=s+r,\qquad v=s-r,\qquad du\,dv=2\,ds\,dr,
\]
every polynomial in \(s,r\) is a polynomial in \(u,v\).
The double integrals below therefore converge absolutely, and the
same argument pays every finite time or Fourier derivative.
Evenness and positivity of \(\mathscr K_m\) follow directly from
evenness and strict positivity of \(w_t\); for \(m>0\) the integrand
vanishes only at \(r=0\).

The full-line formula for \(h_k\) has multiplier \((iu)^k/2\).
Symmetrizing the product defining \(\Lag_j\), for \(j=0,1\), gives
\[
 \Lag_j(t,x)=\frac18\iint_{\R^2}
     (-uv)^j(u-v)^2w_t(u)w_t(v)e^{ix(u+v)}\dd u\dd v.
\]
The change of variables has \((u-v)^2=4r^2\),
\(-uv=r^2-s^2\), and phase \(e^{2ixs}\).
Its Jacobian cancels the remaining factor \(1/2\), proving
\eqref{eq:corr-L0}--\eqref{eq:corr-L1} after taking real parts.
The imaginary parts vanish by evenness in \(s\).

Since \((s+r)^2+(s-r)^2=2(s^2+r^2)\), differentiation of
\eqref{eq:corr-heat-factor} and the binomial theorem give
\eqref{eq:corr-heat-hierarchy}.  At \(\ell=1,m=1\),
\[
 \partial_t\mathscr K_1=2s^2\mathscr K_1+2\mathscr K_2.
\]
Inserting this into \eqref{eq:corr-L0}, and using
\(\partial_x^2\cos(2xs)=-4s^2\cos(2xs)\), proves
\eqref{eq:corr-L-heat}.
\end{proof}

The last identity has the independent jet verification
\[
 \partial_t\Lag_0=h_2^2-2h_1h_3+h_0h_4,\qquad
 \partial_x^2\Lag_0=h_2^2-h_0h_4.
\]
It does not make the second stationary observation redundant: a
nonnegative function need not have a nonnegative value of
\(\partial_t+\partial_x^2\).

\begin{proposition}[The energy in the same correlation family]
\label{prop:correlation-energy}
Define
\begin{equation}
 \mathscr E(t,s)=\frac12
       \{(1-s^2)\mathscr K_0(t,s)+\mathscr K_1(t,s)\}.
 \label{eq:corr-energy-kernel}
\end{equation}
Then
\begin{equation}
 E(t,x)=\int_{\R}\mathscr E(t,s)\cos(2xs)\dd s.
 \label{eq:corr-energy}
\end{equation}
Writing
\(\mathscr A_0=\mathscr K_1\) and
\(\mathscr A_1=\mathscr K_2-s^2\mathscr K_1\), a prescribed
coefficient \(\kappa_j(t,x)\ge0\) satisfies the stationary loss
\(\Lag_j\ge-\kappa_jE\) if and only if
\begin{equation}
 \int_{\R}\{\mathscr A_j(t,s)+
       \kappa_j(t,x)\mathscr E(t,s)\}\cos(2xs)\dd s\ge0
       \quad\hbox{on }S_j=\{h_{j+1}=0\}.
 \label{eq:corr-stationary-target}
\end{equation}
\end{proposition}
\begin{proof}
The product expression for \(h_0^2+h_1^2\) has multiplier
\((1-uv)/4\) in the \(u,v\) coordinates.  The same substitution gives
\[
 E=\frac12\iint_{\R^2}
       (1-s^2+r^2)w_t(s+r)w_t(s-r)e^{2ixs}\dd r\dd s,
\]
which proves \eqref{eq:corr-energy}.  Subtracting the target
inequalities and using
\eqref{eq:corr-L0}--\eqref{eq:corr-L1} proves the equivalence.
\end{proof}

In \eqref{eq:corr-stationary-target}, stationarity and energy are
still those of the full physical function.  The coefficient must
be obtained independently of an unknown energy floor and have the
local boundedness required by \eqref{fc:stationary-loss}.
Because \(\kappa_j\) may depend on \(x\), this is generally not the
Fourier transform of one fixed kernel as \(x\) varies.

The positivity \(\mathscr K_1>0\) is not positive definiteness.
For a continuous integrable kernel, positive definiteness imposes
nonnegativity of its Fourier transform; pointwise positivity alone
does not.  In particular the associated-kernel criteria in
\cite{csordas} retain a substantive hypothesis.  Requiring a
positive-definite decomposition everywhere would be stronger than
the stationary bounded-loss target
\eqref{eq:corr-stationary-target}.  All complete source correlations
have been retained, but their oscillatory sign is still to be proved.

\begin{remark}[Gaussian normalization and a signed kernel]
\label{rem:corr-gaussian}
For the control \(w(u)=e^{-\alpha u^2}\), \(\alpha>0\),
\begin{align*}
 H(x)&=\frac{\sqrt\pi}{2\sqrt\alpha}e^{-x^2/(4\alpha)},&
 \mathscr K_0(s)&=\sqrt{\frac{\pi}{2\alpha}}e^{-2\alpha s^2},\\
 \mathscr K_1&=\frac{\mathscr K_0}{4\alpha},&
 \mathscr K_2&=\frac{3\mathscr K_0}{16\alpha^2}.
\end{align*}
The formulas above give
\[
 \Lag_0=\frac{\pi}{8\alpha^2}e^{-x^2/(2\alpha)}
       =\frac{H^2}{2\alpha},\qquad
 \Lag_1=\Lag_0\left(\frac1{2\alpha}
                              +\frac{x^2}{4\alpha^2}\right).
\]
Both agree with direct differentiation.  Nevertheless
\(\mathscr K_2-s^2\mathscr K_1
 =\mathscr K_1(3/(4\alpha)-s^2)\) changes sign.
Thus a demand that this correlation integrand be pointwise positive
would already reject this favorable control.
\end{remark}

\subsection{The existing staircase cutoff as a cross-correlation}

We now retain exactly the cutoff
\eqref{eq:signed-staircase-domain}, with \(a=\pi/16\),
integer \(N\ge8\), \(V_N=\frac12\log N\),
\(0\le t\le1/2\), and \(x\ge0\).  No new source truncation is
introduced.  In the original, rather than shared, integration
coordinate \(y=v-d_n\), define
\begin{equation}
 \psi_{a,N}(t,y)=
 \begin{cases}
 e^{t(y+ia)^2}
   \displaystyle\sum_{\substack{n\ge1\\y+d_n\le V_N}}
                  \phi_n(y+ia),&0\le y\le V_N,\\
 0,&\text{otherwise}.
 \end{cases}
 \label{eq:corr-cutoff-source}
\end{equation}
The staircase is precisely \(n\le N e^{-2y}\).
Endpoint values at its finitely many jumps do not affect any integral.
The source translation \eqref{eq:theta-translation} gives
\(W_\alpha=e^{t(y+ia)^2}\phi_n(y+ia)e^{ixy}\)
when \(y=v-d_n\).  Reversing that change of variables therefore gives
\begin{equation}
 h_k^{(N)}(t,x)
  =e^{-ax}\Re\int_0^{V_N}
         (-a+iy)^k\psi_{a,N}(t,y)e^{ixy}\dd y.
 \label{eq:corr-cutoff-jet}
\end{equation}
This reversal takes place in the retained, already shifted integral;
it does not shift individual theta summands across an unpaid contour.

\begin{proposition}[Full-amplitude hybrid cross-correlation]
\label{prop:corr-hybrid}
For \(j=0,1\), put
\begin{equation}
 \begin{split}
 \mathscr C_{j,N}^{a}(t,s)
 =-\int_{s-V_N}^{s}
 &[i(s+r)]^j[-a+i(s-r)]^{j+2}\\
 &\quad\cdot w_t(s+r)\psi_{a,N}(t,s-r)\dd r.
 \end{split}
 \label{eq:corr-hybrid-kernel}
\end{equation}
The full-amplitude observation, defined at every point by
\(\mathcal J_j^{(N)}=-h_jh_{j+2}^{(N)}\), has the exact representation
\begin{equation}
 \mathcal J_j^{(N)}(t,x)
   =e^{-ax}\Re\int_{\R}
          \mathscr C_{j,N}^{a}(t,s)e^{2ixs}\dd s.
 \label{eq:corr-hybrid-transform}
\end{equation}
On the full stationary set \(S_j\) it obeys
\begin{equation}
 \left|\mathcal J_j^{(N)}-\Lag_j\right|
 \le\sqrt E\,e^{-ax}T_{j+2}^{(N)}
 \le\mathsf D_0e^{-ax}\sqrt E\,e^{-cN^2},
 \label{eq:corr-hybrid-tail}
\end{equation}
with \(c,\mathsf D_0\) from \eqref{eq:signed-cutoff-constants}
and the complete omitted moments
\eqref{eq:signed-cutoff-tail-moments}.
\end{proposition}
\begin{proof}
Multiply the full-line formula for \(h_j\) by
\eqref{eq:corr-cutoff-jet}.  Since \(h_j\) is real, its full-line
integral can be placed inside the real part:
\[
 -h_jh_{j+2}^{(N)}
 =-\frac{e^{-ax}}2\Re\int_{\R}\int_0^{V_N}
       (iu)^j(-a+iy)^{j+2}
       w_t(u)\psi_{a,N}(t,y)e^{ix(u+y)}\dd y\dd u.
\]
The finite retained source is bounded on its compact integration
interval, and \(w_t\) has every absolute moment.  Thus Fubini and
the change of variables \(u=s+r,\ y=s-r\) apply.  The Jacobian is
two and the interval \(0\le y\le V_N\) becomes
\(s-V_N\le r\le s\), proving
\eqref{eq:corr-hybrid-transform}.

Before stationarity,
\[
 \mathcal J_j^{(N)}-\Lag_j
   =h_j(h_{j+2}-h_{j+2}^{(N)})-h_{j+1}^2.
\]
On \(S_j\) the square vanishes and \(|h_j|\le\sqrt E\).
The complete omitted-moment bound proves
\eqref{eq:corr-hybrid-tail}.
\end{proof}

The discontinuities and entering-index boundaries remain in
\(\psi_{a,N}\); the cross-correlation formula has not removed them.
Any integration by parts must include their boundary or distributional
terms.  There is also no replacement of the full stationary derivative
by its retained approximation.  At \(E=0\) both
\(\mathcal J_j^{(N)}\) vanish by their full-amplitude factorization,
independently of the sign of the retained source.
An all-\(N\) signed estimate for
\eqref{eq:corr-hybrid-transform} is therefore another way to address
the existing cutoff transfer, with exactly the same tail obligation.

\subsection{An optional common observation for the two stationary sets}

The two hybrid products can be combined if a common scalar observation
is useful.  Put
\begin{equation}
 \mathcal B_N=-h_0h_2^{(N)}-h_1h_3^{(N)},\qquad
 \varepsilon_k^{(N)}=e^{-ax}T_k^{(N)}.
 \label{eq:corr-common-cutoff}
\end{equation}
This is an optional bookkeeping device; the separate products can
give sharper estimates.

\begin{corollary}[Common cutoff and its full energy derivative]
\label{cor:corr-common-cutoff}
At every point in the preceding cutoff domain,
\begin{equation}
 \left|\mathcal B_N-\frac12E_t\right|
 \le\sqrt E\,
       \sqrt{(\varepsilon_2^{(N)})^2+
                         (\varepsilon_3^{(N)})^2}.
 \label{eq:corr-common-tail}
\end{equation}
Moreover,
\begin{align}
 \mathcal B_N&=\mathcal J_0^{(N)}
                 &&\text{on }S_0,
 \label{eq:corr-common-S0}\\
 \mathcal B_N-\mathcal J_1^{(N)}
    &=h_0(h_2-h_2^{(N)})
                 &&\text{on }S_1.
 \label{eq:corr-common-S1}
\end{align}
Consequently, for \(j=0,1\),
\begin{equation}
 |\mathcal B_N-\Lag_j|
 \le\sqrt2\,\mathsf D_0e^{-ax}\sqrt E\,e^{-cN^2}
                    \quad\text{on }S_j.
 \label{eq:corr-common-stationary-tail}
\end{equation}
\end{corollary}
\begin{proof}
Subtract \eqref{eq:stationary-energy-intro} before imposing
stationarity.  The difference is
\[
 h_0(h_2-h_2^{(N)})+h_1(h_3-h_3^{(N)}).
\]
Cauchy--Schwarz gives \eqref{eq:corr-common-tail}.
Setting \(h_1=0\) or \(h_2=0\) proves
\eqref{eq:corr-common-S0}--\eqref{eq:corr-common-S1}.
The complete Gaussian moment bounds and
\eqref{eq:stationary-energy-intro} then give
\eqref{eq:corr-common-stationary-tail}.
\end{proof}

In the associator notation of
Lemma~\ref{lem:hc-twisted-associators}, set
\(G_2^{(N)}=h_2^{(N)}+\jj h_3^{(N)}\) and
\(\mathbf u=\ell\jj\).  Then
\[
 -\frac12[F\ii,G_2^{(N)}\kk,\ell]=\mathbf u\,\mathcal B_N.
\]
The coefficient, rather than its absolute value, is still the
observation.  Both the scalar and octonionic forms vanish when \(E=0\).
On \(S_1\), the combined cutoff pays an additional second-derivative
tail which the separate \(\mathcal J_1^{(N)}\) avoids.

For example, if a new source estimate were to prove
\(\mathcal B_N\ge-A_jN^pE\) on \(S_j\) for every \(N\ge N_0\ge8\),
with the prescribed coefficients and local boundedness required by
Lemma~\ref{lem:polynomial-cutoff-transfer}, Young's inequality would
give
\[
 \Lag_j\ge-(A_jN^p+1/2)E
       -\mathsf D_0^2e^{-2ax}e^{-2cN^2}.
\]
This is within the existing conditional logarithmic-loss transfer.
It is not a proof of the proposed lower bound for \(\mathcal B_N\).
Neither the combined notation nor the correlation transform strengthens
the arithmetic input already available.

\subsection{What modular reflection and exceptional arithmetic would add}

For \(\Theta(y)=\sum_{n\in\mathbb Z}e^{-\pi n^2y}\), \(y>0\),
Poisson summation gives
\[
 \Theta(y)=y^{-1/2}\Theta(1/y),\qquad
 P(u)=e^u\Theta(e^{4u})=P(-u),\qquad
 \Phi=(P''-P)/16.
\]
The factor \(e^{tu^2}\) preserves this reflection.
In the complete pair \(w_t(s+r)w_t(s-r)\), inversion of either
theta variable amounts to exchanging \(s,r\), up to signs, after
the prefactors and differentiations are included.  That exchange is
already a consequence of evenness for any source.  It therefore
cannot by itself distinguish the theta family from the collision
controls of Proposition~\ref{prop:hc-positive-source-control}.

The two theta arguments \(e^{4(s+r)}\) and \(e^{4(s-r)}\) have
product \(e^{8s}\) and are reciprocal only when \(s=0\).
There is no automatic diagonalization of the full correlation.
Nor can the heat deformation be suppressed in a modular identity:
under \(y=e^{4u}\) it is
\(\exp\{t(\log y)^2/16\}\).
It is invariant under \(y\mapsto1/y\), but is not a fixed modular
weight.  With \(\tau=ie^{4u}\), even the modular translation
\(\tau\mapsto\tau+1\) changes the logarithmic variable to
\(\frac14\log(e^{4u}-i)\), on a consistent branch, and changes both
the heat multiplier and the Fourier phase.  The explicit parity identity and its localized form in
Section~\ref{upd:section} retain these factors and all differential
terms. Their separate global modes are not positive-heat integrable;
the required cancellations must precede integration.

Integral octonions offer genuine additional arithmetic, but through
a different source.  The Coxeter order has the \(E_8\) quadratic
lattice \cite[Section~1]{elkiesgross}; in the standard lattice
normalization its theta series is
\[
 \Theta_{E_8}(\tau)=E_4(\tau)
     =1+240\sum_{n\ge1}\sigma_3(n)e^{2\pi in\tau}.
\]
Termwise Mellin integration and \(n=dm\) in the divisor sum give
\begin{equation}
 \int_0^\infty(\Theta_{E_8}(iy)-1)y^{s-1}\dd y
  =240(2\pi)^{-s}\Gamma(s)\zeta(s)\zeta(s-3),
          \qquad\Re s>4.
 \label{eq:corr-E8-mellin}
\end{equation}
Thus replacing the scalar theta source by the octonionic lattice
theta series changes both its multiplicities and its Mellin
transform.  Positive representation counts do not by themselves
give the signed products in \eqref{eq:corr-stationary-target}.

Exceptional theta forms and their classical restrictions are
developed in \cite{kim93,pollackg2,pollackspin8,shan26}.
They justify investigating an arithmetic bridge, rather than
identifying the two sources merely because both involve theta
functions.  For example, \cite[Theorem~B]{shan26} proves nonvanishing
of the stated global lift for a cuspidal Hecke eigenform in
\(S_{2n+12}(\Gamma_0(N))\), \(n\ge0\), in its
\(\mathrm{PGL}_2\) setting.  Its Theorem~E spans
\(S_{2n+12}(\mathrm{SL}_2(\mathbb Z))\) by the specified weighted
exceptional theta series for \(n>0\).
These are genuine classical modular bridges, with their stated
weight and level restrictions.

For the present problem an arithmetic bridge must reproduce $H$, its
physical derivatives, or the complete correlations, and preserve the heat
deformation and full stationary constraints. Section~\ref{sf:e8-bridge}
supplies an exact eight-character restriction of an anisotropic $E_8$
theta series to the required two-coordinate source. Its finite sectors
are globally integrable with positive heat. The remaining step is a
signed estimate paid at the physical energy scale; separate sector norms
do not provide it.
The modular Fourier-sign constructions of \cite{viazovska} and the
modular inequalities of \cite{romik,lee24} give a useful model for
the last step: retain a specifically identified transformed expression
and prove its sign.  The polynomial/quasimodular inequalities of
\cite{lee26polynomial} and Fourier-eigenfunction estimates of
\cite{lee26positive} offer further tools when their hypotheses can
be verified.  None applies solely from the presence of the theta
source or the number eight.  The exact finite identification in Section~\ref{sf:e8-bridge} establishes no new
stationary sign estimate.

\subsection{Analytic tools and the signed physical observation}
\label{subsec:octonionic-analysis-literature}

The available octonionic function theories provide integral representations,
boundary projections, and regularity estimates for the fields above.
Their application to collision exclusion must preserve the distinguished
first pair and estimate the signed radial observation at its energy scale.
Proposition~\ref{prop:hc-positive-source-control} supplies a control family
against which any claimed consequence of generic monogenicity can be tested.

For differential calculations there is an associative real-matrix
realization.  Let $L_a v=av$ on $\OO\cong\R^8$, and let
$\epsilon_1,\ldots,\epsilon_7$ be an orthonormal basis of imaginary
octonions. Alternativity gives
\[
 L_{\epsilon_i}L_{\epsilon_j}+L_{\epsilon_j}L_{\epsilon_i}=-2\delta_{ij}I
 \qquad(1\le i,j\le7),
\]
so the imaginary left-multiplication operators form a Clifford
representation; see \cite[Section~2.3]{baez}.  With the declared basis,
$\mathcal D_8$ is therefore a real matrix Dirac operator and its conjugate
composition is $\Delta_8 I$.  This is useful for elliptic estimates and
symbolic checks.  It is not a multiplicative embedding of the octonion
algebra: generally $L_{ab}\ne L_aL_b$.  The actual associators and their
orientation must still be retained in the physical observations.
The component conventions are treated in \cite{kauhanenorelma}.

The Cauchy formulas of Li--Peng \cite{lipeng} depend on parenthesization.
Colombo--Krau\ss har--Sabadini \cite[Proposition~2.2 and Sections~3--4]{colombokraussharsabadini} distinguish monogenic and slice-monogenic
reproducing spaces and construct the relevant Hardy/Bergman framework
using real parts of appropriately weighted para-linear inner products.
Constales--Krau\ss har \cite{constaleskrausshar} develop boundary projection
methods in the same nonassociative setting.  The Teodorescu and
projection calculus of \cite{teodorescu24} includes the necessary
associator compensation and can be useful when an approximate field has
a quantified monogenic defect.  These tools provide norm and evaluation
estimates.  To yield the present criterion, a boundary or defect estimate
must additionally control the adverse radial component relative to $E$;
an absolute error does not give such control near a candidate zero.
The present lifts are already exactly monogenic, so an integral
reconstruction needs a specific further estimate to justify its use.

The regularity class is another substantive hypothesis.  Octonionic
slice regularity \cite{gentilistruppa}, quaternionic-slice Fueter
regularity \cite{ghiloni}, and the partial-slice theory of
\cite{xusabadini26} have distinct defining equations.  The latter's
Cauchy--Pompeiu, maximum-modulus and extension results organize these
classes but do not imply that a distinguished two-component trace of a
nonzero field is nonvanishing.  The integral Fueter mapping in
\cite{fuetermapping} transfers slice data to monogenic functions in the
Clifford/quaternionic setting; such a transfer changes the field and
requires preservation or reconstruction of the collision observation.
Differential transforms can annihilate low-degree collision controls,
so the transformed norm cannot be substituted for $E$ without that step.

Refined Kato inequalities \cite{feehan} concern norm gradients in the
ambient variables.  Frequency and nodal-set results, including
\cite{borrelliwu26}, describe the structure and dimension of possible
Dirac zeros; they do not exclude an intersection with the physical
slice.  Moreover, an ordinary double collision has $E=0$ and $h_2\ne0$,
so the doubled four-jet field need not vanish there.  The explicit
ambient log-energy calculation above identifies the stronger property
of this particular holomorphic lift, while also showing the loss of
orientation in the resulting squared determinants.

The phase/Wronskian description similarly connects the physical signs
with familiar complex-analytic positivity questions.  A Hermite--Biehler
or reproducing-kernel hypothesis must nevertheless be proved for the
specific source.  Conrey--Li \cite{conreyli} give counterexamples to
particular proposed de Branges shift-positivity conditions associated
with zeta and Dirichlet $L$-functions; their result concerns those
conditions, not every de Branges approach.  Here the operative next step
is a complete theta identity or factorization yielding the declared
signed energy-relative estimate, with both stationary conditions and
all cutoff payments preserved.

% Integration of research Notes 46--48, 50--51, reviewed in Reviews 29--31.
% Prepared for Edward Baker with substantial LLM assistance.
% Recorded model gpt-6-astra; reasoning effort ultra; 15 October 2026.
\section{Signed boundary currents and integrated positivity}
\label{sf:section}

The preceding identities retain the physical observation.  We now prove
boundary and averaged estimates for that observation, and give an exact
finite \(E_8\)-character realization of its complete source.  The
arithmetic energy payment in the unresolved region remains a separate
hypothesis.  In particular, every stationary set below is
\(S_j=\{h_{j+1}=0\}\) for the full field, and every relative loss is paid
by \(E=h_0^2+h_1^2\).

\subsection{A spectral-circle current and its exact defect}

Fix real \(t,x\), \(R>0\), and write
\[
 z_\theta=x+Re^{i\theta},\qquad
 \langle v\rangle=\frac1{2\pi}\int_0^{2\pi}v(\theta)\dd\theta,\qquad
 q(\theta)=1-\cos2\theta.
\]
For a physical quaternion \(a=A+\jj B\) and an arbitrary
\(b=C+iD+\jj E+\kk F\), the multiplication convention of
Lemma~\ref{lem:hc-twisted-associators} gives
\[
 \frac12\langle\mathbf u,[a\ii,b\kk,\ell]\rangle=AC+BE
                 =\Re(\bar a b),\qquad\mathbf u=\ell\jj .
\]
Indeed, the \(\jj\)-coefficient of \([a\ii,b\kk]\) is
\(-2(AC+BE)\).  Other off-slice components of the associator need not
vanish; only the stated projection is being observed.

\begin{proposition}[Exact circle reconstruction and negative defect]
\label{sf:boundary}
Define the combined and separate currents by
\[
 \varphi(\theta)=
 \Re\{h_0(t,x)H(t,z_\theta)+h_1(t,x)H_z(t,z_\theta)\},\qquad
 \varphi_j(\theta)=h_j(t,x)\Re h_j(t,z_\theta).
\]
Then, at every real point,
\begin{equation}
 \langle\varphi\rangle=E,\qquad
 \frac12E_t=-\frac4{R^2}E+\frac4{R^2}\langle q\varphi\rangle .
 \label{sf:combined-circle}
\end{equation}
For \(j=0,1\), on \(S_j\),
\begin{equation}
 \langle\varphi_j\rangle=h_j^2,\qquad
 \Lag_j=-\frac4{R^2}h_j^2+
                      \frac4{R^2}\langle q\varphi_j\rangle .
 \label{sf:stationary-circle}
\end{equation}
With \(v_-=\max(0,-v)\), set
\(\mathfrak D_R=\langle q\varphi_-\rangle\) and
\(\mathfrak D_{j,R}=\langle q(\varphi_j)_-\rangle\).  Consequently
\begin{equation}
 \frac12E_t\ge-\frac4{R^2}(E+\mathfrak D_R),\qquad
 \Lag_j\ge-\frac4{R^2}(h_j^2+\mathfrak D_{j,R})
                  \quad\hbox{on }S_j.
 \label{sf:boundary-defect}
\end{equation}
\end{proposition}
\begin{proof}
If \(F\) is holomorphic on the disk, its Taylor series gives
\[
 \langle\Re F(z_\theta)\rangle=\Re F(x),\qquad
 \langle\cos2\theta\,\Re F(z_\theta)\rangle
                         =\frac{R^2}{4}\Re F''(x).
\]
Apply this first to \(h_0H+h_1H_z\), whose second derivative is
\(-E_t/2\), and then to \(h_jh_j(t,\cdot)\).  On \(S_j\) the latter
has second derivative \(h_jh_{j+2}=-\Lag_j\).  The inequalities follow
from \(q\ge0\).  No division by \(E\) is used.
\end{proof}

The separate currents are precisely the projected associators with
\(a=G_j(t,x)\), \(b=G_j(t,z_\theta)\) on \(S_j\).  On \(S_1\), the
combined current has an extra \(h_0\Re H\) term with zero second Fourier
coefficient, but possibly a larger negative part.  Thus the separate
current can yield a sharper sufficient hypothesis.

A source-prescribed bound
\begin{equation}
 \mathfrak D_{j,R(t,x)}\le\eta_j(t,x)E
 \quad\Longrightarrow\quad
 \Lag_j\ge-\frac{4(1+\eta_j)}{R(t,x)^2}E
 \label{sf:defect-target}
\end{equation}
would provide the desired relative loss.  Here \(\eta_j\ge0\) and
\(R,1/R,\eta_j\) must be locally bounded through the covered candidate
events, on a domain containing the required predecessor neighborhoods.
The weaker signed requirement
\(\langle q\varphi_j\rangle\ge-C_jE\) also suffices.  A radius allowed
to collapse with unknown \(E\) does not give a locally bounded
coefficient.  The combined defect may be used similarly on both sets.

All source coordinates remain present in these currents.  With
\(w_t(u)=e^{tu^2}\Phi(u)\), \(a_\theta=x+R\cos\theta\), and
\(b_\theta=R\sin\theta\),
\begin{align}
 \varphi(\theta)&=\int_0^\infty w_t(u)\cosh(b_\theta u)
       \{h_0\cos(a_\theta u)-h_1u\sin(a_\theta u)\}\dd u,
       \label{sf:source-current}\\
 \varphi_0(\theta)&=h_0\int_0^\infty w_t(u)
                   \cosh(b_\theta u)\cos(a_\theta u)\dd u,\nonumber\\
 \varphi_1(\theta)&=-h_1\int_0^\infty u w_t(u)
                   \cosh(b_\theta u)\sin(a_\theta u)\dd u.\nonumber
\end{align}
The double-exponential source bound gives compact-uniform convergence
in \(t,x,R\).  The negative part must be taken after the complete source
is summed and integrated.  Absolute source moments give only
\[
 \mathfrak D_{j,R}\le |h_j|
       \int_0^\infty u^j w_t(u)\cosh(Ru)\dd u,
\]
which has the \(\sqrt E\) scale, not the \(E\) scale in
\eqref{sf:defect-target}.

The constant four is optimal under nonnegative boundary current, even
at a stationary center.  At a fixed time the polynomial
\[
 P_n(z)=1+2\sum_{k=1}^n\left(1-\frac{k}{n+1}\right)(z/R)^{2k}
\]
has \(P_n(0)=1\), \(P_n'(0)=0\), and nonnegative boundary real part,
the Fej\'er kernel at angle \(2\theta\).  Its signed value is
\(-P_n(0)P_n''(0)=-4n/((n+1)R^2)\); its polynomial backward-heat
extension is exact.  For the quadratic collision
\(H=z^2-2(t-T)\), at \(t=T-\delta,x=0\),
\[
 E=4\delta^2,\qquad
 \varphi_0=2\delta(2\delta+R^2\cos2\theta),\qquad
 \frac{\mathfrak D_{0,R}}E\sim
 \frac{R^2}{2\delta}\left(\frac1\pi+\frac14\right).
\]
The cubic \(H=z^3-6(t-T)z\) gives the identical defect ratio on
\(S_1\).  These follow from
\(\langle(-\cos2\theta)_+\rangle=1/\pi\) and
\(\langle\cos^22\theta\,1_{\{\cos2\theta<0\}}\rangle=1/4\).
Conversely, the slice \(H=\varepsilon-z^2\) has favorable
\(\Lag_0=2\varepsilon\) at zero while
\(\mathfrak D_{0,R}/E\to\infty\).  An eight-dimensional harmonic-ball reconstruction under a nonnegative
scalar-boundary hypothesis gives the weaker constant
\((8+2)(8-1)/R^2=70/R^2\); enlarging the ambient domain does not
improve the spectral-circle observation.  Thus payment of all negative mass
is a strictly stronger sufficient mechanism than the signed bound;
it can discard favorable cancellation.

\subsection{The original staircase on a complex circle}

The cutoff in \eqref{eq:corr-cutoff-source} has an entire spectral
extension.  More generally fix \(0<a<\pi/8\), \(0\le t\le T<\infty\),
\(V_N=\frac12\log N\), and retain exactly
\(n\le Ne^{-2y}\), \(0\le y\le V_N\).  Define
\begin{equation}
 \mathscr H_N(t,z)=\frac12\int_0^{V_N}
 \left\{\psi_{a,N}(t,y)e^{(-a+iy)z}
       +\overline{\psi_{a,N}(t,y)}e^{(-a-iy)z}\right\}\dd y ,
 \label{sf:entire-cutoff}
\end{equation}
where \(\psi_{a,N}=e^{t(y+ia)^2}\sum_{n\le Ne^{-2y}}\phi_n(y+ia)\);
conjugation refers to real \(t\).  Its derivatives on the real axis
are the retained jets in \eqref{eq:corr-cutoff-jet}.  The complete
version equals \(H(t,z)\), first by the complete contour identity on
the real axis and then by entireness.  The complete-source vertical
cancellation is used before separating theta indices.

\begin{proposition}[Complete Gaussian circle payment]
\label{sf:circle-tail}
Put \(\beta=\pi\cos4a>0\), \(c_a=\beta/2\), and
\begin{equation}
 Q_j(T,R,a)=\int_0^\infty(a+y)^je^{Ty^2+Ry}
 \sum_{n\ge1}(2\pi^2n^4e^{9y}+3\pi n^2e^{5y})
                     e^{-c_an^2e^{4y}}\dd y .
 \label{sf:Q}
\end{equation}
These constants are finite and locally bounded for finite \(T,R\)
and \(a\) in the open strip.  Uniformly on \(|z-x|\le R\),
\begin{equation}
 |h_j(t,z)-\partial_z^j\mathscr H_N(t,z)|
 \le e^{-a(x-R)}Q_j(T,R,a)e^{-c_aN^2}
 =:\epsilon_{j,N}.
 \label{sf:circle-tail-bound}
\end{equation}
Retaining the full real factor \(h_j(t,x)\), set
\[
 \mathfrak D_{j,R,N}
 =\left\langle q
 \{h_j(t,x)\Re\partial_z^j\mathscr H_N(t,z_\theta)\}_-\right\rangle .
\]
On the full stationary set,
\begin{equation}
 |\mathfrak D_{j,R}-\mathfrak D_{j,R,N}|
                       \le\sqrt E\,\epsilon_{j,N}.
 \label{sf:defect-stability}
\end{equation}
\end{proposition}
\begin{proof}
The omitted domain is exactly \(n^2e^{4y}>N^2\).
Extract \(e^{-c_aN^2}\) from \(e^{-\beta n^2e^{4y}}\);
the remaining exponential is that in \eqref{sf:Q}.  Both conjugate
branches cost at most \(e^{-a(x-R)+Ry}\), a spectral derivative
costs \((a+y)^j\), and the heat factor costs \(e^{Ty^2}\).
Their factors \(1/2\) pay the two equal bounds.  For finiteness, use
\[
 e^{-c_an^2e^{4y}}
 \le e^{-(c_a/2)n^2}e^{-(c_a/2)e^{4y}}.
\]
The resulting theta-index sums and remaining integral converge
locally uniformly.  Finally \(v\mapsto v_-\) is 1-Lipschitz,
\(\langle q\rangle=1\), and \(|h_j|\le\sqrt E\) for \(j=0,1\).
Only spectral derivatives were taken, so no staircase edge was
differentiated or discarded.
\end{proof}

This yields a quantitative conditional route into the cutoff transfer.
Suppose, for every integer \(N\ge N_{0,j}\ge8\), one proves
\[
 \mathfrak D_{j,R,N}\le
 \alpha_j\mathfrak C_{j,N}E+B_je^{-d_jN^2},\qquad
 \mathfrak C_{j,N}=N^pe^{\sigma_j(\log N)^2},
\]
where \(p\ge0\) is fixed, \(\alpha_j,B_j,\sigma_j\ge0\), \(d_j>0\),
and the coefficients, inverse rates, \(N_{0,j},R,1/R\) are
source-prescribed and locally bounded through candidate events.
Young's inequality and \(\mathfrak C_{j,N}\ge1\) give
\begin{equation}
 \begin{split}
 \Lag_j\ge{}&
 -\frac4{R^2}(\alpha_j+\tfrac32)\mathfrak C_{j,N}E\\
 &-\frac4{R^2}\left\{B_j+
             \tfrac12e^{-2a(x-R)}Q_j(T,R,a)^2\right\}
                         e^{-\min(d_j,2c_a)N^2}.
 \end{split}
 \label{sf:conditional-boundary-cutoff}
\end{equation}
For \(\sigma_j=0\) this has the polynomial transfer form.
The additional factor \(e^{\sigma_j(\log N)^2}\), after optimizing
a Gaussian tail, has subpower growth in \(1/E\) and is also covered
by the corresponding analytic-rigidity transfer.  The original
manuscript cutoff is recovered by \(a=\pi/16\), \(T=1/2\).
The circle tail and its payment are proved here; the all-cutoff signed
hypothesis is not.

\subsection{Sharp translated averages and source certificates}

For a real \(f=h_j(t,\cdot)\) and a positive \(C^1\) weight \(\chi\),
direct differentiation gives
\begin{equation}
 \chi\Lag_j+\frac{\chi'^2}{8\chi}f^2+
       \partial_x(\chi ff')
 =2\chi\left(f'+\frac{\chi'}{4\chi}f\right)^2.
 \label{sf:spatial-flux}
\end{equation}
Thus, if the endpoint flux vanishes and
\(|(\log\chi)'|\le a\), then
\(\int\chi\Lag_j\ge-(a^2/8)\int\chi h_j^2\).
On an interval with \(ff'=0\) at both endpoints the unweighted
identity is \(\int\Lag_j=2\int h_{j+1}^2\ge0\).
Multiplying \eqref{sf:spatial-flux} by \(2\mathbf u\) is an exact
oriented associator identity; the derivative remains a signed
spatial flux.

\begin{theorem}[Sharp hyperbolic-secant average]
\label{sf:sech}
For any \(a>0\), \(x_0\in\R\), set
\[
 \chi(x)=\operatorname{sech}(a(x-x_0)),\qquad
 \nu=\frac{\sqrt2-1}{2},\qquad c_*=2\nu^2=\frac{3-2\sqrt2}{2}.
\]
For the complete theta field, at every real time,
\begin{equation}
 \int_{\R}\chi\Lag_j\dd x
       \ge-c_*a^2\int_{\R}\chi h_j^2\dd x,\qquad
 \int_{\R}\chi(\Lag_0+\Lag_1)\dd x
       \ge-c_*a^2\int_{\R}\chi E\dd x .
 \label{sf:sech-bound}
\end{equation}
The constant \(c_*\) is sharp in the general weighted Sobolev class.
\end{theorem}
\begin{proof}
For \(g=\sqrt\chi f\), integration by parts and
\(\nu(\nu+1)=1/4\) give
\[
 \int\chi(f'^2-ff'')=
 2\int g'^2+\frac12\int(\log\chi)''g^2,
\]
and therefore
\begin{equation}
 \int\chi\Lag_j+c_*a^2\int\chi f^2
   =2\int\{g'+\nu a\tanh(a(x-x_0))g\}^2\ge0 .
 \label{sf:sech-square}
\end{equation}
Adding the two branches uses \(h_0^2+h_1^2=E\).
Equality in the general class holds for
\(g=\operatorname{sech}^{\nu}(a(x-x_0))\),
\(f=\cosh^{1-1/\sqrt2}(a(x-x_0))\), and compactly supported smooth
truncations approach equality.  For theta jets, a fixed complete-source
contour shift with \(0<a_0<\pi/8\) gives
\(|h_k(t,x)|\le C_{k,a_0}e^{-a_0|x|}\) on compact real-time
intervals.  The source tail pays the heat and polynomial factors.
All integrations converge and their endpoint fluxes vanish.
\end{proof}
Each individual bound also has right side at least
\(-c_*a^2\int\chi E\).  Sharpness is not asserted within the theta
subfamily.  Narrowing the weight to width \(1/a\) costs \(a^2\);
the theorem therefore gives no locally bounded pointwise coefficient.

There is a useful source target which does not penalize angular
motion.  Identify \(F=h_0+\jj h_1\) with \((h_0,h_1)\), and let
\(J(A,B)=(-B,A)\).  If \(M=S+\omega J\), \(S=S^*\), then
\(F\cdot MF=F\cdot SF\).  Hence, on \(S_j\), an independently proved
identity
\[
 F_t=MF+c_jh_{j+1}+R_j,\qquad
 S\succeq-\kappa_jI,\qquad F\cdot R_j\ge-b_jE
\]
would give \(\Lag_j\ge-(\kappa_j+b_j)E\).
The identity must be defined on its domain, but the magnitude of
\(\omega\) need not enter the loss.  A full residual-norm bound is
a stronger condition than the stated signed pairing.

For explicit source certificates put \(w=e^{tu^2}\Phi\) on
\(u\ge0\), and define
\[
 w\mathsf T_wA=(wA)',\qquad
 \mathsf T_wA=A'+(2tu+\Phi'/\Phi)A .
\]
Require \(wA\to0\) at both endpoints.  If prescribed real
coefficients and fields satisfy
\begin{align}
 (u^2-a_0)\cos(xu)+b_0u\sin(xu)
       &=\mathsf T_wA_0,\label{sf:certificate0}\\
 -u^3\sin(xu)-a_1\cos(xu)+b_1u\sin(xu)+c_1u^2\cos(xu)
       &=\mathsf T_wA_1,\label{sf:certificate1}
\end{align}
then integration gives
\[
 -h_2=a_0h_0+b_0h_1,\qquad
 -h_3=a_1h_0+b_1h_1+c_1h_2.
\]
On the respective full stationary sets this proves
\begin{equation}
 \Lag_0=a_0E,\qquad
 \Lag_1=a_1h_0h_1+b_1h_1^2
 \ge-\frac{\sqrt{a_1^2+b_1^2}-b_1}{2}E .
 \label{sf:certificate-losses}
\end{equation}
The second constant is the negative of the smallest eigenvalue of
\(\left(\begin{smallmatrix}0&a_1/2\\a_1/2&b_1\end{smallmatrix}\right)\).
An arithmetic certificate should prescribe \(A_j\) as a finite
differential expression in complete theta data and phase parameters,
with paid endpoint and residual terms.  Defining it by an incomplete
integral merely hides the unknown cancellation at an endpoint.
Unknown physical quotients such as \(-h_2/h_0\) do not prescribe
locally bounded coefficients.  No theta certificate
\eqref{sf:certificate0}--\eqref{sf:certificate1} is established here.

Frequency-dependent angular motion has an additional exact term.
Set
\[
 v_x(u)=(\cos xu,-u\sin xu),\qquad
 Q_\omega=\int_0^\infty w(u)\omega(u)Jv_x(u)\dd u.
\]
Symmetrization yields
\begin{equation}
 F\cdot Q_\omega=\frac12\iint_{[0,\infty)^2}
 w(u)w(v)\{\omega(u)-\omega(v)\}
 \{u\sin(xu)\cos(xv)-v\cos(xu)\sin(xv)\}\dd u\dd v .
 \label{sf:angular-commutator}
\end{equation}
For \(x>0\), a source supported in \(0\le u\le\pi/(2x)\), and
nondecreasing \(\omega\), this is nonnegative: the displayed
determinant has the sign of \(u-v\), as follows from monotonicity of
\(u\tan(xu)\) in the interior and continuity at the endpoints.
The full theta source has unbounded support, and its exterior cannot
be omitted or paid at energy scale by an absolute moment.  Thus this
compact-support observation does not prove the theta target.

\subsection{An exact finite exceptional-lattice source bridge}
\label{sf:e8-bridge}

Use the coordinate model
\[
 L=\{\xi\in\mathbb Z^8:\textstyle\sum_i\xi_i\in2\mathbb Z\}
 \ \cup\
 \{\xi\in(\mathbb Z+\tfrac12)^8:
                       \textstyle\sum_i\xi_i\in2\mathbb Z\}
\]
of \(E_8\).  For independent \(u=s+r,v=s-r\) put
\[
 y_+=e^{4u},\quad y_-=e^{4v},\quad
 Q_\Lambda=\operatorname{diag}(y_+/4,y_-/4,\Lambda,\ldots,\Lambda),
 \quad \Lambda>0,
\]
and define the eight twisted theta series
\[
 T_{a,b}^{\Lambda}(u,v)=
 \sum_{\xi\in L}e^{-\pi\xi^TQ_\Lambda\xi}
             e^{\pi i[(a-b)\xi_1+b\xi_2]},
 \qquad 0\le a\le3,\quad 0\le b\le1.
\]
Here and below
\(\Theta_3(y)=\sum_ne^{-\pi n^2y}\),
\(\Theta_4(y)=\sum_n(-1)^ne^{-\pi n^2y}\), and
\(\Theta_2(y)=\sum_ne^{-\pi(n+1/2)^2y}\).

\begin{proposition}[Exact index-eight character projection]
\label{sf:E8-projection}
With
\(A_6(\Lambda)=(\Theta_3(\Lambda)^6+\Theta_4(\Lambda)^6)/2>0\),
\begin{equation}
 \frac18\sum_{a=0}^3\sum_{b=0}^1T_{a,b}^{\Lambda}(u,v)
       =A_6(\Lambda)\Theta_3(y_+)\Theta_3(y_-).
 \label{sf:E8-theta}
\end{equation}
Consequently, for
\[
 \mathcal D=(\partial_s^2-\partial_r^2)^2
                    -8(\partial_s^2+\partial_r^2)+16,
\]
the complete physical product is
\begin{equation}
 \Phi(s+r)\Phi(s-r)=\frac1{4096A_6(\Lambda)}
 \mathcal D\left[\frac{e^{2s}}8\sum_{a,b}
                      T_{a,b}^{\Lambda}(s+r,s-r)\right].
 \label{sf:E8-source}
\end{equation}
Every differentiated character sector is a finite linear combination
of products of translated complete theta sources.  Its full
polynomially weighted heat transform is absolutely convergent on
compact real-time and complex-spectral parameter sets.
\end{proposition}
\begin{proof}
The character average
\[
 \frac18\sum_{a=0}^3\sum_{b=0}^1
        e^{\pi i[a\xi_1+b(\xi_2-\xi_1)]}
\]
selects \(\xi_1\in2\mathbb Z\) and then \(\xi_2\in2\mathbb Z\).
It is the indicator of
\(M=2\mathbb Z^2\oplus D_6\), where
\(D_6=\{n\in\mathbb Z^6:\sum n_i\in2\mathbb Z\}\).
This sublattice has index eight.  Summing the selected coordinates
proves \eqref{sf:E8-theta}.  Since
\((\partial_u^2-1)(\partial_v^2-1)=\mathcal D/16\) and
\(\Phi=(P''-P)/16\), \(P(u)=e^u\Theta_3(e^{4u})\), the source
normalization is \(16\cdot16^2=4096\).

For the sector assertion set \(c=(\log2)/2\),
\(P_\nu(u)=e^u\Theta_\nu(e^{4u})\),
\(\Phi_\nu=(P_\nu''-P_\nu)/16\), and
\(\widetilde\Phi_2=(\partial_u^2-1)
[e^u\Theta_2(e^{4u}/4)]/16\).  Splitting even and odd indices gives
\begin{equation}
 \begin{split}
 \Phi_2(u)&=e^c\Phi(u-c)-\Phi(u),\\
 \Phi_4(u)&=e^c\Phi(u+c)-\Phi(u),\\
 \widetilde\Phi_2(u)&=e^{2c}\Phi(u-2c)-e^c\Phi(u-c).
 \end{split}
 \label{sf:duplication}
\end{equation}
For completeness an explicit formula for each character is obtained
by writing \(\Theta_{[m]}=\Theta_3\) for even \(m\), \(\Theta_4\)
for odd \(m\), \(r_0=a-b\), and \(C_6=\Theta_2(\Lambda)^6/2\):
\begin{align*}
 T_{a,b}^{\Lambda,\mathrm{int}}
 &=\tfrac12\{
 \Theta_{[r_0]}(y_+/4)\Theta_{[b]}(y_-/4)\Theta_3(\Lambda)^6\\
 &\hspace{18mm}+
 \Theta_{[r_0+1]}(y_+/4)\Theta_{[b+1]}(y_-/4)
                                      \Theta_4(\Lambda)^6\},\\
 T_{a,b}^{\Lambda,\mathrm{half}}
 &=1_{\{b=0,\ a\ {\rm even}\}}(-1)^{a/2}
                     C_6\Theta_2(y_+/4)\Theta_2(y_-/4).
\end{align*}
The parity-weighted half-integer sums vanish under \(n\mapsto-n-1\);
even phases give the displayed sign.  Applying \((\partial_u^2-1)/16\) to
\(e^u\Theta_3(e^{4u}/4)\) gives \(e^c\Phi(u-c)\); applying the
same operator to \(e^u\Theta_4(e^{4u}/4)\) gives
\(e^{2c}\Phi(u)-e^c\Phi(u-c)\).  Together with
\eqref{sf:duplication} these prove the finite translated-source
description.  Each translate retains double-exponential decay in
both directions, paying every polynomial, heat factor and bounded
complex spectral displacement.  The eight-sector sum is finite,
so no global infinite Poisson-mode interchange is needed.
\end{proof}

The bridge is exact at every \(\Lambda>0\) and preserves both
independent physical coordinates.  In the correlations its heat
multiplier is still \(e^{2t(s^2+r^2)}\).  Nevertheless the projector
does not preserve exceptional mixing of the physical plane.  The
norm-two roots of \(M\) are exactly the \(D_6\) roots; their span is
the transverse six-plane.  Any lattice isometry preserving \(M\)
preserves this span and its orthogonal physical plane, acting there
by signed coordinate permutations.  More general transformations
must carry both the anisotropic form and the projector with them.
Coefficientwise positivity of a selected theta series does not
establish positivity after \(\mathcal D\).

Indeed the untwisted sector has an explicit adverse region.  Put
\(B_6=(\Theta_3(\Lambda)^6-\Theta_4(\Lambda)^6)/2\).  Splitting the
integer parity and half-integer sectors gives
\begin{align}
 T_{0,0}^{\Lambda}={}&
 A_6\{\Theta_3(y_+)\Theta_3(y_-)+\Theta_2(y_+)\Theta_2(y_-)\}
 \nonumber\\
 &+B_6\{\Theta_3(y_+)\Theta_2(y_-)+
                    \Theta_2(y_+)\Theta_3(y_-)\}
       +C_6\Theta_2(y_+/4)\Theta_2(y_-/4).
 \label{sf:E8-untwisted}
\end{align}
After multiplying by \(e^{2s}\) and applying \(\mathcal D/4096\),
this is the same combination of \(\Phi,\Phi_2,\widetilde\Phi_2\)
products.  Modular reflection gives \(\Phi_2(-u)=\Phi_4(u)\).
For \(u\ge0\),
\[
 \Phi_4(u)=2e^u\sum_{n\ge1}(-1)^n
 z_n(z_n-\tfrac32)e^{-z_n}<0,\qquad z_n=\pi n^2e^{4u}.
\]
The inequality follows because \(z(z-\tfrac32)e^{-z}\) is positive
and strictly decreasing for \(z>3\); its derivative is
\(-e^{-z}(z-3)(z-\tfrac12)\).  In contrast \(\Phi_2(u)>0\) when
\(e^{4u}\ge6/\pi\).
As \(R\to\infty\),
\[
 \Phi(R)\sim2\pi^2e^{9R-\pi e^{4R}},\quad
 \Phi_2(R)\sim\frac{\pi^2}{8}e^{9R-(\pi/4)e^{4R}},\quad
 \Phi_4(R)\sim-2\pi^2e^{9R-\pi e^{4R}}.
\]
First let \(\Lambda\to\infty\), so \(A_6\to1\), \(B_6,C_6\to0\).
At \((u,v)=(R,-R)\) the differentiated untwisted expression becomes
\begin{equation}
 \Phi(R)^2+\Phi_2(R)\Phi_4(R)
 \sim-\frac{\pi^4}{4}e^{18R-(5\pi/4)e^{4R}}<0,
 \label{sf:E8-negative}
\end{equation}
and the adverse term divided by \(\Phi(R)^2\) is asymptotic to
\(-\frac1{16}\exp\{(3\pi/4)e^{4R}\}\).
For each sufficiently large fixed \(R\), continuity gives the same
negativity for sufficiently large finite \(\Lambda\).  This order
of limits proves a pointwise signed obstruction to untwisted
positivity, not a failure of heat integrability.

Translation also identifies exactly which auxiliary fields a
sector norm would observe.  Define
\[
 \mathcal T_cH(t,z)=e^{c+tc^2+icz}H(t,z-2itc).
\]
Changing the source variable in \eqref{sf:duplication} gives
\begin{equation}
 \int_{\R}e^{tu^2}\Phi_2(u)e^{izu}\dd u
          =2\{\mathcal T_cH(t,z)-H(t,z)\}.
 \label{sf:heat-translation}
\end{equation}
At \(t=0\) the result is proportional to \(H(0,z)\).  At positive
time its imaginary displacement is indispensable.  At a possible
double collision the displaced value contains the complete series
\[
 H(t,x-2itc)=\sum_{k\ge2}\frac{(-2itc)^k}{k!}h_k(t,x);
\]
it need not vanish with \(E\).

This is an exact obstruction to a universal auxiliary-norm payment.
For fixed \(t_*>0\), take
\(\Phi_{\rm ctrl}(u)=e^{-(1+t_*)u^2}(u^2-\frac34)^2\).  Gaussian
differentiation gives
\[
 H_{\rm ctrl}(t_*,z)=\frac{\sqrt\pi}{32}
                          e^{-z^2/4}(z^2-3)^2 .
\]
At \(x_*=\sqrt3\), \(H=H_z=0\), \(h_2\ne0\), whereas
\(H(t_*,x_*-2it_*c)\ne0\).  Replacing the square in the source by
\((u^2-\frac34)^2+\varepsilon\) makes it strictly positive.
The Jacobian of \((H,H_z)\) in \((t,x)\) has determinant
\(-h_2^2\ne0\), so for every sufficiently small fixed
\(\varepsilon>0\) the real implicit-function theorem preserves a
positive-time double collision and the nonzero shifted amplitude.
Thus no finite coefficient can bound
\(|\mathcal T_cH-H|^2\) by \(E\) at that collision, even for that
fixed strictly positive source.  This Gaussian control is locally
heat-analytic; it is not claimed to have the theta tail envelope.

Finally define \(\mathscr F_\Lambda\) by the right side of
\eqref{sf:E8-source}, and put
\(q_0=r^2\), \(q_1=r^4-s^2r^2\),
\(b=(1-s^2+r^2)/2\).  The exact source gate is
\begin{equation}
 \iint_{\R^2}e^{2t(s^2+r^2)}
       \{q_j+\kappa_jb\}\cos(2xs)\mathscr F_\Lambda(s,r)
                         \dd s\dd r\ge0\quad\hbox{on }S_j.
 \label{sf:E8-gate}
\end{equation}
By the complete-correlation formulas this is precisely
\(\Lag_j+\kappa_jE\ge0\).  The boundary current likewise has the
complete two-source formula
\[
 \varphi_j(\theta)=\frac14\Re\iint_{\R^2}
 (iu)^j(iv)^je^{t(u^2+v^2)}e^{iz_\theta u+ixv}
                              \Phi(u)\Phi(v)\dd u\dd v .
\]
The eight-character packet can be inserted into this formula
exactly, but must be reconstructed before a negative-part estimate.
The bridge proves the identification required for an exceptional
arithmetic approach; it does not prove either signed gate.

\subsection{Closed self-dual forms and a sharp disk completion}

Unpublished quaternion-form working notes of Edward B.~Baker~III
dated 21 January and 14 February 2014 motivate retaining a polarized
form pairing.  They were circulated as incomplete exploratory notes.
The construction here uses direct exterior algebra and Stokes'
theorem; it assumes none of their unfinished singular-kernel or
nonabelian boundary formulas.

On oriented Euclidean four-space,
\[
 \Omega\wedge\Omega=(|\Omega_+|^2-|\Omega_-|^2)\,\mathrm{vol},
 \qquad
 \Omega\wedge*\Omega=(|\Omega_+|^2+|\Omega_-|^2)\,\mathrm{vol}.
\]
The wedge pairing has signature \((3,3)\).
For orthonormal self-dual forms \(\sigma_1^+,\sigma_2^+\) and an
orthonormal anti-self-dual form \(\sigma_1^-\),
\[
 \Omega=b\sigma_1^++\tfrac12(a-c)\sigma_2^+
                    +\tfrac12(a+c)\sigma_1^-
\]
has wedge square \((b^2-ac)\mathrm{vol}\), but norm
\(b^2+(a^2+c^2)/2\).  Substituting three consecutive jets therefore
realizes \(\Lag_j\) algebraically while introducing an uncontrolled
higher jet into its positive norm.  This variable-coefficient
mixed-Hodge form also need not be closed.

There is a closed realization of the inputs.  Use coordinates
\((x,y,u,v)\), orientation \(dx\wedge dy\wedge du\wedge dv\),
\(z=x+iy\), and \(w=u+iv\).  For holomorphic \(f(z)\), put
\begin{equation}
 \alpha_f=\frac1{\sqrt2}\Re\{f(z)\dd z\wedge\dd w\}.
 \label{sf:alpha}
\end{equation}
With
\(A=dx\wedge du-dy\wedge dv\),
\(B=dx\wedge dv+dy\wedge du\), one has
\(*A=A\), \(*B=B\), \(A^2=B^2=2\,\mathrm{vol}\), \(A\wedge B=0\).
Thus
\begin{equation}
 d\alpha_f=0,\qquad *\alpha_f=\alpha_f,\qquad
 \alpha_f\wedge\alpha_g=\Re(f\bar g)\,\mathrm{vol}.
 \label{sf:closed-pairing}
\end{equation}
Closedness also follows immediately from
\(d(f(z)dz\wedge dw)=0\).  Heat time is external to this exterior
derivative.  Taking \(\alpha_r=\alpha_{h_r(t,\cdot)}\) yields
\[
 \alpha_{j+1}\wedge\alpha_{j+1}-\alpha_j\wedge\alpha_{j+2}
 =\ell_{h_j}\,\mathrm{vol},\qquad
 \ell_f=|f'|^2-\Re(f\overline{f''}).
\]
At \(y=0\) this coefficient is \(\Lag_j\).  This is evaluation of
a coefficient, not pullback of a four-form onto a line.
The coefficient-space matrix has entries
\(Q_{13}=Q_{31}=-1/2\), \(Q_{22}=1\), all others zero, and spectrum
\(\{1,1/2,-1/2\}\); closed self-dual inputs do not make this
polarized combination positive.  In particular, the quadratic and cubic
backward-heat collision controls above retain closed self-dual inputs
while their stationary ratios diverge as \(-1/(T-t)\).

For a holomorphic primitive \(P_f'=f\), the one-form
\(\beta_f=\Re(P_f\,dw)/\sqrt2\) has \(d\beta_f=\alpha_f\).
On any smooth bounded four-domain with these forms defined nearby,
\[
 \int_\Omega\alpha_f\wedge\alpha_g
                =\int_{\partial\Omega}\beta_f\wedge\alpha_g .
\]
Theta jets have global primitives.  Taking the first factor to be the
constant \(a=h_j(t,x_0)\) realizes exactly
\(\varphi_j(\theta)\) as the coefficient of
\(\alpha_a\wedge\alpha_{h_j}\) on the spectral circle.  Thus the
closed-form lift preserves the signed current already isolated above.

\begin{theorem}[Exact disk Dirichlet completion]
\label{sf:disk}
Let \(f\) be holomorphic near the closure of
\(D_R(x_0)=\{|z-x_0|<R\}\), \(x_0\in\R\), and set
\[
 g_R(z)=\left(1-\frac{(z-x_0)^2}{R^2}\right)f(z).
\]
Then
\begin{equation}
 \int_{D_R(x_0)}\ell_f\dd A+
       \frac1{2R}\int_{\partial D_R(x_0)}|f|^2\dd s
       =\frac12\int_{D_R(x_0)}|g_R'|^2\dd A\ge0.
 \label{sf:disk-completion}
\end{equation}
The boundary coefficient is sharp even among real-reflection
polynomials stationary at \(x_0\).
\end{theorem}
\begin{proof}
Expand \(f(x_0+R\zeta)=\sum_{n\ge0}b_n\zeta^n\), with negative-index
coefficients zero.  Angular orthogonality gives
\[
 \pi^{-1}\int_{D_R}\ell_f=
 \sum_{n\ge0}n|b_n|^2
       -\Re\sum_{n\ge0}(n+2)b_n\overline{b_{n+2}},
\]
while the normalized boundary term is \(\sum|b_n|^2\) and
\(\pi^{-1}\int|g_R'|^2=\sum_{n\ge1}n|b_n-b_{n-2}|^2\).
Expanding this square proves the identity.  For
\[
 f_N(z)=\sum_{m=0}^N\left(1-\frac m{N+1}\right)
                         ((z-x_0)/R)^{2m},
\]
one has \(f_N(x_0)=1\), \(f_N'(x_0)=0\), and, writing
\(B_R=(2R)^{-1}\int_{\partial D_R}|f_N|^2ds\),
\[
 B_R=\pi\frac{(N+2)(2N+3)}{6(N+1)},\qquad
 \frac{\int_{D_R}\ell_{f_N}dA}{B_R}
                         =-\frac{2N}{2N+3}\longrightarrow-1 .
\]
This proves sharpness in the stated general class.
\end{proof}

For \(U(x,y)=|f(x+iy)|^2\), the Cauchy--Riemann equations give
\begin{equation}
 \ell_f=\tfrac12U_{yy}=2|f'|^2-\tfrac12U_{xx},\qquad
 \int_{D_R}\ell_f\dd A=\tfrac12\int_{\partial D_R}U_y n_y\dd s .
 \label{sf:normal-density}
\end{equation}
Real-reflection data have \(U_y(x,0)=0\), which gives no sign for
\(U_{yy}(x,0)\).  On an annulus the inner boundary contributes its
own signed flux.  The sharpness sequence also shows that the disk
loss cannot be paid by the central first-jet energy for arbitrary
holomorphic fields.  Neither a boundary norm nor a positive
Dirichlet remainder establishes central pointwise positivity.

\subsection{Positive centered averages from the complete source}

\begin{theorem}[Centered source positivity]
\label{sf:centered}
Let \(\chi\) be a nonzero real, even, nonnegative integrable weight,
and put
\[
 \widehat\chi(\xi)=\int_{\R}\chi(x)\cos(\xi x)\dd x,\qquad
 W_j(t,y)=\int_{\R}\chi(x)\ell_{h_j}(t,x+iy)\dd x .
\]
For every real \(t,y\), \(W_0(t,y)>0\) if
\(\widehat\chi\ge0\).  If in addition \(\widehat\chi\) is
nonincreasing on \([0,\infty)\), then \(W_1(t,y)>0\).
Both conclusions hold for every centered Gaussian
\(\chi(x)=e^{-a^2x^2}\), \(a>0\).
\end{theorem}
\begin{proof}
With \(\rho_t(u)=w_t(|u|)/2\), symmetrizing the Hermitian products
in the full-line Fourier representation gives
\[
 \ell_{h_j}(t,x+iy)=\frac12\iint_{\R^2}
 \rho_t(u)\rho_t(v)(uv)^j(u+v)^2
                 \cos((u-v)x)e^{-(u+v)y}\dd u\dd v .
\]
Integration against \(\chi\) and folding the four quadrants yield
\begin{equation}
 \begin{split}
 W_j(t,y)=\frac14\int_0^\infty\!\int_0^\infty
 &w_t(u)w_t(v)(uv)^j\\
 {}\times\bigl\{&
 (u+v)^2\widehat\chi(u-v)\cosh((u+v)y)\\
 &+(-1)^j(u-v)^2\widehat\chi(u+v)\cosh((u-v)y)
                      \bigr\}\dd u\dd v .
 \end{split}
 \label{sf:centered-fold}
\end{equation}
For \(j=0\) both terms are nonnegative.  For \(j=1\), the first
dominates the second since \(u+v\ge|u-v|\), the Fourier weight
decreases with absolute frequency, and \(\cosh\) increases with
absolute argument.  Strict positivity holds on a region near
\(u=v>0\), using continuity and \(\widehat\chi(0)>0\).
The complete double-exponential source pays every polynomial and
\(e^{|y|(|u|+|v|)}\), locally uniformly in \(t,y\); integrability
of \(\chi\) therefore justifies every interchange.  For a Gaussian,
\(\widehat\chi(\xi)=\sqrt\pi\,a^{-1}e^{-\xi^2/(4a^2)}\).
\end{proof}

At \(y=0\), these are strictly positive centered averages of the
physical branches, without the loss in \eqref{sf:sech-bound}.
The origin is essential: replacing \(\chi(x)\) by \(\chi(x-x_0)\)
inserts \(\cos((u-v)x_0)\) in the first folded term and
\(\cos((u+v)x_0)\) in the second.  They can have independent signs.

An exact positive atomic-source control quantifies the localization
gap.  Take \(a,k>0\) and
\[
 H(t,z)=a+b(t)\cos(kz),\qquad b(t)=ae^{k^2(t-T)}.
\]
At \(x_*=\pi/k\), \(\delta=T-t>0\), and
\(\varepsilon=a-b>0\), one has
\[
 h_1=0,\quad E=\varepsilon^2,\quad
 \Lag_0=-bk^2\varepsilon,\quad
 \frac{\Lag_0}{E}=-\frac{k^2}{e^{k^2\delta}-1}
                           =-\frac1\delta+O(1).
\]
On the centered disk use \(\zeta=z-x_*\) and write
\(f(\zeta)=a-b\cos(k\zeta)\).  Direct
calculation gives
\(\ell_f=b^2k^2\cosh(2k\Im\zeta)-abk^2\Re\cos(k\zeta)\).
The harmonic mean-value formula and termwise integration yield
\begin{equation}
 \frac1{\pi R^2}\int_{D_R(0)}\ell_f\dd A
 =k^2b^2\{M(kR)-e^{k^2\delta}\},\qquad
 M(s)=\frac{I_1(2s)}s
       =\sum_{n\ge0}\frac{s^{2n}}{n!(n+1)!}.
 \label{sf:atomic-disk}
\end{equation}
Every fixed positive-radius disk has positive average sufficiently
near the event, even though its central value is negative.
Since \(M\) strictly increases for positive arguments and
\(M(s)=1+s^2/2+O(s^4)\), the unique sign-changing radius obeys
\(R_{\rm crit}^2=2\delta+O(\delta^2)\); the average is negative only
for smaller radii.  Detecting this adverse center by its disk mean
therefore forces localization at the collision scale.  This is a
control, not the theta source.  Closed self-dual forms and positive
source averages retain useful signed information, but do not
supply the remaining translated, stationary, energy-relative bound.


% Prepared 15 October 2026 for Edward Baker with substantial LLM assistance.
% Recorded model: gpt-6-astra; reasoning effort: ultra.
% Integrates Notes 52--53 and their internal cross-review in Review 31.
% Suggested insertion: after signed forms, before bounded-region certificates.

\section{A positive theta primitive and source-shape obstructions}
\label{pp:section}

This section gives a positive complete-theta auxiliary channel and its
exact signed reconstruction, then separates that structure from two
insufficient replacements for the missing estimate: finite rational
jet closure and general source-shape assumptions. All results are
analytic; they add no finite certificate or new pointwise theta sign
domain.

\subsection{A positive theta primitive and its signed reconstruction}
\label{pp:primitive-section}

The complete theta source admits a positive auxiliary transform whose
relation to the physical field retains an essential lower-boundary term.
This is useful source structure, but its positivity does not imply the
sign of the differentiated physical observation.

\begin{proposition}[Positive complete theta primitive]
\label{pp:positive-primitive}
For \(u\ge0\), put
\begin{equation}
 \psi(u)=e^u\sum_{n\ge1}e^{-\pi n^2e^{4u}},\qquad
 J(t,z)=\int_0^\infty e^{tu^2}\psi(u)\cos(zu)\dd u.
 \label{pp:primitive-definition}
\end{equation}
The integral defines a jointly entire function of \((t,z)\).
For every real \(t,x\),
\begin{equation}
 8H(t,x)=\frac12+(2t-x^2-1)J(t,x)
                    +4txJ_x(t,x)-4t^2J_{xx}(t,x).
 \label{pp:primitive-reconstruction}
\end{equation}
Moreover,
\begin{equation}
 J(t,x)>0\qquad(0\le t\le1/2,\ x\in\R).
 \label{pp:primitive-positive}
\end{equation}
Uniformly for \(0\le t\le T<\infty\), as real \(|x|\to\infty\),
\begin{equation}
 J(t,x)=\frac1{2x^2}-\frac{1+6t}{2x^4}+O_T(x^{-6}).
 \label{pp:primitive-asymptotic}
\end{equation}
\end{proposition}
\begin{proof}
Let \(\vartheta(r)=\sum_{n\in\mathbb Z}e^{-\pi n^2r}\).
Poisson summation gives
\[
 P(u):=e^u\vartheta(e^{4u})=P(-u),\qquad P=e^u+2\psi.
\]
Differentiating the complete series gives
\begin{equation}
 \Phi=\frac{P''-P}{16}=\frac{\psi''-\psi}{8},\qquad
 \psi^{(2k+1)}(0)=-\frac12\quad(k\ge0).
 \label{pp:primitive-source}
\end{equation}
The last identities follow by differentiating the even function \(P\).
The complete series for \(\psi\), and every fixed derivative, decay
double exponentially as \(u\to+\infty\). On compact sets of complex
\((t,z)\), this decay dominates every polynomial derivative factor,
the heat multiplier, and \(e^{|\operatorname{Im}z|u}\). This proves
joint entirety and justifies the following differentiations and
integrations by parts.

With \(F(u)=e^{tu^2}\cos(xu)\), two integrations by parts give
\[
 8H=-\psi'(0)F(0)+\psi(0)F'(0)
                      +\int_0^\infty\psi(u)(F''(u)-F(u))\dd u.
\]
Here \(F(0)=1\), \(F'(0)=0\), and all upper boundary terms vanish.
Expanding \(F''\) and using spectral derivatives of \(J\) gives
\eqref{pp:primitive-reconstruction}, including the constant \(1/2\).

To prove positivity, fix \(0\le t\le1/2\) and write
\[
 f(u)=e^{tu^2}\psi(u)=\sum_{n\ge1}f_n(u),\quad
 f_n(u)=e^{tu^2+u-q_n(u)},\quad q_n(u)=\pi n^2e^{4u}.
\]
For \(u\ge0\), \(0<2tu+1\le u+1\le q_n(u)\). Consequently
\[
 \frac{f_n''(u)}{f_n(u)}
  =(2tu+1-4q_n)^2+2t-16q_n
  \ge9q_n^2-16q_n>0.
\]
Thus \(f''>0\), while \(f'(0)=-1/2\) and
\(f'(u)\to0\) at infinity. For \(x\ne0\), integration by parts and
\(\int_0^\infty f''=1/2\) yield
\begin{equation}
 J(t,x)=\frac1{x^2}\int_0^\infty
                  f''(u)(1-\cos(xu))\dd u>0.
 \label{pp:primitive-convex-representation}
\end{equation}
Strict positivity at \(x=0\) follows from \(f>0\).

Finally, \eqref{pp:primitive-source} gives
\(f'(0)=-1/2\) and \(f'''(0)=-(1+6t)/2\).
Repeated integration by parts therefore gives
\[
 J=-\frac{f'(0)}{x^2}+\frac{f'''(0)}{x^4}
       -\frac{f^{(5)}(0)}{x^6}
       -\frac1{x^6}\int_0^\infty f^{(6)}(u)\cos(xu)\dd u.
\]
The endpoint value and the last derivative integral are uniformly
bounded for \(0\le t\le T\), proving
\eqref{pp:primitive-asymptotic}.
\end{proof}

This is a direct convex-kernel argument of the type in
\cite[Theorem~3.4]{csordas}. The even extension of
\(e^{tu^2}\psi(u)\) has a cusp at zero, with one-sided derivative
\(-1/2\). It is different from the smooth even physical source
\(e^{tu^2}\Phi(u)\). One cannot instead integrate the modular-even
\(P\) at positive time: \(P(u)\sim e^u\), so its proposed heat
integral diverges. The constant in
\eqref{pp:primitive-reconstruction} retains precisely the boundary
contribution that the decaying half-line primitive requires.

\begin{proposition}[Operator factorization and the scale of the primitive]
\label{pp:primitive-limitations}
For \(D_t=2t\partial_x-x\),
\begin{equation}
 8H=\frac12-(D_t^2+1)J.
 \label{pp:primitive-factorization}
\end{equation}
For \(t>0\), the closure of \(D_t\) initially defined on
\(C_c^\infty(\R)\) is skew-adjoint in
\(L^2(\R,e^{-x^2/(2t)}\dd x)\). The self-adjoint operator
\(-D_t^2-1\) has spectrum \([-1,\infty)\), and its quadratic
form takes both signs. For each fixed real \(t\), no bound
\(J(t,x)^2\le C(1+|x|)^M E(t,x)\), with finite
\(C,M\ge0\), holds for all sufficiently large \(|x|\).
\end{proposition}
\begin{proof}
Expanding \(D_t^2\) proves
\eqref{pp:primitive-factorization}. The unitary map from the weighted
space to ordinary \(L^2\) is
\[
 (U_tf)(x)=e^{-x^2/(4t)}f(x),\qquad
 U_tD_tU_t^{-1}=2t\partial_x.
\]
The derivative on its \(H^1\) domain is skew-adjoint. Thus the
conjugated second-order operator is
\(-4t^2\partial_x^2-1\), with Fourier multiplier
\(4t^2\xi^2-1\). This proves the spectral and form assertions.

For any fixed \(0<a<\pi/8\), complete theta-strip holomorphy
permits shifting the full-line Fourier integral to
\(u+ia\) when \(x>0\), and to \(u-ia\) when \(x<0\).
On each shifted line the positive-end series decays like a polynomial
in \(e^{4|u|}\) times
\(\exp\{-\pi\cos(4a)e^{4|u|}\}\); complete evenness supplies
the same bound at the other end. Every fixed derivative is integrable
after multiplication by the heat factor. Hence
\[
 |h_k(t,x)|\le C_{t,a,k}e^{-a|x|},\qquad k=0,1.
\]
The same endpoint integration argument as above gives
\(J(t,x)\sim1/(2x^2)\) at any fixed real \(t\).
Exponential decay of \(E\) is incompatible with the proposed
polynomial-height payment of \(J^2\).
\end{proof}

The last statement concerns a uniform or polynomial-height budget;
it does not exclude coefficients merely locally bounded at every
finite height. There is a separate local limitation at a hypothetical
real collision in \(0\le t\le1/2\): there \(E=0\) but
\(J>0\), so a bound \(J^2\le C E\) cannot hold through that
event with locally bounded \(C\). At time zero the exact identity is
\begin{equation}
 J(0,x)=\frac{1}{2(1+x^2)}-\frac{8H(0,x)}{1+x^2}.
 \label{pp:primitive-time-zero}
\end{equation}
Thus the physical field is a small difference involving a positive
auxiliary channel. At positive time the full differential operator
in \eqref{pp:primitive-reconstruction} performs that cancellation.
A useful estimate must retain this complete operator and its signed
products with the physical jets. Positivity or separate derivative
norms of \(J\) do not give the required stationary lower bound.

\subsection{The complete theta field has no rational finite-jet closure}
\label{pp:closure-section}

\begin{proposition}[Complete-source rational closure obstruction]
\label{pp:no-rational-closure}
Fix any real \(t\). There are no rational functions
\(r_0(x),\ldots,r_m(x)\), not all zero, for which
\begin{equation}
 \sum_{k=0}^m r_k(x)h_k(t,x)=0
 \label{pp:forbidden-closure}
\end{equation}
on a nonempty open real interval free of their poles.
\end{proposition}
\begin{proof}
Set \(w_t(u)=e^{tu^2}\Phi(u)\). For \(u\to+\infty\),
the first theta index gives
\[
 \Phi(u)=2\pi^2e^{9u-\pi e^{4u}}\{1+O(e^{-4u})\}.
\]
For each fixed integer \(k\ge0\), direct differentiation of the
complete series gives
\begin{equation}
 \frac{w_t^{(k)}(u)}{w_t(u)}
       =(-4\pi)^k e^{4ku}\{1+o(1)\}.
 \label{pp:source-derivative-asymptotic}
\end{equation}
Indeed, the leading term differentiates
\(e^{-\pi e^{4u}}\) at every step. Differentiation of the
Gaussian-polynomial prefactors, or of factors generated at earlier
steps, lowers the power of \(e^{4u}\). After any fixed number of
derivatives, each index \(n\ge2\) still carries the additional
factor \(e^{-\pi(n^2-1)e^{4u}}\); the polynomial factors in
\(n\) and \(e^{4u}\) are dominated by this gap, including in
the sum over \(n\).

Suppose \(\sum_{k=0}^M p_k(u)w_t^{(k)}(u)=0\), where the
\(p_k\) are polynomials and \(p_M\ne0\). For all sufficiently
large \(u\), divide by \(p_M(u)w_t(u)e^{4Mu}\).
By \eqref{pp:source-derivative-asymptotic}, the highest term tends
to \((-4\pi)^M\), whereas every lower term tends to zero.
Ratios of polynomials cannot compensate for
\(e^{-4(M-k)u}\). This contradiction excludes every nonzero
polynomial-coefficient differential operator annihilating \(w_t\).

Now suppose \eqref{pp:forbidden-closure} holds. Clearing its
denominators gives a nonzero polynomial-coefficient differential
operator annihilating \(H(t,\cdot)\). Entirety extends this
identity from the given interval to the whole spectral line.
The complete even source is a Schwartz function, and
\[
 H(t,x)=\frac12\int_{\R}w_t(u)e^{ixu}\dd u.
\]
For this Fourier convention, conjugation sends multiplication by
\(x\) to \(i\partial_u\), and \(\partial_x\) to
multiplication by \(iu\), respecting the order of composed factors.
This is an invertible correspondence on polynomial differential
operators. Fourier injectivity therefore produces a nonzero
polynomial differential operator annihilating \(w_t\), contrary
to the preceding paragraph.
\end{proof}

This is distinct from the compact-filter obstruction of
Proposition~\ref{prop:no-rational-filter-closure}: it applies to the
untruncated complete field itself. In particular, exact identities
\(-h_2=a_0h_0+b_0h_1\) or
\(-h_3=a_1h_0+b_1h_1+c_1h_2\) cannot hold on an open spectral
interval with all coefficients rational in \(x\). The proposition
does not exclude identities restricted to a stationary set,
transcendental coefficients, signed residuals, inequalities, or
infinite and nonlocal relations. The auxiliary transform \(J\)
is a new channel and does not supply the excluded finite closure
among the original jets.

\subsection{A smooth concave source with a positive-time collision}
\label{pp:collision-section}

The next control rules out a deduction of the needed energy-relative
bound from positivity, strong log-concavity, square-variable
log-concavity, and fast source decay alone. It does not reproduce
the theta arithmetic and is not a counterexample for \(\Phi\).

\begin{theorem}[A double-exponential source-shape control]
\label{pp:smooth-collision}
There exists a positive, even, smooth source \(\phi\), with
double-exponential decay of every derivative, such that
\[
 H_\phi(t,z)=\int_0^\infty e^{tu^2}\phi(u)\cos(zu)\dd u
\]
is jointly entire and obeys \(\partial_tH_\phi=-\partial_z^2H_\phi\).
For every real \(t\),
\[
 v\longmapsto\log\{e^{tv}\phi(\sqrt v)\}
\]
has strictly negative second derivative for \(v\ge0\).
For any fixed \(\tau\in(1/2,4/5)\), the deformed source
\(e^{tu^2}\phi(u)\) is strictly decreasing on \(u>0\) and
uniformly strongly log-concave in \(u\) for \(t\le\tau\).
Nevertheless, there is a nondegenerate double collision
\((T,x_*)\), with \(0<T<1/2\) and \(x_*>0\), and a genuine
stationary predecessor curve \(x=x(t)\) on which
\begin{equation}
 \frac{\Lag_{0,\phi}(t,x(t))}{E_\phi(t,x(t))}
           =-\frac1{T-t}+O(1)\qquad(t\uparrow T).
 \label{pp:smooth-adverse-pole}
\end{equation}
Here \(h_{k,\phi}=\partial_x^kH_\phi\),
\(E_\phi=h_{0,\phi}^2+h_{1,\phi}^2\),
\(\Lag_{0,\phi}=h_{1,\phi}^2-h_{0,\phi}h_{2,\phi}\), and
\(h_{1,\phi}(t,x(t))=0\).
\end{theorem}
\begin{proof}
Start with an explicitly computable Gaussian-polynomial source:
\[
 b=\frac15,\quad c=\frac1{50},\quad b^2=2c,\qquad
 \phi_0(u)=e^{-u^2}(1+bu^2+cu^4).
\]
For \(a=1-t>0\), differentiation of the Gaussian cosine integral
gives
\begin{align}
 H_0(t,x)&=\frac{\sqrt\pi}{2\sqrt a}e^{-x^2/(4a)}
                           \frac{P_a(x^2)}{16a^4},
 \label{pp:gaussian-control}\\
 P_a(y)&=cy^2-(4ba^2+12ca)y+16a^4+8ba^3+12ca^2.
 \label{pp:gaussian-polynomial}
\end{align}
For verification, before collecting denominators its multiplier is
\[
 1+b\left(\frac1{2a}-\frac{x^2}{4a^2}\right)
   +c\left(\frac3{4a^2}-\frac{3x^2}{4a^3}
                         +\frac{x^4}{16a^4}\right).
\]
The discriminant of \(P_a\), as a polynomial in \(y\), is
\[
 \Delta(a)=16a^2\{(b^2-4c)a^2+4bca+6c^2\}
          =32ca^2(-a^2+2ba+3c).
\]
It vanishes at
\begin{equation}
 a_c=\frac{2+\sqrt{10}}{10},\qquad
 T_0=1-a_c=\frac{8-\sqrt{10}}{10}\in(0,1/2).
 \label{pp:gaussian-time}
\end{equation}
At this value,
\[
 P_{a_c}(y)=c(y-y_c)^2,\qquad
 y_c=\frac{2ba_c^2}{c}+6a_c=4+\frac{7\sqrt{10}}5>0.
\]
With \(x_c=\sqrt{y_c}\), one has
\(H_0(T_0,x_c)=\partial_xH_0(T_0,x_c)=0\), and
\begin{equation}
 d_0:=\partial_x^2H_0(T_0,x_c)
   =\frac{\sqrt\pi}{2\sqrt{a_c}}e^{-y_c/(4a_c)}
                        \frac{c y_c}{2a_c^4}>0.
 \label{pp:gaussian-second-derivative}
\end{equation}

For \(\delta>0\), replace this source by
\begin{equation}
 \phi_\delta(u)=e^{-u^2-\delta\cosh(4u)}(1+bu^2+cu^4).
 \label{pp:damped-source}
\end{equation}
For any fixed positive \(\delta\), compact parameter sets and
any fixed numbers of heat and spectral derivatives are dominated by
\[
 C(1+u^M)\exp\{A u^2+B u-\delta\cosh(4u)\},
\]
with finite \(A,B,C,M\). This proves joint entirety of its transform
and the backward heat equation. Differentiating the source in \(u\)
adds only polynomial factors and finite products of hyperbolic
functions. These are absorbed by the same double-exponential damping,
which proves the derivative-decay assertion.

On a compact neighborhood of \((T_0,x_c)\) with \(t<1\),
\(e^{-\delta\cosh(4u)}\le1\) permits a fixed Gaussian majorant
for every required finite derivative. Dominated convergence therefore
gives local uniform convergence of these derivatives as
\(\delta\downarrow0\). Apply the continuous-parameter implicit
function theorem to
\[
 \mathscr F_\delta(t,x)
       =(H_\delta(t,x),\partial_xH_\delta(t,x)).
\]
Its derivative in \((t,x)\) at \((0,T_0,x_c)\) is
\[
 D_{t,x}\mathscr F_0=
 \begin{pmatrix}
 -d_0&0\\ -\partial_x^3H_0(T_0,x_c)&d_0
 \end{pmatrix},\qquad \det D_{t,x}\mathscr F_0=-d_0^2<0.
\]
Consequently, for all sufficiently small positive \(\delta\),
there are collision points \((T_\delta,x_\delta)\to(T_0,x_c)\)
with
\begin{equation}
 H_\delta(T_\delta,x_\delta)
  =\partial_xH_\delta(T_\delta,x_\delta)=0,\qquad
 d_\delta:=\partial_x^2H_\delta(T_\delta,x_\delta)>0.
 \label{pp:damped-collision}
\end{equation}
After reducing the parameter interval, \(0<T_\delta<1/2\) and
\(x_\delta>0\). Only continuity in the one-sided parameter and
continuous first derivatives in \((t,x)\) are needed. Alternatively,
the one-sided parameter derivative is Gaussian dominated even with
the extra factor \(\cosh(4u)\). Extending \(\mathscr F\) to
negative parameters by \(2\mathscr F_0-\mathscr F_{-\delta}\)
gives a \(C^1\) map to which the ordinary implicit function theorem
applies; it does not integrate a negatively damped source.

It remains to verify all shape properties for the same fixed positive
\(\delta\). Write
\[
 Q(v)=1+bv+cv^2,\quad C(v)=\cosh(4\sqrt v),\quad
 \eta_{\delta,t}(v)=-(1-t)v-\delta C(v)+\log Q(v).
\]
The entire power series
\(C(v)=\sum_{k\ge0}16^kv^k/(2k)!\) gives
\(C'(v)\ge8\) and \(C''(v)>0\) for \(v\ge0\).
Also
\begin{equation}
 (\log Q)''(v)
  =\frac{2c-b^2-2bcv-2c^2v^2}{Q(v)^2}
  =-\frac{2bcv+2c^2v^2}{Q(v)^2}.
 \label{pp:shape-log-polynomial}
\end{equation}
Thus \(\eta_{\delta,t}''(v)<0\) for all real \(t\) and
\(v\ge0\). Since \((\log Q)'\le b\),
\(\eta_{\delta,t}'\le-(1-t)+b-8\delta\). For
\(t\le\tau<4/5=1-b\), this implies
\begin{align*}
 \partial_u\log(e^{tu^2}\phi_\delta(u))
       &=2u\eta_{\delta,t}'(u^2)<0 &&(u>0),\\
 \partial_u^2\log(e^{tu^2}\phi_\delta(u))
       &=2\eta_{\delta,t}'(u^2)+4u^2\eta_{\delta,t}''(u^2)
         \le-2(1-\tau-b)<0.
\end{align*}
Positivity and evenness follow directly from
\eqref{pp:damped-source}. These bounds hold throughout a time
neighborhood of the collision and its predecessors.

Finally fix any sufficiently small positive \(\delta\), and set
\(\phi=\phi_\delta\), \(T=T_\delta\), \(x_*=x_\delta\).
By \eqref{pp:damped-collision}, the implicit function theorem for
\(\partial_xH_\delta=0\) gives a smooth stationary curve through
the collision. Along this curve,
\[
 \frac{\dd}{\dd t}H_\delta(t,x(t))
  =\partial_tH_\delta+\partial_xH_\delta x'(t)
  =-\partial_x^2H_\delta(t,x(t)).
\]
Writing \(\varepsilon=T-t\downarrow0\), Taylor expansion gives
\[
 h_{0,\phi}=d_\delta\varepsilon+O(\varepsilon^2),\quad
 h_{1,\phi}=0,\quad h_{2,\phi}=d_\delta+O(\varepsilon),
\]
and hence
\[
 E_\phi=d_\delta^2\varepsilon^2+O(\varepsilon^3),\qquad
 \Lag_{0,\phi}=-d_\delta^2\varepsilon+O(\varepsilon^2).
\]
This proves \eqref{pp:smooth-adverse-pole}. All predecessors use
the same fixed, strictly positive damping and the complete transform.
\end{proof}

No numerically certified damping size is asserted: the theorem is an
existence result for every sufficiently small positive \(\delta\).
The explicitly solvable source at \(\delta=0\) establishes the
nondegeneracy used to obtain a fixed double-exponentially decaying
source. The adverse limit then varies time, not the source.

The source properties in this theorem imply no locally bounded
coefficient through the event satisfying
\(\Lag_{0,\phi}\ge-\kappa E_\phi\) on its full stationary
predecessors. Failure on this branch already excludes a deduction
of both branch bounds from those properties alone. The complete
correlation identity \eqref{eq:corr-L0}, applied to this source,
has a positive source kernel but a negative cosine transform on
these predecessors. Positive moments, admissibility of associated
kernels, and Fourier positivity are different assertions; compare
\cite[Theorems~3.5 and~4.2(b), Remark~2.1, and Example~3.12]{csordas}.
The exact closed-form lifts apply to the entire jets of this control
as well. Additional theta-specific structure must therefore enter
at the step that bounds the signed physical observation by its
first-jet energy.


\section{The bounded region: complete certificates through height 112}
\label{sec:fc-certificate}

The bounded calculation supplies a genuine norm gap, so it excludes
collisions directly.  Its role is to join the central analytic estimates
to the large-height theorem, while leaving the intervening and
small-$t\log|x|$ regions explicitly visible.

\begin{theorem}[Certified closed bounded region]
\label{fc:height104}
For the complete theta source,
\begin{equation}
 |F(t,x)|=\sqrt{h_0(t,x)^2+h_1(t,x)^2}>10^{-16}
 \quad(0\le t\le1/2,\ |x|\le104).
 \label{fc:height104-gap}
\end{equation}
On this rectangle, $h_1=0$ implies $\mathcal L_0>0$, and $h_2=0$
implies $\mathcal L_1>0$.  The norm bound combines an analytic central
estimate with the validated interval records described below.
\end{theorem}

\subsection{A central estimate with no theta truncation}

\begin{proposition}[Central norm gap]
\label{fc:central}
For $0\le t\le1$ and $|x|\le8$,
$|F(t,x)|>1/120000$.  Both stationary signs in
Theorem~\ref{fc:height104} hold on this larger time interval.
\end{proposition}
\begin{proof}
Put $w_t(u)=e^{tu^2}\Phi(u)$ and
$M_k(t)=\int_0^\infty u^kw_t(u)\,du$.
The complete source obeys
$0<\Phi(u)<4\pi^2e^{9u-\pi e^{4u}}$ for $u\ge0$:
after dropping each summand's negative term, the ratio of successive
$n^4e^{-\pi e^{4u}n^2}$ terms is at most $16e^{-3\pi}<1/2$.
The following elementary lower cores and upper tails hold uniformly
for $0\le t\le1$:
\begin{equation}
 \begin{array}{c|c@{\qquad}c|c}
 \text{core }[0,d]&\Phi(u)>c_d&\text{tail }u\ge b&w_t(u)<A_be^{-q_b(u-b)}\\ \hline
 d=1/10&6/25&b=7/25&(1/20)e^{-28(u-b)}\\
 d=3/20&3/25&b=3/10&(1/49)e^{-32(u-b)}\\
 d=1/5&1/20&&
 \end{array}
 \label{fc:source-envelopes}
\end{equation}
For completeness, the lower bounds follow from
$\phi_1(u)=e^ug(\pi e^{4u})$ with
$g(q)=q(2q-3)e^{-q}$ and $g'(q)=-(2q-1)(q-3)e^{-q}$.
On the three cores the respective arguments are less than
$33/7,6,36/5$.  The bounds
$e<11/4$, $e^{2/7}>9/7$, and $e^{1/5}<5/4$ give the three stated
rational lower bounds for $g$.
For the tails, $\psi_t(u)=9u+tu^2-\pi e^{4u}$ is strictly concave.
At $b=7/25$ one has $\psi_t(b)<-69/10$, $\psi_t'(b)<-28$;
at $b=3/10$ one has $\psi_t(b)<-38/5$, $\psi_t'(b)<-32$.
Together with $4\pi^2<40$, $e^{69/10}>800$, and $e^{38/5}>1960$,
these prove the full upper tails in \eqref{fc:source-envelopes}.
All comparisons use rational Taylor bounds for the indicated
exponentials; in particular they do not truncate the theta source.

Integrating $u^2(1-25u^2/2)$ over the first positive core and paying
its entire negative tail with $b=7/25$ gives
\begin{align}
 M_2-\frac{25}{2}M_4
 &>\frac{37}{500000}
   -\frac58\frac{3497259}{33614000000}
   =\frac{2413193}{268912000000}>\frac1{120000}.
 \label{fc:moment24}
\end{align}
Indeed the polynomial is nonnegative up to $b$ and is bounded below
on the tail by $-(25/2)u^2(u^2-b^2)$.
Likewise, with $f(u)=u^4(1-32u^2/3)$ and $b=3/10$, the three
successive cores and the entire negative tail give
\begin{align}
 M_4-\frac{32}{3}M_6
 &>\frac{111799}{33600000000}
   -\frac{32}{147}\frac{85837761}{6710886400000}
   =\frac{2092283587}{3853516800000000}>\frac1{2000000}.
 \label{fc:moment46}
\end{align}
The positive term is
$(6/25)\int_0^{1/10}f+(3/25)\int_{1/10}^{3/20}f
 +(1/20)\int_{3/20}^{1/5}f$;
the tail integral is
$\int_b^\infty u^4(u^2-b^2)e^{-32(u-b)}\,du$.
Thus both displayed differences are exact rational integrations.

The global inequalities $\sin y\ge y-y^3/6$ for $y\ge0$ and
$\cos y\ge1-y^2/2$ now imply, for $0<x\le8$,
\[
 -h_1\ge x(M_2-x^2M_4/6)>x/120000,
 \qquad h_3\ge x(M_4-x^2M_6/6)>x/2000000.
\]
For $|x|\le1$, integrate the first core using
$\cos(xu)\ge199/200$, discard the nonnegative middle interval
$[1/10,7/25]$, and pay the entire remaining tail.  This gives
\[
 H(t,x)>\frac{597}{25000}-\frac1{560}
 =\frac{7733}{350000}>\frac1{50}.
\]
The value bound covers $|x|\le1$, and the slope bound covers
$1\le|x|\le8$, proving the norm gap by evenness.
The first stationary set is exactly $x=0$, where
$\mathcal L_0=H(t,0)M_2>0$.
For $0<|x|\le8$,
$\mathcal L_1=h_2^2-h_1h_3>x^2/240000000000$;
at $x=0$ one has $h_2=-M_2<0$, so the second stationary set is absent.
\end{proof}

\subsection{What the interval certificates enclose}

The records use Arb/Acb interval arithmetic to enclose the complete
physical jets on closed space--time cells.  Each enclosure includes
the returned finite-quadrature radius, every omitted theta index,
the complete integration tail, rounding, and spatial and temporal
Taylor remainders.  The older compensated calculation also pays the
complete Gaussian transforms and their truncation remainders.
Finite compensated channels are transported using the raw heat
equation followed by their explicit time-dependent compensation;
they are not assigned a homogeneous heat equation.

For a heat-transported sequence $Y_r$, a cell with half-widths
$d_t,d_x$ uses
\[
 Y_r(t_c+\tau,x_c+\eta)=
 \sum_{\ell=0}^{q-1}\sum_{m=0}^{p-1}
 \frac{(-\tau)^\ell\eta^m}{\ell!m!}Y_{r+2\ell+m}(t_c,x_c)
 +\mathcal R_r.
\]
If $M_k$ bounds the complete $k$th derivative on the cell, Taylor's
theorem first in time and then in space gives
\begin{equation}
 |\mathcal R_r|\le
 \sum_{\ell=0}^{q-1}\frac{d_t^\ell}{\ell!}
 \frac{d_x^p}{p!}M_{r+2\ell+p}
 +\frac{d_t^q}{q!}M_{r+2q}.
 \label{fc:taylor-payment}
\end{equation}
The higher bands sum the retained theta coefficients before applying
the common Taylor monomials.  This preserves cancellation between
different indices while adding every positive remainder payment.

Stationary decisions concern genuine derivatives.  If
$h_k=p_k+r_k$ with $|r_k|\le\varepsilon_k$, then at $h_{j+1}=0$
one has
\begin{equation}
 \mathcal L_j\ge p_{j+1}^2-p_jp_{j+2}
 -\bigl(\varepsilon_{j+1}^2+|p_j|\varepsilon_{j+2}
       +|p_{j+2}|\varepsilon_j+\varepsilon_j\varepsilon_{j+2}\bigr).
 \label{fc:stationary-payment}
\end{equation}
Here $p_{j+1}^2\le\varepsilon_{j+1}^2$ follows from genuine
stationarity; the square payment has a negative sign.
The products retain all cross terms between retained channels.
A cell is discharged either by exclusion of the stationary derivative
or by a strictly positive paid signed lower bound.
Monotone endpoint exclusions, where used, employ the next physical
derivative and uniform endpoint signs.  No zero of an ordinary
derivative of an exponentially rescaled function is substituted.

Independently, if $I_0,I_1$ enclose the genuine $h_0,h_1$ throughout
a cell, then
\begin{equation}
 E\ge \operatorname{dist}(0,I_0)^2+\operatorname{dist}(0,I_1)^2.
 \label{fc:interval-norm}
\end{equation}
Exact rational comparisons and downward square-root rounding turn
these energy floors into the reported norm floors.  No division by
an unproved norm is used.

\subsection{Coverage, the centered-residual pilot, and reproducibility}

Every numerical band covers the full closed time interval $[0,1/2]$
using eight slabs $[r/16,(r+1)/16]$.  The records give the following
conservative norm bounds; the central row is analytic.
\begin{center}
\begin{tabular}{c r c c}
Positive-height band&Cells&Precision&Proved lower bound for $|F|$\\ \hline
$[0,8]$&---&exact estimates&$1/120000$\\
$[7,41]$&272&256 bits&$9/5000000$\\
$[40,100]$&510&256 bits&$10^{-16}$\\
$[100,104]$&32&192 bits&$4\cdot10^{-16}$
\end{tabular}
\end{center}
The 510 cells comprise 160, 170, and 180 cells in $[40,60]$,
$[60,80]$, and $[80,100]$.  Their exact norm floor is
\[
 \frac{22271745300542300149807}
 {85070591730234615865843651857942052864}>10^{-16}.
\]
Exact rational cover checks include the spatial subdivisions,
all shared boundaries, and both time endpoints.  The spatial overlaps
$[7,8]$, $[40,41]$, and the shared boundary $x=100$, together with
evenness, prove \eqref{fc:height104-gap}.

The last band also certifies the centered-residual mechanism.  On each
of its 32 cells, freeze
\[
 a=\frac\pi{16},\qquad
 b=\frac14\log\frac{x_c+2t_ca}{4\pi\sin(4a)},\qquad
 \mu=a^2+b^2,
 \qquad C_j=-h_{j+2}-2ah_{j+1}-\mu h_j.
\]
The complete even source is shifted before theta indices are separated.
With eight retained indices, integration cutoff $3/2$, spatial order
12 and time order 8, the computation encloses all center jets through
order 28 and complete derivative majorants through order 29.
Centered coefficients are combined before Taylor transport, with
the parameters frozen on the cell.  The cell-constant normalization
$e^{\pi x_c/8}$ rescales each physical jet equally.
For normalized enclosures $P_0,P_1,Q_0,Q_1$ of $h_0,h_1,C_0,C_1$,
the actual decisions are
\[
 E_{\rm low}=\operatorname{dist}(0,P_0)^2+
 \operatorname{dist}(0,P_1)^2>0,\qquad
 E_{\rm low}/9-\sup|Q_0|^2-\sup|Q_1|^2>0.
\]
Thus $C_0^2+C_1^2\le E/9$ and $\mu\ge2/5$ throughout every cell.
The exact stationary identity
$\mathcal L_j=\mu h_j^2+h_jC_j$ at $h_{j+1}=0$ yields
\[
 \mathcal L_0\ge E/15\quad(h_1=0),\qquad
 \mathcal L_1\ge-5E/72\quad(h_2=0).
\]
The second strict stationary sign in
Theorem~\ref{fc:height104} is a separate validated signed decision;
it is stronger than the bounded-loss conclusion supplied by this cone.
At each level the band has 24 derivative-exclusion cells and eight
strict signed cells, with no claim that a signed cell alone proves
existence or uniqueness of a stationary branch.

Appendix~\ref{app:replay} identifies the retained drivers, records,
exact auditors, and earlier independent internal replays. In particular,
the higher-height auditor checks 6682 assertions and the centered-band
auditor checks 547. Those checks concern arithmetic, decisions, complete
coverage, and source bindings. The complete-source contour bounds used
for the latest band are proved in Appendix~\ref{app:payments}.
The computations establish the stated closed bounded region; they do
not extend it to arbitrary heights.


% Prepared for Edward Baker, 14 October 2026.
% LLM assistance: gpt-6-astra (Codex), reasoning effort ultra.
% Replacement for the former Section 7, integrating Notes 21--22.

\subsection{Extension through height 112 and all earlier-time comparisons}
\label{subsec:height112}

\begin{theorem}[Extended closed bounded certificate]
\label{fc:height112}
For the complete theta source,
\[
 \sqrt{h_0(t,x)^2+h_1(t,x)^2}>10^{-17}
       \qquad(0\le t\le1/2,\ |x|\le112).
\]
Both stationary implications $h_{j+1}=0\Rightarrow\Lag_j>0$,
$j=0,1$, hold on this rectangle.  With the fixed centering
$a=\pi/16$, $b=1/4$, $\mu=a^2+b^2$, the following stronger delayed
statements hold on the positive-height extension $104\le x\le112$:
at every current stationary point and every $0\le s\le t$,
\begin{equation}
 h_j(t,x)C_j(s,x)>0,\qquad
 h_j(t,x)\{C_j(t,x)-C_j(s,x)/2\}>0.
 \label{eq:height112-causal}
\end{equation}
There is exactly one $h_1=0$ point in $(107,108)$ and exactly one
$h_2=0$ point in $(109,110)$ at each time in $[0,1/2]$; there are
no other stationary points of these respective observations in
$[104,112]$.
\end{theorem}
\begin{proof}[Proof from the complete-source interval certificates]
The two band records in the numerical archive use the complete source
and error bounds already specified above:
\begin{quote}\small\raggedright
\path{numerics/FIXED_CUSP_BAND_104_108_20261014.json}\\
\path{numerics/FIXED_CUSP_BAND_108_112_20261014.json}.
\end{quote}
Each
has four unit spatial cells and eight closed time slabs
$[k/16,(k+1)/16]$, $0\le k\le7$.  Their fixed numerical parameters
are 192-bit arithmetic in python-flint 0.9.0, $a=\pi/16$, retained
index $N=8$, integration endpoint $U=3/2$, spatial order 12, and time
order 8.  Center physical jets through order 28 and complete moment
payments through order 29 include quadrature, full-index, integration,
rounding, spatial, and temporal errors.

The exact audit proves the whole-cell physical norm bound and the
normalized lower bound $e^{\pi x/4}E>3000$ on both new bands.  The stated
physical floor follows from the stronger cell-specific norm enclosures
and a rational exponential estimate using $\pi<22/7$; the displayed
coarse normalized bound alone is not used to obtain that floor.  The 128 stationary decisions are
\begin{center}
\begin{tabular}{ccl}
Band&Level&Complete closed-cell decisions\\\hline
$[104,108]$&$j=0$&24 exclusions; 8 strict signs in $(107,108)$\\
$[104,108]$&$j=1$&32 exclusions\\
$[108,112]$&$j=0$&32 exclusions\\
$[108,112]$&$j=1$&24 exclusions; 8 strict signs in $(109,110)$.
\end{tabular}
\end{center}
In each candidate spatial cell, paid endpoint signs of $h_{j+1}$ are
opposite and the full-cell $h_{j+2}$ enclosure has one strict sign.
This proves existence and uniqueness throughout the closed time slab.
The other cells exclude stationarity, so the implications cover branch
creation as well as regular branches.  The shared boundaries at 104 and
108 and the preceding certificate leave no spatial gap.  Physical parity
reflects the norm and stationary signs to negative heights.

For the delayed products the auditor reconstructs the residual from the
physical jets with $b=1/4$, rather than the saddle centering used for
the initial stationary test.  For current slab $k$, the earlier slabs
$p=0,\ldots,k$ cover all $0\le s\le t$, conservatively including
some $s>t$ in slab $k$.  The current and earlier time displacements
$d,e\in[-1/32,1/32]$ are independent; the spatial displacement
$\eta\in[-1/2,1/2]$ is the same in both observations.  If their paid
Taylor polynomials are $T_k(\eta,d)$ and $T_p(\eta,e)$, form the
coefficients of $T_k(\eta,d)-T_p(\eta,e)/2$ before ranging the common
$\eta$.  Explicitly, for coefficients $c_{k,r,m}$, the polynomial is
\[
 \sum_{m=0}^{11}\left[c_{k,0,m}-\frac12c_{p,0,m}
 +\sum_{r=1}^{7}\left(c_{k,r,m}d^r-\frac12c_{p,r,m}e^r\right)
 \right]\eta^m.
\]
Add the full current residual payment plus half the earlier payment;
only at the current time intersect with
$C_j=-h_{j+2}-\mu h_j$.  No earlier stationarity is imposed.
Taking the union over all earlier slabs and multiplying the complete
integrated observations preserves all source cross terms.

In the common constant normalization $Z_j=e^{\pi x_c/8}h_j$, with
the same factor for each residual, the final lower bounds for all
earlier times are
\begin{center}
\begin{tabular}{ccc}
Current branch&$Z_j\widehat C_j(s,x)$&
$Z_j(\widehat C_j(t,x)-\widehat C_j(s,x)/2)$\\\hline
$h_1=0$, $x_c=107.5$&$>15511$&$>7559$\\
$h_2=0$, $x_c=109.5$&$>6910$&$>3371$.
\end{tabular}
\end{center}
These enclosures prove \eqref{eq:height112-causal}.  They use complete
earlier fields, so they do not apply a multiplier bound restricted to
$\delta\le1/4$ when $t-s$ is larger.
\end{proof}

The earlier uniform residual-cone enclosure remains undecided on 24
cells in the first band and 16 in the second, all outside the candidate
stationary cells.  Those original verdicts are retained.  The successful
norm and stationary decisions do not turn the failed cone enclosures
into cone proofs or negative physical signs.  Likewise the delayed
products in \eqref{eq:height112-causal} are asserted for the positive-height
protocol actually computed; fixed drift in $C_j$ prevents an automatic
unchanged-drift reflection of that numerical claim.

The generating driver is
\path{numerics/certify_fixed_cusp_band_extension.py}, run with
\texttt{--x-left 104} and \texttt{--x-left 108}, each with an
\texttt{--output} path.  The driver
\path{numerics/audit_fixed_cusp_band_extension.py} takes both records
through repeated \texttt{--record} arguments and an \texttt{--output}
path.  Its retained
\path{numerics/FIXED_CUSP_BAND_EXTENSION_AUDIT_20261014.json}
contains 1321 passing checks, all 128 stationary observations, and all
16 genuine-branch comparisons.  Fresh 192-bit replays reproduced both
generator records after omitting only \texttt{elapsed\_seconds}; the
exact audit was reproduced in full.  Source bindings and these historical
replay claims are documented in
\path{numerics/PHASE_CONTROLLED_CUSP_CONTINUATION_CHECKS_20261014.json}
and Review 19.  They are internal validations, not formal verification
of the analytic source lemmas or the interval library.

\section{One-sided stationary estimates and their obstructions}
\label{sec:remaining}

The outer theorem and the finite cover give complementary information.
The unresolved positive-time domain is contained in the transition strip
$112<|x|<X_*$ and the high-height inner region
\begin{equation}
 0<t\le\tfrac12,\qquad |x|\ge X_*,\qquad
 t\log(|x|/(4\pi))<\frac{61}{24}.
 \label{eq:remaining-domain}
\end{equation}
The closed-boundary assertion is asymptotic: it does not supply a usable
numerical value of $X_*$. The separate explicit threshold-five bound is
not a numerical implementation of this sharper theorem. Even the
conditional completion down to $5/2$ leaves a smaller unbounded cusp.
The remaining analytic problem is a signed estimate in this domain,
with coefficients locally bounded through a candidate positive-time
collision. The following obstructions delimit that problem.

\subsection{Absolute dominance and pointwise positivity are too strong}

At the exact common-carrier coordinate $v$, put
$M=\lfloor e^{2v}\rfloor$. For $\lceil M/2\rceil\le n\le M$,
\[
 td_n^2-2td_nv=t[(d_n-v)^2-v^2]\ge-tv^2.
\]
Consequently
\begin{equation}
 \sum_{n\le e^{2v}}|A_n|
 \ge(M/2-1)M^{-1/2}e^{-tv^2}.
 \label{eq:absolute-channel-obstruction}
\end{equation}
With fixed $a$, $v=b_1$, and $t=\vartheta/\log x$ for fixed
$0\le\vartheta<4$, this grows as a positive constant times
$x^{1/4-\vartheta/16}$, whereas $|A_1|=1$. Absolute first-channel
dominance therefore fails in this scaling. Cancellation within the sum
and after integration remains available; indeed it is essential to the
improved outer theorem.

Pointwise positivity is also unavailable. For fixed $v=b$, write
\[
 D_j=\sum_{n\le e^{2b}}A_n(t,x;b)[-a+i(b-d_n)]^j,\qquad p_j=\Re(UD_j),
\]
where $U$ is the unit phase of
$e^{t(b+ia)^2}\phi_1(b+ia)e^{ixb}$. Put
\[
 \mathfrak H_j=|D_{j+1}|^2-\Re(D_j\overline{D_{j+2}}),\qquad
 \mathfrak W_j=D_{j+1}^2-D_jD_{j+2}.
\]
Then
\begin{equation}
 2(p_{j+1}^2-p_jp_{j+2})
 =\mathfrak H_j+\Re(U^2\mathfrak W_j).
 \label{eq:phase-projection}
\end{equation}
The stronger phase-independent condition
$\mathfrak H_j\ge|\mathfrak W_j|$ fails for the actual common-carrier
sums. With $a=\pi/16$ and $b=b_1$ from
\eqref{eq:channel-centers}, retained Arb/Acb enclosures give
\begin{center}
\begin{tabular}{@{}lll@{}}
\toprule
$x$ & $t$ & $\mathfrak H_0-|\mathfrak W_0|$ (approximate value)\\
\midrule
$211$ & $1/\log211$ & $-0.0277198$\\
$4898.7$ & $2/\log4898.7$ & $-1.26105$\\
\bottomrule
\end{tabular}
\end{center}
At the second point the actual phase also gives
$p_1^2-p_0p_2\approx-0.01099197<0$. The stored enclosures are strictly
negative; Appendix~\ref{app:replay} identifies their replay. These are
pointwise integrand forms, not the integrated Laguerre forms, and no
physical stationarity is asserted at these points. The paired identities
in Section~\ref{sec:residuals} retain correlations between distinct
coordinates as well as distinct indices. A proof requiring every
integrand to be positive would discard that possibility.
The staircase boundary identity \eqref{eq:staircase-cancellation} must
likewise be retained: entering terms have an exact cancellation of their
arithmetic $x$-phase, not a freely discardable endpoint error.

\subsection{A sufficient split that preserves favorable curvature}

Throughout this section let $0<a<\pi/8$, prescribe a real $b$, and put
$\mu=a^2+b^2$. At each point the centered observations are
\begin{equation}
 C_k=-h_{k+2}-2ah_{k+1}-\mu h_k,\qquad k\ge0.
 \label{eq:all-centered-observations}
\end{equation}
Whenever these observations are differentiated or compared at different
times below, $a,b$ are held fixed. The exact identity is
\[
 \Lag_j=\mu h_j^2+h_jC_j+h_{j+1}(h_{j+1}+2ah_j).
\]
Thus the remaining target can be stated as
\begin{equation}
 \boxed{h_jC_j\ge-\mu h_j^2-K_jE
 \quad\hbox{on }h_{j+1}=0,\qquad j=0,1,}
 \label{eq:final-target}
\end{equation}
where $K_j\ge0$ is prescribed without dividing by an unknown energy and
is locally bounded through every candidate positive-time collision.
This is precisely the stationary loss condition in
Section~\ref{sec:fc-collision}.

\begin{lemma}[Favorable curvature and a residual square]
\label{lem:favorable-split}
Prescribe an exact split $C_j=G_j+R_j$ and parameters
$0<\kappa\le\mu$, $\rho\ge0$. Define
\begin{equation}
 \Delta_j=4\kappa h_jG_j+\rho^2E-R_j^2.
 \label{eq:split-slack}
\end{equation}
Without imposing stationarity,
\begin{align}
 4\kappa\Lag_j+\rho^2E
 &=(R_j+2\kappa h_j)^2+\Delta_j
       +4\kappa(\mu-\kappa)h_j^2 \notag\\
 &\quad+4\kappa h_{j+1}(h_{j+1}+2ah_j).
 \label{eq:split-stationary-ideal}
\end{align}
In particular, $\Delta_j\ge0$ on $h_{j+1}=0$ implies
$\Lag_j\ge-\rho^2E/(4\kappa)$. This loss is locally bounded if
$\rho$ is locally bounded and $\kappa$ is locally bounded below by a
positive constant.
\end{lemma}
\begin{proof}
Expand the square and use $G_j+R_j=C_j$. All terms are polynomial in
the physical jets, so neither $E>0$ nor division by a jet is required.
On stationarity the last term vanishes and the other displayed squares
and $4\kappa(\mu-\kappa)h_j^2$ are nonnegative.
\end{proof}

The split must expose an independently estimable source contribution.
For the exact domain and fields of \eqref{eq:pair-domain}, choose a
prescribed real multiplier $m_j(\alpha)$ with at most polynomial growth
in $y_\alpha$, and set
\begin{align}
 G_j&=e^{-ax}\Re\int_\Omega m_j(\alpha)W_\alpha q_\alpha^j\dd\nu(\alpha),
       \notag\\
 R_j&=e^{-ax}\Re\int_\Omega r_j(\alpha)W_\alpha q_\alpha^j\dd\nu(\alpha),
 \qquad r_j(\alpha)=y_\alpha^2-b^2-m_j(\alpha).
 \label{eq:prescribed-source-split}
\end{align}
A proper source filter takes
$m_j=(y_\alpha^2-b^2)\mathbf1_A$, with the retained set $A$ intersected
with the original staircase. The filter and its parameters are fixed
before evaluating the physical observations. For real $m_j,r_j$, the
polarization \eqref{eq:paired-J} gives the exact scalar obligation
\begin{align}
 \Delta_j=\frac{e^{-2ax}}2\Re\iint_{\Omega^2}
 \big[&\{2\kappa(m_j(\alpha)+m_j(\beta))
                  -r_j(\alpha)r_j(\beta)\}J_j(\alpha,\beta)\notag\\
       &+\rho^2(J_0(\alpha,\beta)+J_1(\alpha,\beta))\big]
       \dd\nu(\alpha)\dd\nu(\beta).
 \label{eq:paired-split-slack}
\end{align}
Indeed $h_jG_j$ is $e^{-2ax}/4$ times the real part of the integral
of $(m_j(\alpha)+m_j(\beta))J_j$, whereas $R_j^2$ has coefficient
$e^{-2ax}/2$ and integrand $r_j(\alpha)r_j(\beta)J_j$.
The complete source majorants justify these identities and their tail
payments; they do not prove the sign in \eqref{eq:paired-split-slack}.
The formula requires real multipliers. A complex filter requires its
own conjugated polarization or an independent physical-field identity.

Even a quotient-free split can be vacuous: choosing
$m_j=y_\alpha^2-b^2+2\kappa$ makes $R_j=-2\kappa h_j$ and cancels
the first square in \eqref{eq:split-stationary-ideal}. With $\kappa=\mu$,
the resulting slack condition only restates the target and supplies no
theta estimate.

\subsection{Collision controls and the limitation of the square criterion}

For the backward-heat polynomial $P(t,y)=y^2-d^2-2t$, put
$D=d^2+2t>0$. Its Laguerre forms are $2y^2+2D>0$ and $4$.
At $y=0$, $h_0=-D$, $h_1=0$, $E=D^2$, $C_0=-2+\mu D$, and
$C_1=-4a$. The whole-residual cone would require
\begin{equation}
 \eta\ge\frac{\sqrt{(2-\mu D)^2+16a^2}}{D}.
 \label{eq:close-pair-cone}
\end{equation}
For bounded $a,\mu$, this grows without bound as the pair closes,
although the required negative loss is zero. The prescribed split
$G_0=-2$, $R_0=\mu D$ instead gives
\[
 \Delta_0=8\kappa D+(\rho^2-\mu^2)D^2.
\]
It accepts $\rho=0$ whenever $0<D\le8\kappa/\mu^2$, preserving the
large favorable curvature. This control is not a theta counterexample;
growth as $t\downarrow0$ is also compatible with local boundedness at
each positive time.

The same algebra rejects adverse collision predecessors. For
$P=y^2-2(t-T)$ at $y=0$, $t=T-s<T$, the split
$G_0=-2$, $R_0=-2\mu s$ yields
\[
 \Delta_0=-16\kappa s+4(\rho^2-\mu^2)s^2,\qquad
 \Lag_0=-4s,\qquad E=4s^2.
\]
For $P=y^3-6(t-T)y$ at $y=0$, $t=T-s<T$, the second stationary
set $h_2=0$ has $h_1=6s$, $h_3=6$. Taking $G_1=-6$,
$R_1=-6\mu s$ gives
\[
 \Delta_1=-144\kappa s+36(\rho^2-\mu^2)s^2,\qquad
 \Lag_1=-36s,\qquad E=36s^2.
\]
Both slacks are negative for sufficiently small $s>0$ when $\rho$
is bounded and $\kappa$ stays positively bounded below. The second
stationary set cannot be omitted.

There is also a precise restriction at the candidate collision itself:
if $h_0=h_1=0$, then $\Delta_0=-R_0^2$. The square criterion therefore
forces $R_0=0$ exactly; if also $h_2=0$, it forces $R_1=0$.
A small absolute bound on an omitted source does not imply this
vanishing. Thus Lemma~\ref{lem:favorable-split} is a sufficient
criterion with an extra cancellation obligation, not a solution of the
remaining theta problem.

Nor can its sign be justified by arbitrary-source positivity. On
$h_1=0$ the arbitrary-jet form $\Lag_0+KE=-h_0h_2+Kh_0^2$ is
indefinite for every finite $K$; on $h_2=0$ the form is
$-h_1h_3+K(h_0^2+h_1^2)$ and is likewise indefinite. Terms in the
stationary ideal vanish there. Any useful positivity argument must
therefore use relations of the actual theta source.

\subsection{Direct product absorption}

The residual-square condition can be replaced by a one-sided product
condition which imposes no residual vanishing at $E=0$.
\begin{lemma}[Fixed-fraction absorption]\label{lem:direct-absorption}
For a prescribed split $C_j=G_j+R_j$, let $0\le\lambda\le1$ and
$K_j\ge0$. On the genuine stationary set $h_{j+1}=0$, suppose
\begin{equation}
 h_jG_j\ge0,\qquad h_j(R_j+\lambda G_j)\ge-K_jE.
 \label{eq:direct-absorption}
\end{equation}
Then $\Lag_j\ge-K_jE$.
\end{lemma}
\begin{proof}
On this set,
\[
 \Lag_j=\mu h_j^2+(1-\lambda)h_jG_j
                      +h_j(R_j+\lambda G_j)\ge-K_jE.
\]
\end{proof}
At a first-pair zero the displayed products vanish automatically, but
the estimate must still hold with a locally bounded prescribed loss in
a neighborhood of that point. A favorable retained product alone does
not suffice. The finite results below use $\lambda=1/2$ and
$K_0=K_1=0$, without fitting a loss to an observed energy floor.


\section{Proper source filters and complete theta correlations}
\label{sec:source-filters}

We first test the preceding mechanisms on populated stationary branches
inside the existing height-$104$ cover. We then prove why fixed sharp
filters cannot support uniformly small separate payments at arbitrary
heights, and derive a complete-source alternative. The finite tests do
not extend the bounded cover or settle \eqref{eq:remaining-domain}.

\subsection{Two certified stationary branches and a signed comparison}

Fix
\begin{equation}
 a=\pi/16,\qquad b=1/4,\qquad \mu=a^2+b^2,\qquad
 \lambda=1/2,\qquad K_0=K_1=0.
 \label{eq:filter-pilot-parameters}
\end{equation}
Use the constant normalization $Z_k=e^{\pi x_c/8}h_k$, with
$x_c=101.5$ on $[101,102]$ and $x_c=103.5$ on $[103,104]$.
Hats denote the same normalization of centered observations.
The complete-source certificate and paid endpoint Taylor transport give,
for every $t\in[0,1/2]$,
\begin{align}
 Z_1(t,101)>37,\qquad& Z_1(t,102)<-20,\qquad Z_2<0
                       \quad(101\le x\le102),\notag\\
 Z_2(t,103)<-26,\qquad& Z_2(t,104)>14,\qquad Z_3>0
                       \quad(103\le x\le104).
 \label{eq:pilot-stationary-endpoints}
\end{align}
The intermediate value theorem and strict monotonicity therefore prove
exactly one $h_1=0$ point in $(101,102)$ and exactly one $h_2=0$ point
in $(103,104)$ at each time. The implicit function theorem gives smooth
branches. These are stationary points of the physical derivatives,
not derivatives of a variable normalization.

A first signed comparison uses the full residual on the first branch.
Exact rational extraction from the eight stored closed time slabs gives
\begin{equation}
 \begin{gathered}
 Z_0>0,\qquad\widehat C_0>0,\qquad
 Z_0\widehat C_0>6000,\qquad\widehat C_0-\mu Z_0>7/5,\\
 \mu Z_0^2+Z_0\widehat C_0>9400,\qquad
 \mu^2Z_0^2-\widehat C_0^2-\widehat C_1^2<-78.
 \end{gathered}
 \label{eq:full-residual-signed-pilot}
\end{equation}
Thus the signed condition holds with $K_0=0$, whereas every prescribed
whole-residual cone budget $\eta\le\mu$ fails on this branch.
The previously certified saddle-centered $\eta=1/3$ cone still passes;
the comparison concerns the different fixed centering
\eqref{eq:filter-pilot-parameters}, not a failure of every norm cone.
Taking $G=C$, $R=0$ here is a full-residual comparison, not a proper
source split. The next experiment supplies such a split.

\subsection{A proper filter absorbs an adverse complement on 64 cells}

Before evaluating partial integrals, prescribe the four filters
\[
 (N,U)\in\{(1,1),(1,3/2),(2,1),(2,3/2)\}.
\]
For $z=u+ia$, define
\begin{align}
 g_k(t,x)&=e^{-ax}\Re\sum_{n=1}^N\int_0^U
             (iz)^k e^{tz^2+ixu}\phi_n(z)\dd u,\notag\\
 G_j&=-g_{j+2}-2ag_{j+1}-\mu g_j,\qquad R_j=C_j-G_j.
 \label{eq:rectangular-filter}
\end{align}
This is precisely the real source multiplier
$(u^2-b^2)\mathbf1_{n\le N,\,0\le u\le U}$. Its shared-carrier
domain is $n\le N$, $d_n\le v\le d_n+U$, so all entry boundaries
are retained. Each of the eight time slabs
$[r/16,(r+1)/16]$, $0\le r\le7$, is tested on both branches:
four filters, two branches, and eight slabs give $64$ closed-cell tests.
Even-index slabs were designated development and odd-index slabs holdout
before evaluation; no filter, parameter, or loss was changed between them.
Their interiors are disjoint and their common boundary times are covered.

Both values of $U$ give the following conservative bounds on every
corresponding slab. Products are normalized by $e^{\pi x_c/4}$.
\begin{center}
\begin{tabular}{@{}ccrrr@{}}
\toprule
$N$ & $j$ & $Z_j\widehat G_j$ & $Z_j\widehat R_j$
    & $Z_j(\widehat R_j+\widehat G_j/2)$\\
\midrule
$1$ & $0$ & $>0$ & $<0$ & $<-1{,}200{,}000$\\
$1$ & $1$ & $>0$ & $<0$ & $<-250{,}000$\\
$2$ & $0$ & $>7800$ & $(-212,-120)$ & $>3800$\\
$2$ & $1$ & $>3800$ & $(-46,-25)$ & $>1900$\\
\bottomrule
\end{tabular}
\end{center}
Here $j=0$ means the $h_1=0$ branch and $j=1$ the $h_2=0$ branch.
The two-index filters therefore pay a genuinely adverse complementary
product with the prescribed half of a favorable retained product.
Lemma~\ref{lem:direct-absorption} applies with the fixed zero loss.
For one index the same prescribed absorption test actually fails on
nonempty stationary sets: the negative upper bounds are not an interval
precision failure. The large one-index terms nearly cancel in the
complete residual.

For completeness, the certification encloses retained center derivatives
through order $28$ with $192$-bit Arb quadrature, then uses space order
$12$ and time order $8$ for transport. Fixed endpoints give
$\partial_xg_k=g_{k+1}$ and $\partial_tg_k=-g_{k+2}$ exactly.
The complete-source moment payments remain valid for the three masks
$\mathbf1_A$, $1-\mathbf1_A$, and $1-\mathbf1_A/2$, all of which
lie in $[0,1]$. Centered payments are bounded by
$P_{j+2}+2aP_{j+1}+\mu P_j$ for each spatial or temporal remainder.
The full $C$ arrays already include quadrature error, all omitted theta
indices, and the infinite integration tail. The complement is the exact
difference $C-G$, and $C-G/2$ is formed before transport to preserve
cancellation. On current stationarity its enclosure is intersected with
$-h_{j+2}-\mu h_j-G_j/2$; the complement is treated similarly.
The product always uses the full physical jet and an integrated
observation, retaining all cross terms. No individual theta index is
assigned the vanishing vertical contour integral of the complete source.
The complete $64$-case receipt, including failures, is identified in
Appendix~\ref{app:replay}.

\subsection{The endpoint terms of a fixed sharp filter}

Now fix $t\in(0,1/2]$, $0<a<\pi/8$, $b\ge0$, $N\ge1$, and
$U>0$, independently of $x$. Write
\[
 \Phi_N(z)=\sum_{n=1}^N\phi_n(z),\qquad q=-a+iu=iz,\qquad
 f_{N,j}(u)=(u^2-b^2)q^j e^{tz^2}\Phi_N(z).
\]
Then the retained residual in \eqref{eq:rectangular-filter} is
\begin{equation}
 G_j(x)=e^{-ax}\Re\int_0^U f_{N,j}(u)e^{ixu}\dd u,\qquad R_j=C_j-G_j.
 \label{eq:sharp-filter-integral}
\end{equation}
These fixed filtered fields still solve backward heat. The following
effect comes from source integration by parts, not a heat-equation defect.

\begin{lemma}[Complete endpoint expansion]\label{lem:sharp-endpoints}
For $x>0$,
\begin{align}
 e^{ax}G_j(x)&=\frac{B_j(x)}x+\mathcal E_j(x),\qquad
 B_j(x)=\Im\{f_{N,j}(U)e^{ixU}-f_{N,j}(0)\},\notag\\
 |\mathcal E_j(x)|&\le\frac{M_{N,j}}{x^2},\qquad
 M_{N,j}=|f'_{N,j}(U)|+|f'_{N,j}(0)|
                      +\int_0^U|f''_{N,j}(u)|\dd u.
 \label{eq:sharp-endpoint-expansion}
\end{align}
In particular the lower coefficient is
\begin{equation}
 -\Im f_{N,j}(0)=b^2(-a)^j e^{-ta^2}\Im\Phi_N(ia).
 \label{eq:raw-index-endpoint}
\end{equation}
\end{lemma}
\begin{proof}
Twice integrating $\int_0^Ufe^{ixu}\dd u$ by parts gives
\[
 \frac{[fe^{ixu}]_0^U}{ix}
 -\frac{[f'e^{ixu}]_0^U}{(ix)^2}
 +\frac{1}{(ix)^2}\int_0^Uf''e^{ixu}\dd u.
\]
Taking real parts uses $\Re(-iZ)=\Im Z$ and gives the stated bound.
Substitution at $u=0$ proves \eqref{eq:raw-index-endpoint}.
\end{proof}

For the complete theta source, evenness and Schwarz symmetry make
$\Phi(ia)$ real. More generally its centered integrand obeys
$f_j(-u)=\overline{f_j(u)}$. Its even derivatives at zero are real
and its odd derivatives purely imaginary, so every lower-endpoint
coefficient in the real integration-by-parts expansion vanishes.
A finite raw-index sum need not have this symmetry.
If instead a common cutoff $v\le V$ is used, retain
$U_n=V-d_n>0$ for every included index; the leading term becomes
\[
 \Im\left\{\sum_{\substack{n\le N\\d_n<V}}
 f_{n,j}(U_n)e^{ixU_n}
 -\sum_{\substack{n\le N\\d_n<V}}f_{n,j}(0)\right\}.
\]
Here $f_{n,j}$ is defined with the individual term $\phi_n$ in place
of $\Phi_N$. The remainder is the sum of the bounds in
\eqref{eq:sharp-endpoint-expansion}. An index with $d_n=V$ has an
empty interval and contributes zero, not a single unmatched endpoint.
This is the same entry cancellation as
\eqref{eq:staircase-cancellation}.

The upper endpoint is provably nonzero for the pilot family, without a
genericity assumption. More generally suppose
$a=\pi/16$, $U\ge1$, and $U\ne b$. Put $D_4=e^{4U}>50$ and
$\beta=\pi\cos(4a)D_4>100$. For $n\ge2$,
\[
 \left|\frac{\phi_n(U+ia)}{\phi_1(U+ia)}\right|
 \le n^2\frac{2\pi n^2D_4+3}{2\pi D_4-3}
             e^{-\beta(n^2-1)}
 \le2n^4e^{-100(n^2-1)}.
\]
Here $\phi_1(U+ia)\ne0$ since
$|2\pi e^{4(U+ia)}-3|\ge2\pi D_4-3>0$.
Using $n^4\le e^{n^2-1}$ and
$n^2-1\ge3+5(n-2)$ gives
\begin{equation}
 \sum_{n=2}^\infty2n^4e^{-100(n^2-1)}
 \le\frac{2e^{-297}}{1-e^{-495}}<\tfrac12.
 \label{eq:nonzero-sharp-endpoint}
\end{equation}
Thus $\Phi_N(U+ia)\ne0$ for every finite $N$, and
$f_{N,0}(U)\ne0$. This includes $U=1,3/2$, $b=1/4$.

\begin{theorem}[Adverse sharp-filter products on genuine stationary sets]
\label{thm:sharp-filter-obstruction}
Fix the parameters of \eqref{eq:sharp-filter-integral} and assume
$f_{N,0}(U)\ne0$. For each $j=0,1$ there is a sequence
$x_k\to\infty$ of genuine stationary points $h_{j+1}(t,x_k)=0$ such
that $h_jG_j>0$, $h_jR_j<0$, and
\begin{equation}
 \frac{-h_jR_j}{E}(t,x_k)
   =\exp\{(\pi/8-a)x_k+o(x_k)\}.
 \label{eq:sharp-filter-adverse-rate}
\end{equation}
For every fixed $0\le\lambda<1$, along the same sequence eventually
\begin{equation}
 h_j(R_j+\lambda G_j)<0,\qquad
 \frac{-h_j(R_j+\lambda G_j)}E
   =\exp\{(\pi/8-a)x_k+o(x_k)\}.
 \label{eq:sharp-filter-absorption-rate}
\end{equation}
There are also stationary sequences with $h_jR_j>0$ and the same
exponential rate for its absolute ratio. All these points eventually
belong to the proved outer region, where $E>0$ and both complete
Laguerre forms are positive.
\end{theorem}
\begin{proof}
At fixed $t>0$, the outer four-jet estimates apply for all sufficiently
large $x$. With $A(x)=2|Q_t(x)|$, $\omega=\log(x/(4\pi))/4$, and
a continuous phase $\varphi=\arg Q_t$, their fixed-time consequences are
\begin{equation}
 \begin{gathered}
 h_k=A\omega^k\{\cos(\varphi+k\pi/2)+O(1/\log x)\},
       \qquad 0\le k\le3,\\
 \varphi'(x)=\omega+O(1)>0,\qquad
 \log A(x)=-\pi x/8+O_t((\log x)^2).
 \end{gathered}
 \label{eq:fixed-time-phase-windows}
\end{equation}
For $j=0$, take phase windows centered at $\varphi=k\pi$ with
half-width $\pi/6$. For all sufficiently large windows, $h_1$ has
opposite endpoint signs and $h_2$ has a fixed nonzero sign throughout.
There is exactly one $h_1=0$ point in each window, and the sign of
$h_0$ alternates. There $|h_0|\asymp A$ and $E=h_0^2\asymp A^2$.

For $j=1$, center the windows at $\varphi=\pi/2+k\pi$. Now $h_2$
has opposite endpoint signs and $h_3$ has a fixed nonzero sign.
There is exactly one $h_2=0$ point in each window and $h_1$ alternates
sign. At these points $|h_1|\asymp A\omega$, while $h_2=0$ in
\eqref{eq:fixed-time-phase-windows} gives $h_0=O(A/\log x)$;
therefore $E\asymp A^2\omega^2$.

Since $q(U)\ne0$, the hypothesis makes each $B_j$ a nonconstant
sinusoid plus a constant. It has infinitely many fixed-width intervals
on which its sign is fixed and $|B_j|\ge c_j>0$. The spacing of the
stationary points just constructed is $O(1/\log x)$. Each such interval
at sufficiently large height therefore contains stationary points of
both signs of $h_j$. By \eqref{eq:sharp-endpoint-expansion}, $G_j$
has the sign of $B_j$ there and $|G_j|\asymp e^{-ax}/x$.
In contrast, $C_j$ is a fixed linear combination of complete physical
jets, so \eqref{eq:fixed-time-phase-windows} makes it exponentially
smaller than $G_j$. Hence $R_j=-G_j(1+o(1))$ on these intervals.

Choose the stationary points with $h_jG_j>0$ or with $h_jG_j<0$.
The energy estimates above give \eqref{eq:sharp-filter-adverse-rate}
and its opposite-sign version: all powers of $x,\log x$ and the
$O_t((\log x)^2)$ amplitude correction contribute $o(x)$ to the
logarithm. Finally
$h_j(R_j+\lambda G_j)=h_jC_j-(1-\lambda)h_jG_j$, and the same
comparison proves \eqref{eq:sharp-filter-absorption-rate}.
\end{proof}

Thus no constant, polynomial, or other budget $K_j(x)=\exp(o(x))$
can make the second inequality in \eqref{eq:direct-absorption} hold
globally for this fixed sharp filter and a fixed $\lambda<1$.
This conclusion does not exclude losses locally bounded at every finite
positive-time point: exponential growth in $x$ is still locally bounded.
It is a fixed-positive-time obstruction in an already favorable region,
not an obstruction to an adaptive filter in the unresolved moving cusp.
Neither $\lambda=1$ nor a fraction tending exponentially fast to one is
covered. In the complete product $h_jC_j$, the artificial endpoint
terms cancel exactly.

\subsection{A compact filter has no finite rational jet closure}

\begin{proposition}\label{prop:no-rational-filter-closure}
For any fixed $t,a,b,N,U$ as in
\eqref{eq:sharp-filter-integral}, no finite family of real rational
functions $r_0,\ldots,r_m$ can satisfy
\begin{equation}
 R_j(x)=\sum_{k=0}^m r_k(x)h_k(t,x)
 \label{eq:forbidden-rational-closure}
\end{equation}
on a nonempty open real interval free of their poles, for $j=0$ or $1$.
\end{proposition}
\begin{proof}
Since $C_j$ is a constant-coefficient combination of physical jets,
clearing denominators would give
\[
 p(x)G_j(x)=\sum_{k=0}^m p_k(x)h_k(t,x),\qquad p\ne0,
\]
with polynomials. Both sides extend to entire functions, so this
identity holds for all real $x$.
Complete evenness and Schwarz symmetry give
\begin{equation}
 e^{ax}h_k(t,x)=\frac12\int_{\R}
 e^{t(u+ia)^2}\Phi(u+ia)[i(u+ia)]^k e^{ixu}\dd u.
 \label{eq:whole-line-filter-closure}
\end{equation}
The Fourier density is real analytic on all of $\R$, and it and all its
derivatives decay rapidly: holomorphy follows from the open theta strip,
positive-end decay from the theta series, and negative-end decay from
the complete evenness. Polynomial multiplication in $x$ corresponds
to finite derivatives in $u$, preserving real analyticity.

By contrast, $e^{ax}G_j$ is the Fourier transform of the compactly
supported $L^1$ density equal to $f_{N,j}(u)/2$ on $(0,U)$,
$\overline{f_{N,j}(-u)}/2$ on $(-U,0)$, and zero elsewhere.
Multiplication by $p(x)$ gives a compactly supported distribution.
Fourier injectivity for tempered distributions equates this compact
distribution to a real analytic function. That function vanishes on
$(U,\infty)$, hence everywhere. Therefore $pG_j=0$ and $G_j=0$.

This is impossible. The finite sum $\Phi_N$ is a nonzero holomorphic
function, since every $\phi_n$ is positive on the positive real axis.
The other factors in $f_{N,j}$ have only isolated zeros, so its compact
Hermitian extension is not zero. Fourier injectivity implies $G_j$ is
not identically zero. The injectivity used here follows, for example,
by Gaussian regularization and the approximate-identity limit; it
requires no information about zeros of $H$.
\end{proof}

The proposition excludes rational-coefficient finite-jet identities for
these fixed filters, including constant and polynomial coefficients.
It does not exclude general transcendental coefficients, merely locally
bounded coefficients, or moving filters. Nor can complete symmetry be
restored by naively adding $\Phi_N(-z)$: as real $u\to\infty$,
\[
 \Phi_N(-u)=-3\pi\Big(\sum_{n=1}^Nn^2\Big)e^{-5u}+O_N(e^{-9u}),
\]
so multiplication by $e^{tu^2}$ destroys integrability at every
positive time. Complete theta evenness is an infinite-index identity.
A vertical contour correction alone does not restore this raw-index
symmetry. Smooth compact cutoffs still require matched endpoint
derivatives and a new relative signed estimate.

\subsection{Complete-source Gaussian transport and a 14-cell pilot}

There is an exact alternative without sharp source endpoints. For
fixed $\delta\ge0$, multiply the complete source on the shifted
contour by $e^{-\delta z^2}$, $z=u+ia$. With the same frozen $a,b$
at both times, its retained observation and complement are
\begin{equation}
 G_j(t,x)=C_j(t-\delta,x),\qquad
 R_j(t,x)=C_j(t,x)-C_j(t-\delta,x).
 \label{eq:complex-gaussian-transport}
\end{equation}
This follows directly from $e^{tz^2}e^{-\delta z^2}=e^{(t-\delta)z^2}$
and the complete contour identity. The multiplier preserves the
symmetry on the imaginary segment and introduces no raw-index or upper
cutoff boundary. It is generally complex on the horizontal contour:
\[
 e^{-\delta(u+ia)^2}
     =e^{\delta a^2-\delta u^2}e^{-2i\delta au}.
\]
Accordingly \eqref{eq:paired-split-slack}, which assumes real
multipliers, is not applied unchanged. Formula
\eqref{eq:complex-gaussian-transport} supplies the physical products
directly. Numerical truncation errors still require full payment.

For the parameters \eqref{eq:filter-pilot-parameters}, fix
$\delta=1/16$ before testing. The two current-time branches in
\eqref{eq:pilot-stationary-endpoints} and the seven closed time slabs
$[r/16,(r+1)/16]$, $1\le r\le7$, give $14$ predeclared cases.
All satisfy $h_jG_j>0$, $h_jR_j>0$, and
$h_j(R_j+G_j/2)>0$. Conservative normalized margins are
\begin{center}
\begin{tabular}{@{}crrr@{}}
\toprule
$j$ & $Z_j\widehat G_j$ & $Z_j\widehat R_j$
    & $Z_j(\widehat R_j+\widehat G_j/2)$\\
\midrule
$0$ & $>8300$ & $>54$ & $>4200$\\
$1$ & $>4200$ & $>7$ & $>2100$\\
\bottomrule
\end{tabular}
\end{center}
Thus this second finite pilot has favorable retained and complementary
products, with the same prescribed zero loss. Stationarity is imposed
only at $(t,x)$, never at $(t-\delta,x)$.

The proof is exact rational extraction from the stored full-source
certificate and requires no new quadrature. Previous and current slabs
have the same spatial center and time displacements. Center coefficient
differences are formed before transport. The complete spatial and
temporal payments are $P_{\rm prior}$ for $G$,
$P_{\rm current}+P_{\rm prior}$ for $R$, and
$P_{\rm current}+P_{\rm prior}/2$ for $R+G/2$.
All full center arrays already contain the complete index and integration
tails. Intersecting with the current stationary identity
$C_j=-h_{j+2}-\mu h_j$ preserves that constraint, and products of full
integrated observations retain every cross term. The receipt in
Appendix~\ref{app:replay} records all $14$ cases. It proves the stated
finite comparison on $1/16\le t\le1/2$, not persistence of the
correlation near an unknown collision or suitability of a fixed delay
as $t\downarrow0$.

The distinct real Gaussian multiplier
$(u^2-b^2)e^{-\delta u^2}$ also retains the complete source. Denote
its retained residual by $G^{\rm real}_{\delta,j}$. It satisfies
\begin{equation}
 G^{\rm real}_{\delta,j}(t,x)
     =e^{a^2\delta}C_j(t-\delta,x+2a\delta).
 \label{eq:real-gaussian-transport}
\end{equation}
Indeed
$tz^2-\delta u^2=(t-\delta)z^2+2i\delta au-\delta a^2$.
After including the physical prefactor $e^{-ax}$, the corresponding
filtered jet is $e^{a^2\delta}h_k(t-\delta,x+2a\delta)$.
Applying $-\partial_x^2-2a\partial_x-\mu$ proves
\eqref{eq:real-gaussian-transport}.
Both identities converge for all real times; $0\le\delta\le t$
is imposed when using a theorem restricted to nonnegative time.
The real mask obeys the hypotheses of the paired multiplier formula,
but no numerical certificate for this additional mask is asserted.
These are standard Fourier/heat identities specialized to the complete
theta source, not consequences of an unproved positivity property.

Their exact complementary remainders expose the next analytic obligation.
Since the frozen observations satisfy $\partial_tC_j=-C_{j+2}$,
\begin{align}
 R_j(t,x)&=-\int_0^\delta C_{j+2}(t-s,x)\dd s,
       \label{eq:complex-gaussian-remainder}\\
 C_j(t,x)-G^{\rm real}_{\delta,j}(t,x)
   &=-\int_0^\delta e^{a^2s}
       [C_{j+2}+2aC_{j+1}+a^2C_j](t-s,x+2as)\dd s.
       \label{eq:real-gaussian-remainder}
\end{align}
The earlier or displaced point need not be stationary.

\subsection{Complex delay polarization and moving-parameter identities}
\label{subsec:complex-delay-polarization}

The Gaussian comparison uses a complex multiplier on the shifted contour.
Its paired formula must therefore conjugate the multiplier in the second
Hermitian factor.  To state this without identifying source coordinates,
use $\Omega_0=[0,\infty)\times\mathbb N$.  For $\alpha=(u,n)$ set
\[
 z_\alpha=u+ia,\quad q_\alpha=iz_\alpha,\quad
 D_\alpha=u^2-b^2,\quad
 W_\alpha=e^{tz_\alpha^2+ixu}\phi_n(z_\alpha),
\]
and put
$\mathcal F_j(k)=e^{-ax}\Re\int_{\Omega_0}W_\alpha q_\alpha^j
k_\alpha\,d\nu(\alpha)$.  The complete-source contour identity gives
$h_j=\mathcal F_j(1)$, $C_j=\mathcal F_j(D)$.
For $G_j=C_j(t-\delta,x)$, the multipliers for $P_j=h_jG_j$ and
$Q_j=h_j(C_j-G_j/2)$ are respectively
\[
 k^P_\alpha=D_\alpha e^{-\delta z_\alpha^2},\qquad
 k^Q_\alpha=D_\alpha(1-e^{-\delta z_\alpha^2}/2).
\]
The complete source is shifted before its arithmetic terms are separated.
In particular no individual raw theta term is assumed even.

\begin{lemma}[Complex polarization and full pair-tail payment]
\label{lem:complex-delay-polarization}
For the above multipliers, with independent $\alpha,\beta$,
\begin{align}
 h_j\mathcal F_j(k)=\frac{e^{-2ax}}4\Re\iint_{\Omega_0^2}
 \big[&(k_\alpha+\overline{k_\beta})W_\alpha\overline{W_\beta}
                       (q_\alpha\overline{q_\beta})^j\notag\\
     &+(k_\alpha+k_\beta)W_\alpha W_\beta(q_\alpha q_\beta)^j
 \big]d\nu(\alpha)d\nu(\beta).
 \label{eq:complex-delay-polarization}
\end{align}
For $a=\pi/16$, $b=1/4$, $0\le t\le1/2$ and
$0\le\delta\le\min(t,1/4)$, put $M=e^{a^2/4}$.  If $A$ is
a retained set, let $M_{A,k}=\int_A|W||q|^k$ and
$T_k=\int_{A^c}|W||q|^k$.  The absolute multiplier masses satisfy
\[
 B_A^P\le M(M_{A,j+2}+b^2M_{A,j}),\qquad
 B_A^Q\le(1+M/2)(M_{A,j+2}+b^2M_{A,j}),
\]
and the same bounds hold on $A^c$ with $M_{A,k}$ replaced by $T_k$.
Truncating both factors of either product has error at most
\begin{equation}
 e^{-2ax}\{M_{A,j}B_{A^c}+T_jB_A+T_jB_{A^c}\}.
 \label{eq:complex-delay-tail}
\end{equation}
\end{lemma}
\begin{proof}
Use $\Re U\Re V=\frac12\Re(U\overline V+UV)$ and then average
under interchange of the two source coordinates.  This gives the two
sectors and the conjugated multiplier in
\eqref{eq:complex-delay-polarization}.  Absolute convergence follows,
with $c_a=\cos(4a)$ and $\rho_0=16e^{-3\pi c_a}<1$, from
\[
 \sum_n|W_{(u,n)}||q_{(u,n)}|^k
 \le\frac{2\pi^2+3\pi}{1-\rho_0}(u+a)^k
           e^{u^2/2+9u-\pi c_a e^{4u}},
\]
which has every polynomial moment.  The inequality
$|e^{-\delta(u+ia)^2}|\le M$ and
$|u^2-b^2|\le|q|^2+b^2$ prove the multiplier bounds.  Expanding a
product into its retained and omitted factors proves
\eqref{eq:complex-delay-tail}, including both mixed sectors and the
omitted--omitted sector.  The complete index and integration bounds of
Appendix~\ref{app:payments} supply the $T_k$ on their declared
rectangles.  Quadrature and parameter-transport errors remain additional.
\end{proof}

The real horizontal Gaussian $e^{-\delta u^2}$ is a different mask;
its time-and-position transport identity above must not replace the
complex mask in this lemma.  At a fixed observation point the latter
has the exact time identity
\[
 G_j=C_j(t-\delta,x),\qquad
 C_j(t,x)-G_j=-\int_0^\delta C_{j+2}(t-s,x)\,ds.
\]
Earlier stationarity is never imposed.  If $\delta$ depends on $t,x$,
its derivatives also enter transport.  With superscript $-$ denoting
evaluation at $(t-\delta(t,x),x)$ and fixed $a,b$,
\begin{align}
 G_{j,t}&=-(1-\delta_t)C_{j+2}^-,&
 G_{j,x}&=C_{j+1}^-+\delta_x C_{j+2}^-,\notag\\
 (\partial_t+\partial_x^2)G_j
 &=(\delta_t+\delta_{xx})C_{j+2}^-
       +2\delta_xC_{j+3}^-+\delta_x^2C_{j+4}^-.
 \label{eq:general-moving-delay}
\end{align}
These identities follow by the chain rule and $\partial_tC_k=-C_{k+2}$.
All needed derivatives are justified by the complete-source moment
majorant, not by differentiating a four-jet asymptotic.
On a regular current branch $x=\xi(t)$,
\[
 \xi'=h_{j+3}/h_{j+2},\qquad
 \dot h_j=-h_{j+2},\qquad
 \dot G_j=\xi'C_{j+1}^--(1-\delta_t-\xi'\delta_x)C_{j+2}^-.
\]
Thus $\dot P_j=-h_{j+2}G_j+h_j\dot G_j$ and
\[
 \dot Q_j=-h_{j+2}(C_j-G_j/2)
       +h_j[-C_{j+2}+\xi'C_{j+1}-\dot G_j/2].
\]
The denominator-free tangent operator
$\mathcal T_j=h_{j+2}\partial_t+h_{j+3}\partial_x$ satisfies
\[
 \mathcal T_jh_j=-h_{j+2}^2,\quad
 \mathcal T_jG_j=-[h_{j+2}(1-\delta_t)-h_{j+3}\delta_x]C_{j+2}^-
                              +h_{j+3}C_{j+1}^-.
\]
It may vanish at a singular point.  None of these transport identities
alone proves continuation or positivity through branch creation.

\subsection{Why every fixed proportional delay fails its exact retained sign}
\label{subsec:proportional-delay-obstruction}

\begin{theorem}[Proportional-delay phase obstruction]
\label{thm:proportional-delay-obstruction}
Fix $t_0\in(0,1/2]$, $0<\rho<1$, and finite real centering
parameters $a,b$.  Put $\delta=\rho t_0$, $\tau=(1-\rho)t_0$.
For each $j=0,1$, the genuine current set $h_{j+1}(t_0,x)=0$
contains sequences $x\to+\infty$ on which
$h_j(t_0,x)C_j(\tau,x)<0$, and other sequences with the opposite
strict sign.  Nevertheless, for every fixed $\lambda\in[0,1]$,
\[
 h_j(t_0,x)\{C_j(t_0,x)-(1-\lambda)C_j(\tau,x)\}>0
\]
at every sufficiently high current stationary point.  Both full
Laguerre forms are also positive there.
\end{theorem}
\begin{proof}
Use the analytic carrier $Q_v$ in \eqref{eq:analytic-model} at each
fixed positive time $v=t_0,\tau$, and put
$A_v=2|Q_v|$, $\phi_v=\arg Q_v$, $L=\log(x/(4\pi))$,
$\omega=L/4$.  The four-jet theorem applies at both times for all
sufficiently large $x$, giving
\[
 h_k(v,x)=A_v\omega^k\{\cos(\phi_v+k\pi/2)+O(L^{-1})\},
 \quad 0\le k\le3,\qquad \phi_{t_0}'=\omega+O(1)>0.
\]
The exact carrier quotient and
$\alpha(s)=L/2+i\pi/4+O(x^{-1})$, $s=(1+ix)/2$, give
\begin{equation}
 D:=\phi_{t_0}-\phi_\tau=\frac{\delta\pi L}{16}+o(1),\qquad
 r_\delta:=\frac{A_\tau}{A_{t_0}}
  =\exp\{-\delta L^2/16+\delta\pi^2/64+o(1)\}.
 \label{eq:fixed-delay-phase-amplitude}
\end{equation}
Indeed the two zeta arguments have real parts tending to infinity, so
their Dirichlet series and their analytic logarithms tend to 1 and 0,
respectively.  Continuous carrier phases may be chosen consistently.

At current stationarity let $\psi=\phi_{t_0}+j\pi/2$; then
$\sin\psi=O(L^{-1})$.  The lower-order drift and mass terms in the
residual contribute only $O(L^{-1})$ after normalization.  Hence
\[
 \frac{h_j(t_0,x)C_j(\tau,x)}
      {A_{t_0}A_\tau\omega^{2j+2}}=\cos D+O(L^{-1}),\qquad
 \frac{h_j(t_0,x)C_j(t_0,x)}
      {A_{t_0}^2\omega^{2j+2}}=1+O(L^{-1}).
\]
The normalized absorption expression tends to 1 because
$r_\delta\to0$.

To obtain actual stationary points, take each sufficiently high phase
window $|\psi-k\pi|\le\pi/6$.  The signs of $h_{j+1}$ at its
endpoints are opposite, and $h_{j+2}$ has one strict sign throughout,
so the window contains exactly one current stationary point.  On the
logarithmic bands
\[
 \left|L-\frac{16(2n+1)}\delta\right|\le\frac8{3\delta},
\]
the phase $D$ is within $\pi/6+o(1)$ of $(2n+1)\pi$ and
$\cos D\le-1/2$ eventually.  Their lengths in $x$ grow exponentially,
whereas the phase windows have width of order $1/L$, so each sufficiently
high band contains complete windows of both levels.  Bands centered at
$L=32n/\delta$ give the positive sign.  All these points lie in the
already proved outer-sign region.
\end{proof}

The obstruction also excludes a height-independent rule $\delta(t)$
whenever $0<\delta(t_0)<t_0$ at one positive time.  It concerns the
exact retained-sign requirement, not the bounded-loss program: on these
proved nonzero fixed-time branches,
\[
 P_j/E=r_\delta\omega^2\{\cos D+O(L^{-1})\}\longrightarrow0.
\]
The sign alternates while its energy-relative magnitude vanishes.
No first counterexample height is evaluated, and this fixed-time argument
does not settle a moving inner cusp.

For the particular half-time rule $\delta=t/2$, a complete-source pilot
in $(101,102)$ on $h_1=0$ and $(103,104)$ on $h_2=0$ certifies
$P_j,Q_j>0$ in all 16 branch/slab cases over $0\le t\le1/2$.
In the constant branch normalizations its lower bounds for $(P_j,Q_j)$
are $(8108,3999)$ and $(4100,2025)$.  The complement $h_j(C_j-G_j)$
is undecided on the first two time slabs of each branch, and equals zero
at $t=0$; these four undecided intervals are not negative signs.
The exact earlier-time map on slab $k$ is
$d^-=d/2+\epsilon_k$, with $\epsilon_k=-1/64$ for even $k$ and
$1/64$ for odd $k$.  Substitution in the common Taylor polynomial,
before interval ranging and with all original payments retained, proves
the pilot.  Its record is
\path{numerics/PROPORTIONAL_GAUSSIAN_SOURCE_FILTER_PILOT_20261014.json}.
This bounded success is consistent with the asymptotic obstruction.

\subsection{A height-dependent delay with a uniform outer comparison}
\label{subsec:phase-controlled-delay}

Keep $a=\pi/16$, $b=1/4$ fixed, and prescribe
\begin{equation}
 \ell(x)=\log(1+x^2),\quad B=1+t\ell(x),\quad
 \delta(t,x)=\frac{t}{2B},\quad\tau=t-\delta,
 \quad G_j=C_j(\tau,x).
 \label{eq:smooth-delay-rule}
\end{equation}
For $t\ge0$, $0\le\delta\le t/2$ and $t/2\le\tau\le t$.

\begin{theorem}[Phase-controlled outer comparison]
\label{thm:phase-controlled-delay}
There is a finite unevaluated $X$ such that, uniformly for
$|x|\ge X$, $0<t\le1/2$, and
$\theta=t\log(|x|/(4\pi))\ge13/4$, the rule
\eqref{eq:smooth-delay-rule} satisfies at current stationarity
\[
 P_j\ge\frac18A_tA_\tau\omega^{2j+2}>0,\qquad
 Q_j\ge\frac12A_t^2\omega^{2j+2}>0,\qquad j=0,1,
\]
where $A_v=2|Q_v|$ is the analytic carrier amplitude at time $v$ and
$\omega=\frac14\log(|x|/(4\pi))$.
\end{theorem}
\begin{proof}
Take $x>4\pi$ and write $L=\log(x/(4\pi))$, $d=\delta L$.
Since $\ell(x)>2L$,
\[
 0<d<\frac{\theta}{2+4\theta}<\frac14,\qquad
 \tau L>f(\theta):=\theta-\frac{\theta}{2+4\theta}.
\]
The increasing function $f$ satisfies $f(13/4)=91/30=3+1/30$.
Thus the uniform four-jet asymptotic is available at both times with
a fixed margin, even when those times decrease with height.  Write
$s=(1+ix)/2$, $w_v=s+v\alpha(s)/2$.  The full relative phase is
\begin{equation}
 D:=\phi_t-\phi_\tau
 =\Im\left\{\frac\delta4\alpha(s)^2+
       \int_\tau^t\frac{\alpha(s)}2
                \frac{\zeta'}{\zeta}(w_v)\,dv\right\}.
 \label{eq:smooth-delay-full-phase}
\end{equation}
The analytic prime-power logarithm is valid throughout this path:
$\Re w_v\ge5/4$ uniformly once $X$ is large.  In that half-plane,
\[
 \left|\frac{\zeta'}\zeta(w)\right|
 \le\sum_{n\ge2}(\log n)n^{-5/4}<18.
\]
For the last inequality bound the $n=2,3$ terms by 1 each and the
decreasing remainder by
$\int_3^\infty(\log u)u^{-5/4}du<16$.
Since $\alpha=L/2+i\pi/4+O(x^{-1})$ and $d<1/4$, the gamma and
arithmetic parts of \eqref{eq:smooth-delay-full-phase} are bounded by
$\pi/64+o(1)$ and $18/16+o(1)$, respectively.  Increase $X$ until
\[
 |D|<5/4,\qquad \cos D\ge1-D^2/2\ge7/32.
\]
At current stationarity, $\psi=\phi_t+j\pi/2$ has
$\sin\psi=O(L^{-1})$, so
\[
 \frac{P_j}{A_tA_\tau\omega^{2j+2}}
 =\cos\psi\cos(\psi-D)+O(L^{-1})=\cos D+O(L^{-1}).
\]
This proves the bound $1/8$ after a further uniform increase of $X$.

For absorption, the current jets give
$h_jC_j/(A_t^2\omega^{2j+2})=1+O(L^{-1})$.
The delayed/current amplitude ratio is exactly
\[
 r_\delta=\frac{A_\tau}{A_t}
  =e^{-\delta\Re(\alpha^2)/4}
       \left|\frac{\zeta(w_\tau)}{\zeta(w_t)}\right|.
\]
The second factor is at most $\zeta(5/4)^2\le25$, by the
absolutely convergent series for zeta and its reciprocal.  Uniformly
for $\theta\ge13/4$, the fact that $\ell/L\to2$ implies $d\ge1/5$
eventually (it suffices that $\ell/L\le21/10$).  Therefore
\[
 r_\delta\le25e^{-L/80+O(L^{-1})},\qquad
 \frac{Q_j}{A_t^2\omega^{2j+2}}
 =1+O(L^{-1})-\frac{r_\delta}{2}\{\cos D+O(L^{-1})\}
 \ge\frac12
\]
eventually.  For $x=-y<0$, the delay is unchanged and
$h_k(t,-y)=(-1)^kh_k(t,y)$.  The drift in the delayed residual
changes sign under reflection but remains in the same lower-order
$O(L^{-1})$ payment.  The leading margins are consequently unchanged.
\end{proof}

This theorem proves a delayed comparison within the known outer-sign
region.  It does not improve the full-sign threshold.  Using only the
historical condition $\tau L>3$ in the same proof gives, for every
fixed $\varepsilon>0$, the comparison on
$\theta\ge(11+\sqrt{217})/8+\varepsilon$, with an unevaluated
$\varepsilon$-dependent height.  No claim is made at the limiting
value, nor is that historical threshold asserted optimal for the filter.

For completeness, the exact moving derivatives in
\eqref{eq:general-moving-delay} are supplied by
\[
 \delta_t=\frac1{2B^2},\quad
 \delta_x=-\frac{t^2\ell'}{2B^2},\quad
 \delta_{xx}=-\frac{t^2\ell''}{2B^2}
                  +\frac{t^3(\ell')^2}{B^3},
\]
where $\ell'=2x/(1+x^2)$ and
$\ell''=2(1-x^2)/(1+x^2)^2$.  The pointwise theorem requires no
regular stationary branch; its transport formulas do not themselves
handle branch creation.  The finite all-earlier-time result
\eqref{eq:height112-causal} includes this rule and the half-time rule
on the two new positive-height branches.

The exact checkers
\path{numerics/check_proportional_delay_identities.py} and
\path{numerics/check_phase_controlled_delay_identities.py}
verify 20 and 21 finite algebraic statements, respectively.  The affine
transport audit performs 1200 rational point comparisons.  Their
receipts and Reviews 18--19 distinguish these checks from the analytic
phase proof and the complete-source interval calculations.  No check
evaluates the asymptotic height or proves an unbounded inner-cusp estimate.

\subsection{The remaining cross-time estimate}

One optional sufficient target is to prescribe a delay rule $0<\delta(t,x)\le t$,
a fraction $0\le\lambda<1$, and losses locally bounded through a
candidate collision, and prove on the \emph{current} set $h_{j+1}(t,x)=0$
that
\begin{align}
 h_j(t,x)C_j(t-\delta,x)&\ge0,\notag\\
 h_j(t,x)\{C_j(t,x)-(1-\lambda)C_j(t-\delta,x)\}
        &\ge-K_j(t,x)E(t,x),\qquad j=0,1.
 \label{eq:cross-time-target}
\end{align}
Lemma~\ref{lem:direct-absorption} would then supply the required
stationary loss. This is a theta-specific correlation statement: neither
backward heat, complete-source symmetry, nor the transport identities
prove its sign. The finite comparisons establish it only on their declared
cells; Theorem~\ref{thm:phase-controlled-delay} gives the separate
outer comparison. Any adaptive parameters remain prescribed independently of the
unknown energy; differentiating an adaptive filter additionally requires
all chain-rule and moving-boundary terms. The paired source formulas and
the geometric formulation identify what must be controlled, while
\eqref{eq:cross-time-target} remains an open all-height analytic estimate.
The combined signed observation can be bounded directly without
requiring both separate products to have favorable signs;
Section~\ref{pg:research-program} develops that less restrictive target.

\subsection{Polynomial cutoff losses as a quantitative target}
\label{subsec:polynomial-cutoff-transfer}

The useful role of logarithmic slack is to permit deterioration with
source resolution while retaining a collision-excluding estimate.  The
following transfer is independent of any unproved theta sign.

\begin{lemma}[Polynomial cutoff transfer]
\label{lem:polynomial-cutoff-transfer}
On a declared stationary domain, suppose independently established
$A,B\ge0$, $c>0$, a fixed finite $p\ge0$, and an integer $N_0\ge1$
give, for every integer $N\ge N_0$,
\begin{equation}
 \Lag_j\ge-A N^p E-B e^{-cN^2}.
 \label{eq:polynomial-cutoff-hypothesis}
\end{equation}
Assume $A,B,N_0,c^{-1}$ are locally bounded through candidate events.
For every fixed $M>0$ there is an independently prescribed, locally
bounded $C$ such that
\begin{equation}
 \Lag_j\ge-C E[1+\log_+(M/E)]^{p/2},
 \label{eq:polynomial-cutoff-conclusion}
\end{equation}
with value zero on the right at $E=0$.
\end{lemma}
\begin{proof}
For $E>0$ set $\ell_E=1+\log_+(M/E)$,
$d_B=\log_+(B/M)$ if $B>0$ and $d_B=0$ otherwise.  Take
\[
 N=\left\lceil\max\left\{N_0,
       \sqrt{c^{-1}\log_+(B/E)}\right\}\right\rceil,
\]
where the logarithm is defined as zero for $B=0$.  Then
$B e^{-cN^2}\le E$ and
$\log_+(B/E)\le d_B+\log_+(M/E)\le(1+d_B)\ell_E$.
Consequently $N\le D\sqrt{\ell_E}$ for
\[
 D=N_0+1+c^{-1/2}\sqrt{1+d_B},\qquad C=AD^p+1.
\]
This proves \eqref{eq:polynomial-cutoff-conclusion} and its local
uniformity.  At $E=0$, let $N\to\infty$ in the hypothesis.
\end{proof}

The same conclusion follows if a retained observation
$\Lag_j^{(N)}$ obeys the corresponding inequality and its difference
from $\Lag_j$ has a separately paid Gaussian tail.  This is an
all-$N$ family with fixed coefficient rules: proving one cutoff or
choosing coefficients from an unknown energy floor does not suffice.
The final optimization may depend on $E$ because all members of the
family have already been proved with the same prescribed constants.

For the common-carrier source, use the exact staircase cutoff
\[
 \Omega_N=\{(v,n):0\le v\le\tfrac12\log N,\quad n\le e^{2v}\}.
\]
The complete-source majorant in \eqref{eq:shared-majorant} implies
Gaussian omitted moments $O(e^{-cN^2})$ at bounded time, after reducing
$c>0$ to absorb polynomial and logarithmic factors.  The paired error
\eqref{eq:pair-tail-payment} pays both mixed sectors and the
omitted--omitted sector.  Thus a sufficient new arithmetic target on an
unresolved domain is
\begin{equation}
 \Lag_j^{(N)}\ge-A_j(t,x)N^pE-B_j(t,x)e^{-cN^2}
 \quad\text{on the full set }h_{j+1}=0,
 \quad N\ge N_0,\quad j=0,1.
 \label{eq:inner-signed-cutoff-target}
\end{equation}
All retained jet cross terms must remain, and both $E$ and stationarity
refer to the full physical jets.  The coefficients must be derived
independently of $E$ and locally bounded through candidate events.
A formulation with retained energy and a paid stationary defect is also
possible: the tail implies $E_N\le2E+O(e^{-2cN^2})$ and
$|h_{j+1}^{(N)}|=O(e^{-cN^2})$ at full stationarity.
None of these tail estimates proves the signed target.

\subsection{A cumulative signed moment and its stronger sufficient bound}
\label{subsec:cumulative-signed-moment}

An alternative exposes the signed correlation directly on the original
real source.  For $j=0,1$, put
\[
 a_j(u)=e^{tu^2}\Phi(u)u^j\cos(xu+j\pi/2),\quad
 W=h_j=\int_0^\infty a_j(u)\,du,\quad
 F_j(r)=\int_0^r a_j(u)\,du.
\]
At full stationarity, $\Lag_j=W\int_0^\infty u^2a_j(u)\,du$.
An integration by parts, with the full tail retained, gives
\begin{equation}
 \Lag_j=\int_0^R2r\{W^2-WF_j(r)\}\,dr
       +W\int_R^\infty(u^2-R^2)a_j(u)\,du.
 \label{eq:cumulative-signed-identity}
\end{equation}
Indeed integrating $u^2\,dF_j(u)$ on $[0,R]$ gives
$R^2F_j(R)-\int_0^R2rF_j(r)dr$, and
$W-F_j(R)=\int_R^\infty a_j$ supplies the displayed tail.

Suppose, as an additional unproved source estimate, that
\begin{equation}
 WF_j(r)\le W^2+A_jE\qquad(r\ge0)
 \label{eq:cumulative-overshoot-target}
\end{equation}
with a prescribed locally bounded $A_j\ge0$.  For $0\le t\le1/2$,
$0\le k\le3$, and $R\ge1$, the complete real-source envelope yields
\begin{equation}
 \int_R^\infty u^ke^{tu^2}\Phi(u)\,du
 \le2\pi e^{-4R}e^{-(\pi/2)e^{4R}}
 \le2\pi e^{-(\pi/2)e^{4R}}.
 \label{eq:cumulative-real-tail}
\end{equation}
To verify this, use $\Phi(u)<4\pi^2e^{9u-\pi e^{4u}}$ and, for
$u\ge1$, $k\log u+tu^2+9u\le u^2/2+12u\le(\pi/2)e^{4u}$.
Integrating the remaining exponential using its derivative gives the
first bound.  Since $|W|\le\sqrt E$ for both levels,
\eqref{eq:cumulative-signed-identity} implies
\[
 \Lag_j\ge-A_jR^2E-2\pi\sqrt E\,e^{-(\pi/2)e^{4R}}.
\]
At $E=0$ the stationary identity gives $\Lag_j=0$ directly. For $E>0$, set $\ell_E=1+\log_+(M/E)$,
$R=1+\tfrac14\log\ell_E$, and $c_0=(\pi/2)e^4$.
Then $R^2\le25\ell_E/16$ and
$\sqrt E e^{-c_0\ell_E}\le e^{-c_0}E/\sqrt M$, both for
$E\le M$ and for $E\ge M$.  The conditional consequence is
\[
 \Lag_j\ge-\left(\frac{25}{16}A_j+
                 \frac{2\pi e^{-c_0}}{\sqrt M}\right)
                  E[1+\log_+(M/E)].
\]
Keeping $R^2=(1+\frac14\log\ell_E)^2$ gives the sharper squared
iterated-log factor.  This is a deduction from
\eqref{eq:cumulative-overshoot-target}, not a proof of that target.
An estimate on the integrated expression in
\eqref{eq:cumulative-signed-identity} would demand less than pointwise
control for every $r$.

Positive source weights alone cannot establish the needed alignment.
For the positive two-frequency heat family
$H_*(t,x)=e^t\cos x+e^{-3T}e^{4t}\cos(2x)$, there is a double
collision at $(T,\pi)$.  At $t=T-s$, $x=\pi$,
$W=-3e^Ts+O(s^2)$, $E=9e^{2T}s^2+O(s^3)$, and
$\Lag_0=-9e^{2T}s+O(s^2)$.  For a source cutoff $1<r<2$,
$F_0(r)=-e^{T-s}$ and $WF_0(r)=3e^{2T}s+O(s^2)$, contradicting
\eqref{eq:cumulative-overshoot-target} for bounded $A_0$.
This atomic-source control is outside the theta class.  It shows why
an arithmetic relation among the actual theta contributions is needed;
evenness, rapid decay, or an absolute source-tail payment does not supply
an energy-relative signed estimate.

% Prepared for Edward Baker, 15 October 2026.
% Substantial LLM assistance: gpt-6-astra, reasoning effort ultra,
% as verified and supplied by the parent session for this turn.
% eq:wpm, eq:RS-input, eq:shared-majorant, fc:predecessor.

% Prepared for Edward Baker, 15 October 2026.
% Substantial LLM assistance: gpt-6-astra, reasoning effort ultra.
% Integrates research Notes 38--44; internal review is not formal verification.
\section{Poisson localization, sharper remainders, and subpower losses}
\label{upd:section}

The complete correlation identities admit a genuinely two-coordinate
arithmetic transformation.  Its separate sectors, however, lose the
decay needed by positive-time heat integration.  We prove this obstruction,
give two admissible localizations, and quantify why their absolute budgets
do not yet supply the stationary energy estimate.  We also sharpen the
existing weighted remainder and extend the conditional cutoff transfer.
None of these results enlarges the proved sign domain.

\subsection{The complete parity--Poisson identity}

Write $\Theta(y)=\sum_{n\in\mathbb Z}e^{-\pi n^2y}$ and
$P(u)=e^u\Theta(e^{4u})$. For independent $s,r\in\R$, set
\[
 C=2\cosh(4r),\quad \tau=\tanh(4r),\quad f=C^{-1/2},\quad
 \Pi(s,r)=P(s+r)P(s-r),
\]
and let
\[
 \mathscr D=(\partial_s^2-\partial_r^2)^2
                    -8(\partial_s^2+\partial_r^2)+16.
\]
\begin{proposition}[Arithmetic transformation with both coordinates retained]
\label{upd:parity-poisson}
The exact identities
\begin{align}
 \Pi(s,r)&=\frac1{\sqrt C}\sum_{p,k\in\mathbb Z}
 e^{-\pi(e^{4s}p^2+e^{-4s}k^2)/C}
                  (-1)^{pk}\cos(\pi pk\tau),\label{upd:parity-series}\\
 \Phi(s+r)\Phi(s-r)&=\frac1{4096}\mathscr D\Pi(s,r),
 \label{upd:parity-source}\\
 \mathscr K_m(t,s)&=\frac{e^{2ts^2}}{4096}
       \int_{\R}r^{2m}e^{2tr^2}\mathscr D\Pi(s,r)\dd r
 \label{upd:parity-correlation}
\end{align}
hold.  The series and all its finite coordinate derivatives converge
normally on compact subsets of $\R^2$.
\end{proposition}
\begin{proof}
Put $u=s+r,v=s-r,y_+=e^{4u},y_-=e^{4v}$ and change the two
theta indices to $p=m+n,q=m-n$.  This is a bijection onto the parity
lattice $p\equiv q\pmod2$.  With $A=y_++y_-$ and $B=y_+-y_-$,
\[
 y_+m^2+y_-n^2=\frac A4(q+(B/A)p)^2+\frac{y_+y_-}{A}p^2.
\]
Write $q=p+2\ell$ and apply Poisson summation in $\ell$:
\[
 \sum_{\ell\in\mathbb Z}e^{-\pi A(\ell+\alpha p)^2}
 =A^{-1/2}\sum_{k\in\mathbb Z}e^{-\pi k^2/A}e^{2\pi i k\alpha p},
 \qquad \alpha=\frac{1+\tau}{2}.
\]
Since $A=e^{4s}C$ and $y_+y_-=e^{8s}$, taking real parts gives
\eqref{upd:parity-series}.  On a compact coordinate set the Gaussian
coefficients are bounded below; differentiation introduces only fixed
polynomials in the indices.  This proves normal convergence.
Finally $\Phi=(P''-P)/16$ and
$\partial_u=(\partial_s+\partial_r)/2$,
$\partial_v=(\partial_s-\partial_r)/2$ give
\eqref{upd:parity-source}.  The heat factor remains exactly
$e^{2t(s^2+r^2)}$, proving \eqref{upd:parity-correlation}.
\end{proof}

\subsection{Why undamped sectors and modes cannot be integrated separately}

Split the axes, counting $(0,0)$ once, from the mixed modes:
\begin{align}
 \Pi_{\rm ax}&=f\{\Theta(e^{4s}/C)+\Theta(e^{-4s}/C)-1\}\notag\\
 &=e^{-2s}\Theta(Ce^{-4s})+e^{2s}\Theta(Ce^{4s})-f,
 \qquad \Pi_{\rm mix}=\Pi-\Pi_{\rm ax}.\label{upd:axis-definition}
\end{align}
\begin{proposition}[Positive-time nonintegrability]
\label{upd:sector-obstruction}
Uniformly for $s$ in each bounded interval, as $|r|\to\infty$,
\[
 \frac{\mathscr D\Pi_{\rm ax}}{4096}
       =\frac98e^{-10|r|}(1+o(1)),\qquad
 \frac{\mathscr D\Pi_{\rm mix}}{4096}
       =-\frac98e^{-10|r|}(1+o(1)).
\]
Thus for every fixed $t>0,s\in\R,m\ge0$, their separate
$r^{2m}e^{2t(s^2+r^2)}$ integrals diverge to $+\infty$ and
$-\infty$, respectively.  Every individual undamped mode in
\eqref{upd:parity-series}, and every nonempty centered finite
rectangular packet of modes, also has divergent absolute positive-time
heat integral.
\end{proposition}
\begin{proof}
For $|s|\le S$, the axis formula is
$2\cosh(2s)-f+O_S(e^{-\pi Ce^{-4S}})$; its derivatives through order
four have the same bound with an additional factor $O_S(C^4)$.
This follows by differentiating the nonconstant theta terms and
summing their Gaussian majorants.  The first term is annihilated by
$\mathscr D$.  Put $q=\operatorname{sech}^2(4r)$. Direct differentiation
gives
\[
 (\partial_r^2-4)f=-12qf,\quad
 (\partial_r^2-4)(qf)=(96-140q)qf,\quad
 -\mathscr Df=48fq(24-35q).
\]
Consequently $-\mathscr Df/4096=(9/8)e^{-10|r|}
(O(e^{-18|r|}))$.  In contrast, for $r\ge S$, the complete source obeys
\[
 0<\Phi(s+r)\Phi(s-r)
 \le16\pi^4e^{18r}\exp\{-2\pi e^{4r}\cosh(4s)\}.
\]
Subtracting proves the mixed-sector asymptotic and both divergences.

For the mode
$A_{p,k}=f e^{-\pi(e^{4s}p^2+e^{-4s}k^2)/C}
(-1)^{pk}\cos(\pi pk\tau)$, put $z=e^{-4r}$ and
$B=e^{4s}p^2+e^{-4s}k^2$.  At fixed indices, uniformly on a bounded
$s$-strip and with four differentiated remainders,
\[
 A_{p,k}=z^{1/2}\left[1-\pi Bz+
            \left(\frac{\pi^2B^2}{2}-\frac12\right)z^2
            +O_{p,k,S}(z^3)\right].
\]
The symbol of $\mathscr D$ on $e^{\lambda s-ar}$ is
$\lambda^4-(2a^2+8)\lambda^2+(a^2-4)^2$.
It vanishes at $(a,\lambda)=(2,0),(6,\pm4),(10,\pm8)$, and equals
$9216$ at $(10,0)$. Hence
\begin{equation}
 \frac{\mathscr DA_{p,k}}{4096}
 =\left(\frac94\pi^2p^2k^2-\frac98\right)e^{-10r}
                  +O_{p,k,S}(e^{-14r}).\label{upd:mode-asymptotic}
\end{equation}
The leading coefficient never vanishes.  Summing over
$|p|\le P_0,|k|\le K_0$ gives coefficient
\[
 \frac98(2P_0+1)(2K_0+1)
 \left\{\frac{2\pi^2P_0(P_0+1)K_0(K_0+1)}9-1\right\}.
\]
It is negative when either cutoff is zero, and positive otherwise.
Every stated absolute heat integral therefore diverges.
\end{proof}

The opposite signs of the single axis modes and the complete axis sum
are consistent: the fixed-index expansion is not uniform in the indices
as $r$ grows.  Neither this expansion nor the compact normal convergence
permits unbounded heat integration term by term.  The same tail
obstruction persists for $r^2$, $r^4-s^2r^2$, and either of these plus
a fixed nonnegative multiple of $(1-s^2+r^2)/2$.

\begin{proposition}[Limits of finite diagonal compensation]
\label{upd:diagonal-compensation}
Adding $f$ to $\Pi_{\rm ax}$ repairs its fixed-$s$ tail but does not
give an absolutely heat-integrable source in both coordinates.
Moreover, no fixed-order subtraction of the modular Taylor profiles
described below from $\Pi_{\rm ax}+f$ repairs positive-time absolute
heat integrability.
\end{proposition}
\begin{proof}
For fixed $r$ and $s\to+\infty$, scalar Poisson summation gives
$\Pi_{\rm ax}=e^{2s}+o(1)$ with superexponentially small differentiated
remainders.  Thus $\mathscr D(\Pi_{\rm ax}+f)\to\mathscr Df(r)$,
which is nonzero on an open set ($\mathscr Df(0)=528/\sqrt2$).
The $s$ heat factor prevents absolute integration.
For a more precise joint-escape statement put $u=s+r,v=s-r$.
For fixed $K\ge0$, uniformly on compact $v$ intervals,
\[
 \Pi_{\rm ax}+f=e^{u+v}+\sum_{n=0}^{K}e^{-(4n+1)u}F_n(v)
                       +O_K(e^{-(4K+5)u}),\qquad
 F_n(v)=\frac{e^{-v}}{n!}\Theta^{(n)}(e^{-4v}).
\]
This follows by Taylor expansion of the axis theta formula, and remains
valid after four derivatives.  Here $F_0=P$ and
$\mathscr D/4096=(\partial_u^2-1)(\partial_v^2-1)/256$, which
annihilates the first two displayed terms. For $n\ge1$,
\[
 (F_n''-F_n)(0)=\frac{(-1)^n}{n!}\sum_{\ell\ne0}
   (\pi\ell^2)^{n+1}(16\pi\ell^2-24)e^{-\pi\ell^2}\ne0.
\]
Indeed applying $\partial_v^2-1$ to $e^{-v}g(e^{-4v})$ gives
$e^{-v}(24yg'(y)+16y^2g''(y))$, and the terms in the last sum
have one strict sign.  After subtracting profiles $1,\ldots,K$, the
leading differentiated remainder is
\[
 \frac{(4K+5)^2-1}{256}e^{-(4K+5)u}
                        (F_{K+1}''-F_{K+1})(v).
\]
Its coefficient is nonzero near $v=0$, so multiplication by the exact
heat $e^{t(u^2+v^2)}$ forces divergence on that tube.
\end{proof}
This rules out the specified finite subtraction class, not exact
resummations or localized compensation.

\subsection{Two admissible localized signed budgets}

Set $\Omega=s^2+r^2$, $q_0=r^2$, $q_1=r^4-s^2r^2$,
$b_E=(1-s^2+r^2)/2$, and
\[
 W_{j,\kappa}=e^{2t\Omega}(q_j+\kappa_jb_E)\cos(2xs).
\]
The nonnegative coefficient $\kappa_j(t,x)$ is prescribed independently
of the unknown energy and is locally bounded on the stated domain.
For any window $0\le\chi\le1$ for which the following sector integrals
exist, define
\[
 M_j^\chi=\frac1{4096}\iint\chi W_{j,\kappa}\mathscr D\Pi_{\rm ax},
 \quad R_j^\chi=\frac1{4096}\iint\chi W_{j,\kappa}\mathscr D\Pi_{\rm mix},
 \quad T_j^\chi=\iint(1-\chi)W_{j,\kappa}\Phi(s+r)\Phi(s-r).
\]
All double integrals here use $\dd s\dd r$. The complete correlation
and energy identities imply
\begin{equation}
 \Lag_j+\kappa_jE=M_j^\chi+R_j^\chi+T_j^\chi.
 \label{upd:localized-budget}
\end{equation}
The exterior is always the complete unsplit source.

\begin{proposition}[Compact-box and smooth-window payments]
\label{upd:localization}
For $\chi=1_{[-S,S]\times[-R,R]}$, all terms in
\eqref{upd:localized-budget} are finite. Put $U=\min(S,R)$ and
\[
 \mathsf M_k(t)=\int_{\R}|u|^kw_t(u)\dd u,\qquad
 \mathsf T_k(t,U)=\int_{|u|>U}|u|^kw_t(u)\dd u.
\]
Then, for $j=0,1$,
\begin{align}
 |T_j^\chi|\le{}&\frac12(\mathsf T_{2j+2}\mathsf M_0+
                                  \mathsf T_0\mathsf M_{2j+2})\notag\\
 &+\kappa_j\left\{\frac12\mathsf T_0\mathsf M_0+
              \frac14(\mathsf T_2\mathsf M_0+\mathsf T_0\mathsf M_2)
                    \right\}.\label{upd:box-tail}
\end{align}
For $0\le t\le1/2,U\ge1,0\le k\le4$ one may use
\begin{equation}
 \mathsf T_k(t,U)\le4\pi e^{-4U}
                              e^{-(\pi/2)e^{4U}}.\label{upd:box-tail-bound}
\end{equation}
Alternatively, for $\chi_\varepsilon=e^{-\varepsilon\Omega^2}$,
$\varepsilon>0$, both sector sums can be integrated term by term
absolutely, uniformly for $\varepsilon\ge\varepsilon_0>0$, bounded
heat ceiling, and locally bounded $\kappa_j$. With
\[
 \mu_k(t)=2^{-k-1}\sum_{\ell=0}^{k}\binom{k}{\ell}
                   \mathsf M_{2\ell}(t)\mathsf M_{2k-2\ell}(t),
\]
the complete exterior satisfies
\begin{equation}
 |T_j^{\chi_\varepsilon}|\le\varepsilon\left[
 \mu_{j+3}+\frac{\kappa_j}{2}(\mu_2+\mu_3)\right].
 \label{upd:smooth-tail}
\end{equation}
\end{proposition}
\begin{proof}
For the box, compact convergence pays its interior. Under
$u=s+r,v=s-r$, the exterior lies in
$\{\max(|u|,|v|)>U\}$ and $\dd s\dd r=\dd u\dd v/2$.
Use
$|q_0|\le(u^2+v^2)/2$, $|q_1|\le(u^4+v^4)/2$, and
$|b_E|\le1/2+(u^2+v^2)/4$ to obtain \eqref{upd:box-tail}.
The source envelope and
$k\log u+tu^2+9u\le u^2/2+13u\le(\pi/2)e^{4u}$ for $u\ge1$
give \eqref{upd:box-tail-bound} by integrating the remaining
exponential with its derivative.

For smooth localization put $a=e^{4s}/C,b=e^{-4s}/C$ and
$Z=ap^2+bk^2$. All positive-order derivatives, through the required
order, of $\log a,\log b,\log f$ and $\tau$ are bounded.
For $|\alpha|\le4$ it follows that
\[
 |\partial^\alpha A_{p,k}|
 \le c_\alpha f e^{-\pi Z}(1+Z)^4(1+|pk|)^4
 \le c'_\alpha f e^{-\pi Z/2}(1+|pk|^4).
\]
The elementary bounds
\[
 \sum_n e^{-\pi zn^2/2}\le1+\sqrt{2/z},\qquad
 \sum_n n^4e^{-\pi zn^2/2}\le\frac{128}{\pi^2e^2}z^{-5/2}
\]
follow from decreasing Gaussian integral comparison, absorbing
$(zn^2)^2$ into $e^{-\pi zn^2/4}$ for the second. Since $ab=C^{-2}$,
\[
 \sum_{p,k}|\partial^\alpha A_{p,k}|
                  \le C_\alpha(e^{2|s|}+e^{18|r|}).
\]
Thus a uniform integrable majorant is a constant times
$(1+\Omega^2)e^{-\varepsilon_0\Omega^2+2T\Omega}
(e^{2|s|}+e^{18|r|})$. Its negative quartic term controls every
joint escape direction, including bounded-$u$ or bounded-$v$ tubes.
This proves absolute interchange and \eqref{upd:localized-budget}.
Finally $\mu_k=\iint\Omega^ke^{2t\Omega}\Phi(s+r)\Phi(s-r)$ by
the same coordinate change. Apply
$1-\chi_\varepsilon\le\varepsilon\Omega^2$,
$|q_j|\le\Omega^{j+1}$, and $|b_E|\le(1+\Omega)/2$ to obtain
\eqref{upd:smooth-tail}.
\end{proof}

The window multiplies $\mathscr D\Pi$ after differentiation; no
integration by parts, window commutator, or boundary omission is used.
For completeness, if one moves $\mathscr D$ onto $W$ in a box with
primitive $G$, the exact boundary payment is
\[
 \iint W\mathscr DG=\iint G\mathscr DW+\mathfrak B_s+\mathfrak B_r
                                      -2\mathfrak B_{sr},
\]
where
\begin{align*}
 \mathfrak B_s&=\int_{-R}^{R}
 [WG_{sss}-W_sG_{ss}+W_{ss}G_s-W_{sss}G
                      -8(WG_s-W_sG)]_{-S}^{S}\dd r,\\
 \mathfrak B_r&=\int_{-S}^{S}
 [WG_{rrr}-W_rG_{rr}+W_{rr}G_r-W_{rrr}G
                      -8(WG_r-W_rG)]_{-R}^{R}\dd s,\\
 \mathfrak B_{sr}&=\int_{-R}^{R}[WG_{srr}-W_sG_{rr}]_{-S}^{S}\dd r
       +\int_{-S}^{S}[W_{ss}G_r-W_{ssr}G]_{-R}^{R}\dd s.
\end{align*}
These are repeated one-dimensional integration-by-parts identities.
Tangential edge derivatives have not been integrated further; doing so
would also require their corner terms.

\begin{proposition}[Growth of separate absolute sector budgets]
\label{upd:absolute-sector-growth}
Fix $t>0$ and finite $x$, and use $\kappa_j\ge0$. For sufficiently
small $\varepsilon>0$, the absolute axis budget satisfies
\[
 \iint\chi_\varepsilon e^{2t\Omega}|q_j+\kappa_jb_E|
          \left|\mathscr D\Pi_{\rm ax}/4096\right|\dd s\dd r
 \ge c(t/\varepsilon)^{j+1}
       e^{t^2/\varepsilon-10\sqrt{t/\varepsilon}-C}.
\]
For boxes with $S_N=R_N=(\log N)/2$, its analogue retaining
$|\cos(2xs)|$ is at least
\begin{equation}
 c(\log N)^{2j+1}
             e^{(t/2)(\log N)^2-5\log N}.
 \label{upd:box-sector-growth}
\end{equation}
Constants may depend on the fixed parameters. Neither is a bound on
the signed integral.
\end{proposition}
\begin{proof}
On a small fixed strip $|s|\le s_0$ and for large $r$,
$|\mathscr D\Pi_{\rm ax}/4096|\ge ce^{-10r}$,
$b_E\ge0$, and $q_j\ge c_jr^{2j+2}$.
For the smooth window restrict further to
$R\le r\le R+1$, $R=\sqrt{t/\varepsilon}$, and use
\[
 2t\Omega-\varepsilon\Omega^2
       =t^2/\varepsilon-\varepsilon(\Omega-t/\varepsilon)^2.
\]
The final term is bounded there, proving the first claim.  An extra
$|\cos(2xs)|$ can be retained by shrinking $s_0$.
For the box use $R_N-1/R_N\le r\le R_N$ and
$2tr^2-10r\ge2tR_N^2-10R_N-4t$.  The strip width is $1/R_N$,
which leaves the power $R_N^{2j+1}$ and gives
\eqref{upd:box-sector-growth}.
\end{proof}

Choosing $\varepsilon=e^{-cN^2}$ makes the complete smooth exterior
Gaussian-small but this absolute sector budget at least
$\exp\{t^2e^{cN^2}-O(e^{cN^2/2})\}$.  The expanding box has a
Gaussian-small exterior by \eqref{upd:box-tail-bound}, but its absolute
sector cost exceeds every fixed power of $N$.  These facts reject the
corresponding absolute-payment strategies. They do not reject a signed
bound carrying an $E$ factor.  In particular,
\eqref{upd:localized-budget} remains an identity, not a proof that
$R_j^\chi+T_j^\chi\ge-M_j^\chi$ on the full $S_j$.

\subsection{Convex dyadic summation and an explicit conditional budget}

\begin{proposition}[Removal of the dyadic block-count factor]
\label{upd:convex-dyadic}
On the same discs and sufficiently large domain as
Lemma~\ref{lem:weighted-moments}, for every fixed integer $r\ge0$,
\begin{equation}
 \begin{gathered}
 \left|\sum_{e^{\ell/5}<n\le N_{\rm RS}}(\log n)^r
                e^{t\log^2n/4}n^{-w_\pm}\right|
 \le G_0C_r\ell^r e^{-m\ell},\\
 G_0=\frac{962773}{671}<1435,\qquad
 m=\frac{76399}{1273600000}>\frac1{20000}.
 \end{gathered}
 \label{upd:convex-remainder}
\end{equation}
The common imported block constant $C_r$ remains unevaluated.
\end{proposition}
\begin{proof}
Use the exact grid $M_q=e^{\ell/5}2^q$ and assign each block by its
left endpoint $\beta_q=1/5+q\log2/\ell$ to the existing five
intervals and pairs.  In an interval $[a,b]$ of width $\Delta$,
the quadratic exponent satisfies
\[
 F(\beta;\theta_*)=\frac{b-\beta}{\Delta}F(a;\theta_*)
 +\frac{\beta-a}{\Delta}F(b;\theta_*)
 -\frac{\theta_*}{4}(\beta-a)(b-\beta).
\]
With $d(\beta)=\min(\beta-a,b-\beta)$, the paid exponent is at most
$-m-c_\Delta d(\beta)$, $c_\Delta=61\Delta/192$.
On each half interval the distances of successive grid points differ
by $\log2/\ell$, so their total additional weight is at most
$2/(1-2^{-c_\Delta})$, even for a shifted grid.  A crossing block
is valid by the imported maximal-interval estimate; shortened terminal
blocks are included. There is at most one left endpoint above $1/2$,
already paid by the extra $\rho$ in $m$. The five widths are
$1/10,3/40,3/100,33/1400,1/14$. Thus the total multiplier is less than
\[
 1+2\sum_\Delta(1-2^{-c_\Delta})^{-1}
 <11+\frac{768}{61}\sum_\Delta\Delta^{-1}
                         =\frac{962773}{671}.
\]
Here $(1-e^{-v})^{-1}\le1+v^{-1}$ and $\log2>1/2$ were used.
Multiplying by the common block majorant proves the result.
\end{proof}

The five fixed limiting pairs themselves have an exact majorant
bottleneck. Their affine lower envelope changes at
\[
 \frac{1430}{4789},\quad\frac{362}{957},\quad
 \frac{62831}{155153},\quad\frac37.
\]
For $q_{k,l}(\beta)=4\{k+(l-k-1/2)\beta\}/[\beta(1-\beta)]$,
the derivative has the sign of
$(l-k-1/2)\beta^2+2k\beta-k$, which is increasing for $\beta>0$.
Thus each envelope interval has its maximum at an endpoint. Comparing
the endpoint fractions yields
\begin{equation}
 \max_{1/5\le\beta\le1/2}\min_{(k,l)}q_{k,l}(\beta)
    =\frac{4911678521}{1933561194}
    \simeq2.5402239847600088,
 \label{upd:five-pair-bottleneck}
\end{equation}
at the $D,R$ crossover.  This does not pay the positive exponent losses
needed at a closed boundary.  At $\theta=5/2$, all five weighted
majorants have strictly positive exponent on
\[
 \frac{13}{45}<\beta<\frac{105-\sqrt{105}}{210}.
\]
Indeed the $A^2P$ roots are $4/23,13/45$, the $P$ lower root is the
right endpoint displayed, and the other three quadratics have negative
discriminant and positive leading coefficient. At $\beta=2/5$ the
least exponent is $9/3980>0$. These are limitations of these fixed
majorants, not lower bounds for the actual exponential sum.

\begin{proposition}[Elementary constants with the arithmetic input isolated]
\label{upd:explicit-weighted-budget}
For $\eta=\rho=10^{-5}$ suppose the five assigned pairs have numerical
constants $C_i>0$ such that
\[
 \sup_{I\subset(M,2M]}\left|\sum_{n\in I}n^{-iT}\right|
 \le C_iT^{k_i+\eta}M^{l_i-k_i},\qquad M\ge1, T\ge2\pi M.
\]
Put $C_{\rm EP}=\max_iC_i$,
$A_0=36138917/16$, and
\[
 \lambda=\frac{21807}{254720000}-\frac{63}{50}\rho
        =\frac{2324691}{31840000000}>\frac1{14000}.
\]
If $L\ge100002$, $0<t\le1/2$, $tL\ge61/24$, and
$|z-x|\le L^{-1}$, then, with $\ell=\log(\Re z/(4\pi))$,
\begin{equation}
 \left|\sum_{n\le N_{\rm RS}}e^{t\log^2n/4}n^{-w_\pm(z)}
                      -\zeta(w_\pm(z))\right|
 \le\frac{A_0}{\ell}+16e^{-\ell/40}
                    +2G_0C_{\rm EP}e^{-\lambda\ell}.
 \label{upd:explicit-sum}
\end{equation}
If the imported constants require $T\ge T_{\rm EP}$, also assume
$x-L^{-1}\ge3T_{\rm EP}$.  Their numerical values are an unresolved
input; qualitative exponent-pair existence does not assign them a value.
\end{proposition}
\begin{proof}
Write $z=u+iv$, $w_\pm=\sigma+i\tau$. The explicit carrier gives
\begin{gather*}
 |\ell-L|\le[L(x-L^{-1})]^{-1},\quad u\ge10^{12},
 \quad\theta_u\ge\theta_0:=\theta_*-\rho/100,\\
 \sigma\ge\tfrac12+\theta_u/4-\delta,
 \quad\delta=(2L)^{-1}+3/(4u)<\rho/2,\quad
 u/3\le|\tau|\le u,\\
 \log N_{\rm RS}\le\ell/2+e^{-\ell}/64.
\end{gather*}
The rational terms of $\alpha$ have total modulus at most $3/u$,
which proves these bounds directly from \eqref{eq:alpha-Mt}.
The phase condition $|\tau|\ge2\pi M$ follows for every retained
block. Split at $Y=e^{\ell/5}$. The low-index discrepancy is at most
\[
 \frac{\theta_u}{4\ell}\sum_{n\ge2}(\log n)^2
                             n^{-1/2-\theta_u/5+\delta}.
\]
Write $\theta_u=\theta_0+s$. Since
$1/2+\theta_0/5-\rho/2>1+1/121$ and
$(\theta_0+s)n^{-s/5}\le\theta_*+5/(e\log n)$, it suffices to use
the unimodal sum--integral bounds
\[
 \sum_{n\ge2}(\log n)^j n^{-1-1/121}
      \le j!121^{j+1}+j^j=:D_j\quad(j=1,2).
\]
These follow by comparing the increasing and decreasing parts with
adjacent unit integrals and paying the central discrepancy by the
maximum.  As $D_1=14642,D_2=3543126,5/e<2$, the low-index term is
at most $(\theta_*D_2+2D_1)/(4\ell)=A_0/\ell$.  Also
$\sigma\ge9/8$, so the ordinary zeta tail is at most
$8\lfloor Y\rfloor^{-1/8}\le16e^{-\ell/40}$.

The actual heat weight is decreasing on every high block. Abel
summation costs exactly its left endpoint weight, since endpoint plus
variation is $W(R)+\int_M^R|W'|=W(M)$.  Moreover
$T^{k+\eta}\le(4\pi)^{k+\eta}e^{(k+\eta)\ell}
<2e^{(k+\eta)\ell}$. For $\beta\le1/2$, the perturbation of the
weighted exponent is at most
$\eta+\rho/4+\rho/1600<(63/50)\rho$.
Convex summation therefore gives the last term of
\eqref{upd:explicit-sum}.  The possible sole block above $1/2$,
assigned to $P$, has exponent at most
\[
 -\frac{11}{2688}+\rho\left(1+\frac1{1600}+\frac{33}{128}\right)
                 +\frac1{32768}<-\frac1{300}<-\lambda.
\]
This pays its small endpoint displacement separately and completes
the estimate.
\end{proof}

The imported Riemann--Siegel errors and analytic replacement can also
be paid explicitly on this domain. Put $\Lambda=L-1$ and
\begin{equation}
 B_C(L)=\frac{4A_0}{\Lambda}+64e^{-\Lambda/40}
       +8G_0C e^{-\lambda\Lambda}
       +16\Lambda^2e^{-3\Lambda/4}+6e^{-2\Lambda/5}.
 \label{upd:BC-budget}
\end{equation}
Then, for the existing analytic comparison $\mathcal P_t$,
\begin{equation}
 \sup_{z\in D_x}\frac{|H(t,z)-\mathcal P_t(z)|}{\Gamma}
 \le B_{C_{\rm EP}}(L),\qquad
 \max_{0\le k\le3}
 \frac{|\partial_x^k(H-\mathcal P_t)(t,x)|}
                         {2|Q|(L/4)^k}
 \le1728 B_{C_{\rm EP}}(L).
 \label{upd:explicit-analytic-jets}
\end{equation}
Here is the complete error accounting.  In the upper half-disc the
extra coefficient $|\gamma|N^{|\kappa|}n^v$ is below two, since its
logarithm is at most
$1/(50L)+1/[2(u-6)]+e^{-\ell}/(64L)$.
Each positive heat weight is at most $2n^{-1/2}$; the complete positive
weight sum with that coefficient is at most $24e^{\ell/4}$.
The per-index error multiplier is at most $4\ell^2/u$, giving
$e_A+e_B\le8\ell^2e^{-3\ell/4}$.
The correction exponent in the third imported error is at most
$10e^{-\ell/2}+\ell e^{-\ell}<1$; its principal exponent is
$-\ell/4-\theta_u\ell/16\le-2\ell/5$. Thus
$e_{C,0}\le3e^{-2\ell/5}$.  These are bounds for all the errors in
\eqref{eq:RS-AB}--\eqref{eq:RS-C}, not only their principal terms.
The carrier estimates $|\alpha|\le\ell/2+1$, $|\alpha'|\le2/u$
give $|\partial_z\log M_t(s_\pm)|<L/3$, hence
$|M_t(s_\pm(z))|\le2\Gamma$.
Combining both branches with \eqref{upd:explicit-sum}, and using
Schwarz reflection and the monotonicity of each error at
$\ell\ge\Lambda$, proves the first inequality in
\eqref{upd:explicit-analytic-jets}. Since $\zeta(9/8)\le9$,
$|Q|\ge\Gamma/9$; Cauchy's formula gives the second with
$\max_{k\le3}(9/2)k!4^k=1728$.

For example, \eqref{upd:explicit-sum} is at most $\epsilon\in(0,1)$
if
\[
 \ell\ge\max\left\{100001,\frac{3A_0}{\epsilon},
 40\log\frac{48}{\epsilon},
 14000\log\frac{6G_0C_{\rm EP}}{\epsilon}\right\},
\]
with the optional imported threshold retained. This is a conditional
formula in the actual arithmetic constant. The separate differentiation
error from $\mathcal P_t$ to its phase vector is not part of
\eqref{upd:explicit-analytic-jets}.  Hence no numerical full sign
threshold or effective bridge follows; the floor cutoff has never been
differentiated.

\subsection{Six complete-source stationary baseline roots}
\label{upd:stationary-baseline}

For $\theta\in\{1/2,3/4,1\}$ put
$t_\theta(x)=\theta/\log(x/(4\pi))$.  Fresh complete-source interval
calculations certify exactly one root of each indicated full stationary
observation in each final bracket below, and certify a positive
stationary product $\Lag_j=-h_jh_{j+2}$ there.  The final column gives
a conservative lower bound after multiplying by $e^{\pi x_c/4}$,
where $x_c$ is that bracket's midpoint.
\begin{center}\small
\begin{tabular}{@{}llll@{}}
\toprule
$\theta$ & Full condition & Final $x$ bracket & Scaled lower bound\\
\midrule
$1/2$ &$h_1=0$ &$[3249/32,6499/64]$ &$>17532$\\
$1/2$ &$h_2=0$ &$[3313/32,6627/64]$ &$>8914$\\
$3/4$ &$h_1=0$ &$[3247/32,6495/64]$ &$>18207$\\
$3/4$ &$h_2=0$ &$[3311/32,6623/64]$ &$>9226$\\
$1$ &$h_1=0$ &$[6489/64,3245/32]$ &$>18843$\\
$1$ &$h_2=0$ &$[6617/64,3309/32]$ &$>9524$\\
\bottomrule
\end{tabular}
\end{center}
The protocol fixed the initial windows $[101,102]$, $[103,104]$,
six bisections, contour displacement $a=\pi/16$, source cutoff eight,
integration cutoff $3/2$, 192-bit Arb/Acb arithmetic, spatial Taylor
order eight, temporal order four, and $\kappa_0=\kappa_1=0$ before
these evaluations. All six cases are retained. The windows came from
known branch neighborhoods, so this is not an independent holdout of
the prior height cover.

The heat equation gives the curve derivative
\[
 \frac{d}{dx}h_{j+1}(t_\theta(x),x)
 =h_{j+2}+\frac{\theta}{x\log^2(x/(4\pi))}h_{j+3}.
\]
Certified opposite endpoint signs give existence; the derivative is
strictly negative for the three $h_1$ brackets and strictly positive
for the three $h_2$ brackets, proving uniqueness only in the final
brackets.  In particular $h_4$ is required on the second branch.
The computation uses independent paid center jets through order 17.
For rectangle half-widths $d_x,d_t$ and complete normalized contour
moments $\widehat M_k$, its additional Taylor payments are
\[
 R_{x,k}=\sum_{\ell=0}^{3}\frac{d_t^\ell}{\ell!}
              \frac{d_x^8}{8!}\widehat M_{k+2\ell+8},\qquad
 R_{t,k}=\frac{d_t^4}{4!}\widehat M_{k+8}\quad(0\le k\le4).
\]
Moments through order 18 suffice.  Center quadrature, complete omitted
source-index tails, integration tails, and both Taylor remainders are
paid separately.  Rectangle jets use the constant normalization
$e^{\pi x_c/8}$, whereas endpoint sign checks may use their own
positive normalization. No normalization divides by $E$.
The driver \texttt{certify\_theta\_stationary\_baseline.py} and receipt
\texttt{THETA\_STATIONARY\_BASELINE\_20261015.json} in the numerics
directory give the full interval payments and source hashes.
These bounded certificates lie within the existing height-112 cover;
they neither validate the divergent global Poisson sectors nor evaluate
the localized sector contributions.

\subsection{A larger cutoff class still excludes finite analytic collisions}

\begin{proposition}[Conditional subpower cutoff transfer]
\label{upd:subpower-transfer}
On a declared positive-time domain suppose, for both full stationary
sets $S_j$, that one has proved
\begin{equation}
 \Lag_j\ge-A_jN^p e^{a_j(\log N)^2}E-B_je^{-c_jN^2}
 \quad\hbox{for every integer }N\ge N_{0,j}\ge1.
 \label{upd:subpower-premise}
\end{equation}
Here $p\ge0$ is fixed; the nonnegative $A_j,B_j,a_j$, positive $c_j$,
and integer $N_{0,j}$ are source-prescribed independently of $E,N$,
and $A_j,B_j,a_j,c_j^{-1},N_{0,j}$ are locally bounded through every
covered finite positive-time event. Fix $M>0$, put
\[
 d_{B,j}=\log_+(B_j/M)\quad(B_j=0:\ d_{B,j}=0),\quad
 D_j=N_{0,j}+1+c_j^{-1/2}\sqrt{1+d_{B,j}},\quad
 d_j=\frac{\log D_j}{\log2}+\frac12,
\]
and $K_j=A_j+1$, $C_j=a_jd_j^2+pd_j/\log2$. Then
\begin{equation}
 \Lag_j\ge-K_jE\exp\left[
 C_j\{\log(2+\log_+(M/E))\}^2\right],
 \label{upd:subpower-loss}
\end{equation}
with right side zero at $E=0$.  If the domain contains the requisite
predecessor neighborhoods on both stationary sets, this estimate
excludes every finite positive-time analytic collision.
\end{proposition}
\begin{proof}
For $E>0$, only after the entire family
\eqref{upd:subpower-premise} is proved, choose
\[
 N=\left\lceil\max\{N_{0,j},
                 \sqrt{c_j^{-1}\log_+(B_j/E)}\}\right\rceil.
\]
Then $B_je^{-c_jN^2}\le E$ and
$N\le D_j\sqrt{1+\log_+(M/E)}$.
Writing $\eta_E=\log(2+\log_+(M/E))\ge\log2$, we obtain
$\log N\le d_j\eta_E$, whence
$N^p e^{a_j(\log N)^2}\le e^{C_j\eta_E^2}$.
This proves \eqref{upd:subpower-loss}, including $E\ge M$; at
$E=0$ let $N\to\infty$ in the original family.
On a predecessor with $s=T-t\downarrow0$ the already proved local
collision expansion has $E\sim k^2s^{2r}$ and
$-\Lag_j/E=r/s+O(1)$, with finite $r\ge1$.
Local boundedness bounds the right-side coefficient in
\eqref{upd:subpower-loss} by
$\exp\{O((\log\log(1/s))^2)\}=s^{-o(1)}$,
contradicting the pole $r/s$.
\end{proof}

The same transfer applies to the full-amplitude observation
$J_{j,N}=-h_jh_{j+2}^{(N)}$ if its original complete-source payment
$|J_{j,N}-\Lag_j|\le Q_j\sqrt E e^{-c_{0,j}N^2}$ is retained.
Young's inequality replaces
\[
 A_j\mapsto A_j+\tfrac12,\qquad
 B_j\mapsto B_j+\tfrac12Q_j^2,\qquad
 c_j\mapsto\min(c_j,2c_{0,j}).
\]
Use $N_{0,j}\ge8$ when importing the original staircase tail, and
require $Q_j,c_{0,j}^{-1}$ locally bounded. All stationary conditions,
source boundaries, and mixed-sector payments remain those of the full
physical construction; a rectangular or smooth Poisson window is not
a replacement for the original staircase theorem.

\begin{lemma}[Analytic rigidity for every subpower gauge]
\label{upd:subpower-rigidity}
Let $S$ be semianalytic near a finite point, and let $E\ge0,L$ be
real analytic on an ambient neighborhood. Suppose $G(e)>0$ obeys
$G(e)=o(e^{-\epsilon})$ as $e\downarrow0$ for every $\epsilon>0$.
If $L\ge-K E G(E)$ on $S$ near that point, with bounded $K\ge0$
and right side zero at $E=0$, then $L\ge-CE$ on a smaller
neighborhood in $S$, for some finite $C\ge0$.
\end{lemma}
\begin{proof}
Failure gives $E>0,-L>0$ and $E/(-L)\to0$ near a limit point
with $E=0$. Apply real analytic curve selection, as in
Lemma~\ref{lem:analytic-log-rigidity}, to
\[
 \{(z,u):z\in S, E(z)>0, -L(z)>0, u>0, E(z)=u[-L(z)]\}.
\]
The resulting analytic arc has
$u(\rho)=a\rho^k(1+o(1))$ and
$E(z(\rho))=b\rho^m(1+o(1))$, with $a,b>0,k,m\ge1$.
The inequality would give $u(\rho)^{-1}\le KG(E(z(\rho)))$.
Taking $\epsilon=k/(2m)$ makes its right side
$o(\rho^{-k/2})$, contradicting the order $\rho^{-k}$ on the left.
At $E=0$ the assumed inequality already gives $L\ge0$.
\end{proof}
The gauge in \eqref{upd:subpower-loss} is subpower, so this lemma
applies locally after bounding $C_j$. It is a qualitative consequence,
not an effective construction of a source-prescribed linear coefficient.
In particular, the absolute sector growth in
\eqref{upd:box-sector-growth} lacks the required $E$ factor and supplies
none of the hypothesis \eqref{upd:subpower-premise}.

The supporting finite checks comprise 27 exact algebra or sensitivity
assertions and 12 ball comparisons for the parity identity, 260 exact
assertions for mode and cutoff algebra, 129 for convex dyadic arithmetic,
and 50 for the explicit weighted budget. Their saved scripts and receipts
are listed in research Notes 38--44 and internal Reviews 27--28.
These checks support normalizations and arithmetic; the infinite-sum,
integrability, and conditional collision conclusions use the proofs
above. No finite sampling establishes the missing all-cutoff signed
premise, and no RH conclusion follows from this continuation.


\section{Conditional transfers and the remaining research program}
\label{pg:research-program}

The preceding outer theorem does not cover the entire positive-time
collision domain.  This section records three further deductions that
clarify possible next steps.  Their transfer implications are proved;
the arithmetic hypotheses proposed below are not established here.
In particular, this section adds no unconditional sign region.
Write $S_j=\{h_{j+1}=0\}$ for $j=0,1$, and retain
$E=h_0^2+h_1^2$.  All stationary conditions below refer to the full
physical jets.

\subsection{A weighted remainder sufficient for the closed strip}

Retain the branches $s_\pm,w_\pm,Q_\pm$ of
\eqref{eq:alpha-Mt}--\eqref{eq:analytic-model} and the analytic model
$\mathcal P_t=Q_++Q_-$.  For a real center $x$, put
\[
 L=\log(x/(4\pi)),\qquad D_x=\{z:|z-x|\le L^{-1}\},\qquad
 \Gamma(t,x)=|M_t(s_+(x))|.
\]
For $z=u+iv\in D_x$, define
\[
 \ell=\log(u/(4\pi)),\qquad Y=e^{19\ell/100},\qquad
 N_{\rm RS}=\left\lfloor\sqrt{u/(4\pi)+t/16}\right\rfloor.
\]
The notation $N_{\rm RS}$ distinguishes this cutoff from the
staircase source cutoff used below.

\begin{proposition}[Conditional weighted-remainder transfer]
\label{pg:weighted-transfer}
Suppose that, as $x\to\infty$, for both signs one has
\begin{equation}
 \sum_{Y<n\le N_{\rm RS}}e^{t\log^2n/4}n^{-w_\pm(z)}=o(1)
 \label{pg:weighted-target}
\end{equation}
uniformly for real centers satisfying
\[
 0<t\le\tfrac12,\qquad \tfrac52\le tL\le\tfrac83,
 \qquad z\in D_x.
\]
Then, with $Q=Q_+(x)\ne0$, $\omega=L/4$, and $\phi=\arg Q$,
\begin{equation}
 \frac{1}{2|Q|}\left(h_0,\frac{h_1}{\omega},
                 \frac{h_2}{\omega^2},\frac{h_3}{\omega^3}\right)
 = (\cos\phi,-\sin\phi,-\cos\phi,\sin\phi)+o(1)
 \label{pg:weak-four-jet}
\end{equation}
uniformly on this closed strip.  Consequently $E,\Lag_0,\Lag_1$
are strictly positive there at sufficiently large height.  The
hypothesis alone does not imply an $O(L^{-1})$ error in
\eqref{pg:weak-four-jet}.
\end{proposition}

\begin{proof}
The same disc estimates as in the closed-strip argument give
$\Re w_\pm\ge11/10$ at sufficiently large height.  On $n\le Y$,
$e^q-1\le qe^q$ bounds the discrepancy from the ordinary Dirichlet
series by
\[
 \frac t4\sum_{n\ge2}(\log n)^2
 n^{-\{1/2+(81/400)t\ell-O(L^{-1})\}}=O(L^{-1}).
\]
Indeed, at the real lower boundary the exponent tends to
$1/2+(81/400)(5/2)=161/160$, and disc perturbations preserve
the lower bound $1+1/320$.  Also $t=O(L^{-1})$ on this fixed
strip.  The ordinary Dirichlet tail beyond $Y$ is
$O(e^{-19\ell/1000})$.  Adding \eqref{pg:weighted-target} therefore
gives $S_{t,N_{\rm RS}}(w_\pm)-\zeta(w_\pm)=o(1)$ on the disc.

The complete imported value errors in \eqref{eq:RS-input} remain
negligible: $e_A+e_B=O(L^2u^{-3/4})$ and
$e_{C,0}=O(u^{-2/5})$ suffice here.  Carrier magnitudes on $D_x$
are comparable to $\Gamma$.  The upper-half-disc approximation and
Schwarz reflection thus yield
\[
 H(t,z)-\mathcal P_t(z)=o(\Gamma),\qquad z\in D_x.
\]
This difference is analytic even though the intermediate cutoffs
can jump.  Cauchy's estimates, applied only to this analytic
difference, give errors $o(\Gamma L^k)$ for $0\le k\le3$ at $x$.

In $\Re w\ge11/10$, the zeta factors are bounded above and away
from zero, and their fixed-order logarithmic derivatives are bounded.
The explicit carrier differentiation consequently gives
\[
 Q_\pm^{(k)}(x)=(\pm i\omega)^kQ_\pm(x)
                         +O(\Gamma L^{k-1}),\qquad 1\le k\le3.
\]
Together with $Q_-(x)=\overline{Q_+(x)}$ and $|Q|\asymp\Gamma$,
this proves \eqref{pg:weak-four-jet}.  Substitution into the quadratic
forms gives
\[
 h_0^2+(h_1/\omega)^2=4|Q|^2(1+o(1)),\qquad
 \Lag_j=4|Q|^2\omega^{2j+2}(1+o(1)),
\]
so the strict signs follow uniformly in phase.
\end{proof}

Thus a uniform logarithmic saving for the actual weighted remainder
can suffice.  The hypothesis retains its full disc uniformity and
both endpoints; a pointwise real-axis estimate is not a substitute.
An estimate obtained only after combining the two branches may also
suffice if it directly proves the displayed analytic comparison
$H-\mathcal P_t=o(\Gamma)$.  Neither alternative is proved here.

\subsection{A full-amplitude cutoff and its conditional loss transfer}

Recall the exact staircase domain
\[
 \Omega_N=\{(v,n):0\le v\le\tfrac12\log N,\ n\le e^{2v}\},
 \qquad N\ge8,
\]
and its retained jets $h_k^{(N)}$.  Retain $a=\pi/16$, $\mathsf M_k$, $\mathsf D_0$, and $c$ from
\eqref{eq:signed-cutoff-constants}, and put $c_0=c$.
In particular $|h_k|,|h_k^{(N)}|\le e^{-ax}\mathsf M_k$ for
$0\le t\le1/2$ and $x\ge0$.  The complete staircase tail bound is
\begin{equation}
 |h_k-h_k^{(N)}|\le e^{-ax}\mathsf D_0e^{-c_0N^2},\qquad
 \mathsf D_0=\frac{\mathsf C_0}{2\beta_0},\quad c_0=\frac{\beta_0}{2}.
 \label{pg:jet-tail}
\end{equation}
It includes all omitted source indices and coordinates.

On $S_j$ define the full-amplitude observation
\begin{equation}
 J_{j,N}:=\mathcal J_j^{(N)}=-h_jh_{j+2}^{(N)},\qquad j=0,1.
 \label{pg:full-amplitude}
\end{equation}
This notation is distinct from the two-source integrand $J_j(\alpha,\beta)$.
Full stationarity gives the exact identity and estimate
\begin{equation}
 J_{j,N}-\Lag_j=h_j(h_{j+2}-h_{j+2}^{(N)}),\qquad
 |J_{j,N}-\Lag_j|\le D_j\sqrt E\,e^{-c_0N^2},
 \quad D_j=e^{-ax}\mathsf D_0.
 \label{pg:J-tail}
\end{equation}
In particular $J_{j,N}=0$ at $E=0$ for both levels.  No stationarity
condition is imposed on the retained jets.

\begin{proposition}[Conditional full-amplitude cutoff transfer]
\label{pg:J-transfer}
Let $\mathcal D$ be a declared domain in $0<t\le1/2$, $x>0$.
Suppose, on $S_j\cap\mathcal D$, that the source-prescribed inequality
\begin{equation}
 J_{j,N}\ge-A_jN^pE-B_je^{-cN^2}
 \quad\hbox{for every integer }N\ge N_0\ge8
 \label{pg:J-target}
\end{equation}
holds with a fixed finite $p\ge0$, nonnegative $A_j,B_j$, and $c>0$.
Assume $A_j,B_j,c^{-1},N_0$ are locally bounded through candidate
positive-time collisions.  Then, for each fixed $M>0$, a
source-prescribed locally bounded $K_j$ satisfies
\begin{equation}
 \Lag_j\ge-K_jE\,[1+\log_+(M/E)]^{p/2}
 \quad\hbox{on }S_j\cap\mathcal D,
 \label{pg:J-log-loss}
\end{equation}
where the right side is zero at $E=0$ and
$\log_+v=\max\{0,\log v\}$.  If this is proved for both levels
on neighborhoods containing the predecessor branches of every
candidate collision in $\mathcal D$, it excludes those collisions.
\end{proposition}

\begin{proof}
Young's inequality in \eqref{pg:J-tail} gives
\[
 D_j\sqrt E\,e^{-c_0N^2}
 \le\tfrac12E+\tfrac12D_j^2e^{-2c_0N^2}.
\]
Set $\widehat A_j=A_j+1/2$, $\widehat B_j=B_j+D_j^2/2$, and
$\widehat c=\min\{c,2c_0\}$.  Since $N^p\ge1$,
\[
 \Lag_j\ge-\widehat A_jN^pE-\widehat B_je^{-\widehat cN^2}.
\]
For $E>0$ choose
\[
 N=\left\lceil\max\left\{N_0,
          \sqrt{\widehat c^{-1}\log_+(\widehat B_j/E)}
                           \right\}\right\rceil.
\]
Put $\ell_E=1+\log_+(M/E)$ and
$d_j=\log_+(\widehat B_j/M)$.  Then the tail is at most $E$, and
\[
 N\le\Delta_j\sqrt{\ell_E},\qquad
 \Delta_j=N_0+1+\widehat c^{-1/2}\sqrt{1+d_j}.
\]
Thus \eqref{pg:J-log-loss} holds with
$K_j=\widehat A_j\Delta_j^p+1$, which is locally bounded under the
stated hypotheses.  At $E=0$, let $N\to\infty$ in the full-form
family.  All cutoff inequalities precede this optimization; none of
their coefficients is defined using a measured adverse quotient.

For the collision conclusion, the predecessor calculation in
Lemma~\ref{fc:predecessor}, applied to $H$ or $H_x$ according to
the parity of the multiplicity, gives
\[
 E\asymp(T-t)^{2r},\qquad
 -\Lag_j/E=\frac{r}{T-t}+O(1),\qquad r\ge1.
\]
The locally bounded right coefficient in \eqref{pg:J-log-loss}
permits at most a fixed power of $1+|\log(T-t)|$, contradicting
this pole.  This proves exclusion only in the covered neighborhoods.
\end{proof}

The proposition proves a transfer, not \eqref{pg:J-target} for the
theta source in the unbounded inner cusp.  A bound only for one chosen
cutoff, or with coefficients that depend on a sampled energy floor,
does not establish its hypothesis.

\subsection{Localizing the missing estimate near small energy}

The nonzero scale $\Gamma$ defined from \eqref{eq:M0} and
\eqref{eq:alpha-Mt} contains no zeta factor or physical jet denominator.
It can therefore be used even where a zeta-carrier nonvanishing claim
has not been established.  The following observation separates a
simple region from the difficult cancellation regime.

\begin{lemma}[Small-energy localization]
\label{pg:small-energy}
Let $\mathcal D\subset(0,1/2]\times(0,\infty)$ be a declared domain
on which $\Gamma$ is continuous and strictly positive, and fix $\delta>0$.
For either observation
$\mathcal O_{j,N}=J_{j,N}$ or $\mathcal O_{j,N}=\Lag_j^{(N)}$,
suppose a polynomial-cutoff lower bound
\[
 \mathcal O_{j,N}\ge-A_jN^pE-B_je^{-cN^2}
\]
has been proved for all $N\ge N_0\ge8$ on
\begin{equation}
 S_j\cap\mathcal D\cap\{E\le\delta\Gamma^2\}.
 \label{pg:small-energy-domain}
\end{equation}
Here $p\ge0$ is fixed and finite, and the prescribed nonnegative
coefficients are defined on $\mathcal D$ with the local boundedness
conditions of Proposition~\ref{pg:J-transfer}.  The same type of
bound then holds on all of $S_j\cap\mathcal D$ after replacing $A_j$ by
\begin{equation}
 \widetilde A_j=
 \max\left\{A_j,
   \frac{e^{-2ax}\mathsf M_j\mathsf M_{j+2}}
        {\delta\Gamma^2}\right\}.
 \label{pg:extended-coefficient}
\end{equation}
This new coefficient is locally bounded and independent of $N$.
\end{lemma}

\begin{proof}
For $J_{j,N}$ the separate complete and retained jet envelopes give
$|J_{j,N}|\le U_j$, where
$U_j=e^{-2ax}\mathsf M_j\mathsf M_{j+2}$.
For the raw cutoff the complete paired identity gives the same
envelope $|\Lag_j^{(N)}|\le U_j$.  On $E>\delta\Gamma^2$,
\[
 \mathcal O_{j,N}\ge-U_j
       \ge-\frac{U_j}{\delta\Gamma^2}E.
\]
Use $N^p\ge1$ and combine this with the assumed bound on
\eqref{pg:small-energy-domain}.  Local boundedness follows from
continuity and positivity of $\Gamma$ at every finite point.
\end{proof}

For any hypothetical finite collision with $\Gamma>0$ whose sufficiently
nearby predecessor points stay in $\mathcal D$, continuity and $E=0$
at the collision place those points in \eqref{pg:small-energy-domain}.
Thus, wherever the declared domain covers these predecessors, the
restricted small-energy region contains the essential collision obstruction.  The
localization assumes no positive energy floor there and does not prove
the missing lower bound.  Growth of the coefficient as height tends to
infinity or as $t\downarrow0$ is compatible with the local criterion.

\subsection{Arithmetic directions and limits of the present coverage}

The principal cusp target remains \eqref{pg:J-target}, preferably on
the small-energy region \eqref{pg:small-energy-domain}. The complete
correlations in Section~\ref{sec:complete-correlations} give a focused
way to investigate it: derive an exact identity for
$\mathscr A_j+\kappa_j\mathscr E$, or for the hybrid kernel
$\mathscr C_{j,N}^{a}$, and estimate its adverse transform with
source-prescribed coefficients. A positive-definite component with a
paid signed remainder could suffice; global positive definiteness is
stronger than the stationary target. The complete identities, finite-character bridge, and integrable smooth
localization are now available in Sections~\ref{sf:section}--\ref{pp:section}
and~\ref{upd:section}. The next deliverable is a source-prescribed signed
bound for their complete reconstruction, or an all-$N$ estimate on a
declared predecessor-covering domain. The subpower transfer in
Section~\ref{upd:section} permits a larger cutoff cost than
\eqref{pg:J-target}, but still requires its full signed premise. Pointwise positivity of a correlation density
is not that estimate.

Poisson reflection alone exchanges the two correlation coordinates
and is already shared by arbitrary even sources. An additional theta
relation must improve the signed remainder, retain the logarithmic
heat factor, and pay every transformed boundary. The hybrid formula
\eqref{eq:corr-hybrid-transform} keeps the original staircase inside
$\psi_{a,N}$, so it cannot justify discarding entering-index terms.
The reflected lift and planar dictionary are useful finite algebraic
reductions; further triality, integral-operator, or exceptional-theta
work should proceed only after specifying a source- and
observation-preserving identity that could improve this estimate.

The cumulative identity \eqref{eq:cumulative-signed-identity}
provides a related diagnostic. Grouping consecutive oscillatory lobes
before estimating the integrated signed overshoot may demand less than
\eqref{eq:cumulative-overshoot-target}. The complete tail must still be
paid at the physical energy scale; a merely absolute error proportional
to $|h_j|$ is insufficient near a collision.

The central arithmetic gap is not removed by the published systematic
exponent table \cite[Table 1]{taotrudgianyang25}. At block scale
$\beta=2/5$, its bound is
\[
 \frac{18}{199}+\frac{521}{796}\frac25=\frac{701}{1990},
 \qquad \frac{701}{1990}-\frac7{20}=\frac9{3980}>0.
\]
Here $7/20$ is the exponent allowed by the $\theta=5/2$ target before
its additional saving. This is a limitation of the quoted bounds, not
an obstruction to a logarithmic-phase-specific or signed estimate.

The outer program seeks \eqref{pg:weighted-target} through exact
partitions, transformed logarithmic-phase sums, or cancellation across
neighboring scales.  A small diagnostic computation should specify
its observation and coefficient rule before validation, retain complete
tails, and increase the cutoff and precision when the physical carrier
scale requires it.  Along a sampling curve $t=\theta/L(x)$ the derivative
used to isolate a root of the physical stationary observation is
\[
 \frac d{dx}h_{j+1}(\theta/L(x),x)
   =h_{j+2}+\frac{\theta}{xL(x)^2}h_{j+3}.
\]
An undecided enclosure is not a sign violation; favorable finite cells
do not establish an unbounded-domain estimate.

Finally, the nominal zeta argument has real part $1/2+\theta/4$.
The elementary nonvanishing input for the present carrier works with
a fixed margin when $\theta>2$.  At and below two that justification
no longer applies.  A different argument may cross the boundary, but
dividing by a zeta factor whose nonvanishing has not been proved would
import the missing information.  This is a limitation of the present
method, not an impossibility theorem for all carrier methods.

Even a sequence of outer thresholds $\theta_n\downarrow0$ does not
settle the global problem without controlling its associated starting
heights $X_n$.  At a fixed hypothetical collision $(T,x_*)$, choosing
$n$ so that $\theta_n<T\log(x_*/(4\pi))$ may also give $X_n>x_*$.
A finite-height or other global argument must close that quantifier
gap.  Effective bridging should therefore target an evaluated reachable
analytic cutoff; repeated bounded extensions alone do not address the
unbounded small-time cusp.

\appendix
\section{Complete source and integration error bounds}\label{app:payments}

The following bounds make the numerical and shared-source constructions
independent of an unpaid finite theta truncation. Set
\[
 S_k(q,d)=\sum_{r=0}^k\binom kr\frac{d^{k-r}r!}{q^{r+1}}
         =\int_0^\infty(v+d)^ke^{-qv}\dd v\qquad(q>0).
\]
Let $0\le t\le\tau\le1/2$, $0<a<\pi/8$, and $x\ge0$.
For an index cutoff $N$, put $m=N+1$ and
\[
 c_a=\cos4a,\quad \alpha=\pi m^2,\quad \beta=c_a\alpha,
 \quad q_N=4\beta-9,\quad
 \rho_N=((m+1)/m)^4e^{-\pi c_a(2m+1)}.
\]
If $\rho_N<1$, $q_N>0$, and $8\beta\ge\tau$, the complete omitted
index contribution in \eqref{eq:complete-shift} is at most
\begin{equation}
 E_{N,k}=e^{-ax}\frac{(2\alpha^2+3\alpha)e^{-\beta}}{1-\rho_N}
                         S_k(q_N,a).
 \label{eq:tail-index}
\end{equation}
For a cutoff $U\ge0$, put
\[
 \beta_U=\pi c_ae^{4U},\quad \rho_U=16e^{-3\beta_U},\quad
 q_U=4\beta_U-9-2\tau U,\quad B_U=\tau U^2+9U-\beta_U.
\]
If $\rho_U<1$, $q_U>0$, and $8\beta_U\ge\tau$, the entire source
on $u\ge U$ contributes at most
\begin{equation}
 E_{U,k}=e^{-ax}\frac{(2\pi^2+3\pi)e^{B_U}}{1-\rho_U}
                         S_k(q_U,U+a).
 \label{eq:tail-integration}
\end{equation}
Thus $E_{N,k}+E_{U,k}$ pays the omitted domain of an $N$-term integral
on $[0,U]$. Validated retained quadrature and cell variation are
additional errors.

To prove the bounds, use
\[
 |\phi_n(u+ia)|\le(2\alpha_n^2+3\alpha_n)
 e^{9u-c_a\alpha_ne^{4u}},\quad
 |e^{t(u+ia)^2}|\le e^{\tau u^2},\quad |u+ia|^k\le(u+a)^k.
\]
For $n\ge m$, successive coefficient-majorant ratios are bounded by
$\rho_N$ uniformly for $u\ge0$. The inequality
$e^{4u}\ge1+4u+8u^2$ gives
\[
 \tau u^2+9u-\beta e^{4u}
 \le-\beta-q_Nu-(8\beta-\tau)u^2\le-\beta-q_Nu.
\]
Summing the full geometric index bound and integrating proves
\eqref{eq:tail-index}. For $u\ge U$, the exponent
$\tau u^2+9u-\pi c_ae^{4u}$ is concave and has slope $-q_U$ at
$U$. Its tangent bound and the geometric ratio $\rho_U$ give
\eqref{eq:tail-integration}. Adding the bounds overcounts their
intersection, which is harmless.

For the centered residual $C_j$, replace any factor $S_j(q,d)$ in
these payments by $S_{j+2}(q,d)+b^2S_j(q,d)$, since
$|u^2-b^2|\le(u+a)^2+b^2$. For the shared carrier
\eqref{eq:shared-carrier}, set
\[
 \widehat\beta_V=\pi\cos(4a)e^{4V},\quad
 \widehat q_V=4\widehat\beta_V-10-2\tau V,\quad
 \widehat B_V=\tau V^2+10V-\widehat\beta_V.
\]
If $\widehat q_V>0$ and $8\widehat\beta_V\ge\tau$, the majorant
\eqref{eq:shared-majorant} and the same tangent argument give the full
tail payment
\begin{equation}
 e^{-ax}2(2\pi^2+3\pi)e^{\widehat B_V}S_j(\widehat q_V,V+a).
 \label{eq:shared-tail}
\end{equation}
The same degree-two replacement pays the shared expression for $C_j$,
because $0\le v-d_n\le v$.

For completeness, these absolute errors can be made small at the gamma
exponential scale at arbitrary height. For $x\ge41$, $p>0$, and
$\ell_x=\log x$, take
\[
 a_x=\pi/8-1/x,\quad
 m=\left\lceil\sqrt{\frac{x}{2\pi}(p\ell_x+32)}\right\rceil,
 \quad N=m-1,\quad U=\frac{\ell_x+2\log(\ell_x+8)}4.
\]
All preceding payment hypotheses hold, and for fixed $k$,
\begin{equation}
 e^{\pi x/8}E_{N,k}=O_{p,k}(x^{5/2-p}(\log x)^{3/2}),\qquad
 e^{\pi x/8}E_{U,k}=O_{k,A}(x^{-A})\quad(A>0).
 \label{eq:adaptive-payments}
\end{equation}
Indeed $c_a=\sin(4/x)\ge2/x$, $\beta\ge p\ell_x+32$, and
$m^2\ge2x/\pi$. Hence
$-\log\rho_N\ge4\pi m/x-4/m\ge2\pi m/x$ and
$(1-\rho_N)^{-1}\le1+x/(2\pi m)$.
Together with $e^{-a_xx}=e\,e^{-\pi x/8}$, substitution gives the
first bound. For the second,
$\beta_U\ge2\pi(\ell_x+8)^2$, $U\le\ell_x/2+2$, and
$B_U\le-47\ell_x^2/8-181\ell_x/2-364$.
These are absolute error estimates. They neither establish a norm gap
nor replace a signed estimate of the retained source. On a spatial cell,
freeze the contour and cutoffs before differentiating or transporting
physical jets.


% Prepared for Edward Baker; gpt-6-astra, effort ultra; 14 October 2026.
\section{An explicit height at the condition \texorpdfstring{$tL\ge5$}{tL >= 5}}
\label{app:explicit-height}

The exponent-pair argument in Section~\ref{sec:outer-proof} gives the
stronger asymptotic condition $tL\ge61/24$, with an unevaluated
height. The following elementary absolute-sum calculation supplies an
explicit height for $tL\ge5$. These are distinct quantitative statements.

\begin{proposition}[Explicit outer height]\label{prop:explicit-height}
Let $L=\log(x/(4\pi))$, $\omega=L/4$, and let $Q=Q_+(x)$ be the
amplitude in \eqref{eq:analytic-model}. If
\[
 x\ge X_{\mathrm{exp}}:=4\pi e^{10^9},\qquad
 0<t\le\tfrac12,\qquad tL\ge5,
\]
then every coordinate error in \eqref{eq:main-four-jet} is at most
\begin{equation}
 \frac{108000000}{L}\le0.108<\frac18.
 \label{eq:effective-four-jet-budget}
\end{equation}
In particular, for $P_k=h_k/(2|Q|\omega^k)$,
\[
 P_0^2+P_1^2\ge\tfrac12,\qquad
 P_1^2-P_0P_2\ge\tfrac{31}{64},\qquad
 P_2^2-P_1P_3\ge\tfrac{31}{64}.
\]
The same conclusions hold at negative heights after reflection.
\end{proposition}

\subsection{Explicit weights on the complex disc}
All bounds until the last sign step hold for $L\ge1000$,
$0<t\le1/2$, and $tL\ge5$. On $D_x$ from
\eqref{eq:shrinking-disc}, write $z=u+iv$ and $\theta_u=tL_u$. Then
\begin{equation}
 u\ge x/2\ge e^L,\qquad |L_u-L|\le\frac2{xL}<1,
 \qquad\theta_u\ge5-\frac1{xL}\ge\frac{19}{4}.
 \label{eq:effective-disc-geometry}
\end{equation}
For $s=s_\pm(z)$, one has $0\le\Re s\le1$ and
$|s|,|s-1|\ge u/2$. Formula~\eqref{eq:alpha-Mt} therefore yields
\begin{equation}
 \Re\alpha(s)\ge L_u/2-3/u,\qquad
 |\alpha(s)|\le L_u/2+1\le L,
 \qquad |\alpha'(s)|\le1/u+6/u^2\le2/u.
 \label{eq:effective-alpha-disc}
\end{equation}
For the magnitude estimate, use
$|s|\le(u/2)\sqrt{1+4/u^2}$, $|\arg s|\le\pi/2$, and the bound
$3/u$ for the rational terms. It follows that
\begin{equation}
 \sigma_\pm:=\Re w_\pm
 \ge\frac12-\frac1{2L}+\frac{\theta_u}{4}-\frac{3t}{2u}
 \ge\frac32.
 \label{eq:effective-zeta-half-plane}
\end{equation}
The cutoff satisfies
\[
 \tfrac12e^{L_u/2}\le N,\qquad\log N\le L_u/2+1/u.
\]
Indeed the upper correction is at most $\pi t/(8u)$, and flooring
reduces a number at least two by at most a factor of two.
Consequently, for $n\le N$,
\begin{equation}
 e^{t\log^2n/4}n^{-\sigma_\pm}\le n^{-p},\qquad
 p=\frac12+\frac{\theta_u}{8}-\delta,\qquad
 \delta=\frac1{2L}+\frac{7t}{4u}\le\frac1{32}.
 \label{eq:effective-weight}
\end{equation}
Writing $s_0=\theta_u-19/4\ge0$, we obtain
$p\ge17/16+s_0/8$. On the upper half-disc, the same bound holds for
the first branch after multiplication by $n^v$: its leading real part
changes from $(1+v)/2$ to $(1-v)/2$, already included in the estimate.

For $\eta>0$ and an integer $m\ge0$, comparison on $[n-1,n]$ gives
\begin{equation}
 \sum_{n=2}^{\infty}(\log n)^m n^{-1-\eta}
 \le\int_1^\infty(\log y+1)^m y^{-1-\eta}\dd y.
 \label{eq:effective-log-moments}
\end{equation}
At $m=2$, $\eta=1/16$, the right side is
$2/\eta^3+2/\eta^2+1/\eta=8720<9000$. Hence
\begin{align}
 |S_{t,N}(w_\pm)-\zeta(w_\pm)|
 &\le\frac{9000}{4L_u}(19/4+s_0)2^{-s_0/8}+2N^{-1/2}
 \notag\\
 &\le\frac{29250}{L_u}+3e^{-L_u/4}
 \le\frac{31000}{L}.
 \label{eq:effective-sum-discrepancy}
\end{align}
Here $s_0e^{-s_0\log2/8}\le8/(e\log2)<8$, so the time-dependent
numerator is less than $13$. Moreover,
$29250/(L-1)<30000/L$ for $L\ge1000$, while
$e^{(L-1)/4}\ge(L-1)^2/32\ge3L$. Thus the displayed bound includes
both the weighted discrepancy and the complete ordinary zeta tail,
uniformly even as $tL$ grows.

\subsection{All imported approximation errors}
On the upper half-disc, \eqref{eq:RS-ratios} gives $|\gamma|\le2$ and
\[
 |\kappa|\log N
 \le\frac{tv}{2(u-6)}(L_u/2+1/u)<\tfrac12,
 \qquad N^{|\kappa|}\le2.
\]
Since $\zeta(17/16)\le17$, the full positive weight sum in
\eqref{eq:RS-AB} is at most $5\cdot17=85$. Denote that equation's
exponent by
\[
 q_n=\frac{(t^2/16)\log^2(u/(4\pi n^2))+0.626}{u-6.66}.
\]
The estimates $-2/u\le\log(u/(4\pi n^2))\le L_u$,
$(L+1)^2/64+1\le L^2/32$, and $u-6.66\ge u/2$ imply
$0\le q_n\le L^2/(16u)<1$. Using $e^q-1\le2q$ on $[0,1]$ gives
\begin{equation}
 e_A+e_B\le\frac{85L^2}{8u}\le\frac{11L^2}{u}\le\frac1L.
 \label{eq:effective-AB}
\end{equation}
For the final inequality, $u\ge e^L\ge L^4/24\ge11L^3$ suffices.

The positive correction exponent in \eqref{eq:RS-C} is at most
\[
 \frac{1.24(3^v+3^{-v})}{N-0.125}
 +\frac{3|L_u+i\pi/2|+10.44}{u-12}
 \le20e^{-L_u/2}+7L/u\le1.
\]
To verify the last inequality, use
$20e^{-(L-1)/2}\le160/(L-1)^2<1/2$ and
$7Le^{-L}\le168/L^3<1/2$. Its negative main exponent therefore gives
\begin{equation}
 e_{C,0}\le e\exp\{-(1/4+19/64)L_u\}
       \le3e^{-L_u/2}\le\frac1L,
 \qquad e_A+e_B+e_{C,0}\le\frac2L.
 \label{eq:effective-C-and-total}
\end{equation}
All three errors in the imported theorem are included; the constant
$0.626$ in $q_n$ is independent of $t$.

\subsection{Amplitude comparison and all four derivatives}
Let $\Gamma=|M_t(s_+(x))|=|M_t(s_-(x))|$. On $D_x$,
\[
 \left|\frac{d}{dz}\log M_t(s_\pm(z))\right|
 \le\frac{|\alpha|}{2}(1+t|\alpha'|/2)\le L.
\]
Integration along a segment of length at most $1/L$ yields
$|M_t(s_\pm(z))|\le e\Gamma<3\Gamma$. Also
\eqref{eq:effective-zeta-half-plane} and the zeta and reciprocal-zeta
Dirichlet series give
$|\zeta(w)|,|1/\zeta(w)|\le\zeta(3/2)\le3$, hence
$\Gamma\le3|Q|$. Combining \eqref{eq:effective-sum-discrepancy} and
\eqref{eq:effective-C-and-total} with \eqref{eq:RS-input} gives
\begin{align}
 |H(t,z)-\mathcal P_t(z)|
 &\le 2(3\Gamma)\frac{31000}{L}+(3\Gamma)\frac2L
   =\frac{186006\Gamma}{L}
 \notag\\
 &\le\frac{558018|Q|}{L}\le\frac{560000|Q|}{L}.
 \label{eq:effective-disc-error}
\end{align}
Schwarz symmetry extends this upper-half-disc estimate to the full disc.
Cauchy's inequality now gives
\begin{equation}
 \frac{|\partial_x^k(H(t,\cdot)-\mathcal P_t)(x)|}
      {2|Q|\omega^k}
 \le\frac{280000\,k!\,4^k}{L}
 \le\frac{107520000}{L},\qquad 0\le k\le3.
 \label{eq:effective-cauchy-error}
\end{equation}

For the analytic carrier, at $s=(1+ix)/2$ direct differentiation of
\eqref{eq:alpha-Mt} gives
\begin{equation}
 |\alpha-(L/2+i\pi/4)|\le4/x,\quad
 |\alpha'|\le2/x,\quad |\alpha''|\le3/x^2,\quad
 |\alpha'''|\le9/x^3.
 \label{eq:effective-alpha-center}
\end{equation}
The first estimate follows from $s=(ix/2)(1-i/x)$ and
$|\operatorname{Log}(1-i/x)|\le2/x$. For the others one uses
\[
 \alpha'=-\frac1{2s^2}-\frac1{(s-1)^2}+\frac1{2s},\qquad
 \alpha''=\frac1{s^3}+\frac2{(s-1)^3}-\frac1{2s^2},\qquad
 \alpha'''=-\frac3{s^4}-\frac6{(s-1)^4}+\frac1{s^3}.
\]
Let $J(x)=\log M_t(s(x))$. Then
\begin{align*}
 J'&=(i/2)\alpha(1+t\alpha'/2),\\
 J''&=-\tfrac14\{\alpha'+(t/2)[(\alpha')^2+\alpha\alpha'']\},\\
 J'''&=-\tfrac i8\{\alpha''+(t/2)[3\alpha'\alpha''+\alpha\alpha''']\}.
\end{align*}
Equations~\eqref{eq:effective-alpha-disc} and
\eqref{eq:effective-alpha-center} imply
$|J'-i\omega|\le1$, $|J''|,|J'''|\le1$ and
$|w'|,|w''|,|w'''|\le1$. A majorant for the first bound is
$\pi/8+2/x+L/(4x)<1$; the remaining bounds follow by inserting the
stated powers of $x^{-1}$ and $x\ge e^L$.

The prime-power series and \eqref{eq:effective-log-moments} at
$\eta=1/2$ give
\[
 |(\log\zeta)'(w)|\le6,\qquad
 |(\log\zeta)''(w)|\le26,\qquad
 |(\log\zeta)'''(w)|\le158.
\]
Thus for a local logarithm $g=\log Q_+$, the chain rule yields
\[
 |g'-i\omega|\le7,\qquad |g''|\le1+26+6=33,
 \qquad |g'''|\le1+158+3\cdot26+6=243.
\]
The resulting carrier errors are
\begin{equation}
\begin{array}{c|c}
 k&\displaystyle\frac{|Q_+^{(k)}/Q-(i\omega)^k|}{\omega^k}\\[2pt]
 \hline
 1&28/L\\
 2&56/L+1312/L^2\\
 3&84/L+3936/L^2+81856/L^3
\end{array}
\label{eq:effective-carrier-table}
\end{equation}
Every entry in the right-hand column is at most $100/L$ for $L\ge1000$.
At order three, for instance, expanding
$(i\omega+d)^3+3(i\omega+d)g''+g'''$ with $|d|\le7$ bounds the
unnormalized discrepancy by $21\omega^2+246\omega+1279$.
The conjugate branch contributes the same normalized bound; the factor
two is already included in $2|Q|\omega^k$. The carrier error at order
zero is exactly zero.

Adding \eqref{eq:effective-cauchy-error} and
\eqref{eq:effective-carrier-table}, every coordinate error is at most
$107520100/L<108000000/L$. For $L\ge10^9$ this is less than $1/8$,
and \eqref{eq:explicit-perturbation} proves
Proposition~\ref{prop:explicit-height}.

The value $X_{\mathrm{exp}}$ is an explicit consequence of the primary
approximation and elementary inequalities. Its exponential need not be
evaluated numerically. The accompanying exact checker verifies 52 rational
budgets and sufficient polynomial inequalities for every $L\ge1000$ by
shifting $L=1000+y$ and checking nonnegative coefficients. It audits the
arithmetic in this proof, rather than replacing its analytic arguments.
The bound supplies no practical finite-strip coverage between 112 and
$X_{\mathrm{exp}}$, and it does not evaluate the height in the stronger
closed-threshold theorem.

\section{Reproducibility and scope of the checks}\label{app:replay}
\begingroup\raggedright

All paths in this appendix are relative to
\path{papers/investigations-2/active/21_quaternion_theta_geometry/}
in the accompanying repository. The analytic proofs above are distinct
from the finite interval decisions. The latter use complete error
bounds, outward rounding, and closed-cell coverage; hashes identify
the evaluated sources and records, but do not substitute for their
mathematical content.

The bounded results can be regenerated using the drivers
\path{numerics/certify_compensated_stationary_time_extension.py},
\path{numerics/certify_correlated_source_height_extension.py}, and
\path{numerics/certify_phase_centered_residual.py}.
The correlated-height guide
\path{numerics/CORRELATED_SOURCE_HEIGHT_README_20261014.md}
gives the parameters for its three bands. The full-time certificate is
\path{numerics/COMPENSATED_STATIONARY_TIME_EXTENSION_20261014.json};
the final centered-band record is
\path{numerics/THETA_CENTERED_RESIDUAL_100_104_20261014.json}.
The independent internal replays are documented in Reviews 6, 9, and 11
in \path{reviews/}. They reproduced the numerical records after removal
of runtime metadata.

For compact exact audits, run the following from \path{numerics/}:
\begin{quote}\small\ttfamily
python3 audit\_correlated\_source\_height\_records.py\par
python3 audit\_theta\_centered\_residual\_record.py\par
python3 check\_octonion\_jet\_identities.py\par
python3 check\_theta\_centered\_residual\_identities.py
\end{quote}
The centered-band auditor additionally requires
\texttt{--record THETA\_CENTERED\_RESIDUAL\_100\_104\_20261014.json}.
Each requires the option \texttt{--output} followed by a desired JSON
path. The first two inspect the retained records and bound source
versions; the latter two check 248 and nine exact algebraic identities,
respectively. The octonionic identities, centered residual, and outer
analytic argument have separate internal assessments in Reviews 10--12.

The diagnostic
\path{numerics/check_outer_cusp_and_phase_obstructions.py}, also run
with \texttt{--output}, requires python-flint and imports
\path{numerics/check_shifted_contour_payments.py}. Its retained record
\path{numerics/OUTER_CUSP_AND_PHASE_OBSTRUCTION_CHECKS_20261014.json}
contains the pointwise counterexamples in Section~\ref{sec:remaining},
twelve staircase boundary checks, exact four-jet perturbation algebra,
and selected complete-source point enclosures. It does not certify a
new rectangle or prove the uniform asymptotic theorem.

The subsequent paired algebra, explicit outer constants, and signed-split
controls are audited by
\path{numerics/check_signed_carrier_pair_identities.py},
\path{numerics/check_explicit_outer_constants.py}, and
\path{numerics/check_signed_curvature_split.py}.
They require only the Python standard library and an \texttt{--output}
path, and check 34, 52, and 14 exact statements respectively. Their
records are
\path{numerics/SIGNED_CARRIER_PAIR_IDENTITY_CHECKS_20261014.json},
\path{numerics/EXPLICIT_OUTER_CONSTANT_CHECKS_20261014.json}, and
\path{numerics/SIGNED_CURVATURE_SPLIT_IDENTITY_CHECKS_20261014.json}.
These checks support algebra and explicit inequalities; the infinite
source and asymptotic arguments require the proofs in the text.

The full-residual branch comparison is recorded in
\path{numerics/STATIONARY_FAVORABLE_CURVATURE_PILOT_20261014.json}.
The proper source-filter and Gaussian pilots are produced by
\path{numerics/certify_predetermined_source_filter_pilot.py} and
\path{numerics/check_gaussian_source_filter_pilot.py}.
Each accepts \texttt{--record}, \texttt{--original-driver},
\texttt{--contour-source}, and \texttt{--output}; the inputs are the
centered-band record, its driver, and
\path{numerics/check_shifted_contour_payments.py}.
The proper-filter run uses python-flint 0.9.0 at 192 bits and retains
all 64 classifications, including 32 failed absorption conditions.
The Gaussian extraction uses exact rational arithmetic and retains all
14 cases. Their records are
\path{numerics/PREDETERMINED_SOURCE_FILTER_PILOT_20261014.json} and
\path{numerics/GAUSSIAN_SOURCE_FILTER_PILOT_20261014.json}.
Reviews 14--16 document the analytic audits and independent replays,
including the eight full-residual branch slabs. Review 17 audits the
earlier manuscript integration; its fresh exact checks are distinct
from the preceding full interval replays.

The height-112 extension and its all-earlier-time comparisons use
\path{numerics/certify_fixed_cusp_band_extension.py} and the two records
\path{numerics/FIXED_CUSP_BAND_104_108_20261014.json} and
\path{numerics/FIXED_CUSP_BAND_108_112_20261014.json}.
The exact auditor \path{numerics/audit_fixed_cusp_band_extension.py}
accepts both paths through repeated \texttt{--record} options and an
\texttt{--output} path. Its 1321 checks retain all 64 closed cells,
128 stationary observations, and 16 genuine branch/time cases.
The source quadratures were independently replayed in Review 19; the
preceding manuscript integration reran the exact auditor, not those quadratures.

The newly integrated arithmetic and collision material is audited by
\path{numerics/check_proportional_delay_identities.py} (20 checks),
\path{numerics/check_phase_controlled_delay_identities.py} (21),
\path{numerics/check_collision_behavior.py} (1072),
\path{numerics/check_hybrid_cutoff_lemma.py} (84), and
\path{numerics/check_remaining_strip_lemma.py} (97).
The proportional checker takes a positional output path; the phase and
strip checkers accept \texttt{--output}. The collision driver writes its
named receipt beside its source; the integration replay calls its pure
check routines without replacing that historical receipt.
Together with the finite-band audit, the preceding integration executed
2615 exact checks. Its receipt is
\path{numerics/MANUSCRIPT_PROGRAM_INTEGRATION_CHECKS_20261015.json}.
These checks verify finite algebra, rational margins, and recorded
interval decisions. They do not prove the imported arithmetic theorems,
the unproved middle-block or weighted-remainder hypotheses, or an
all-height theta sign estimate. Review 24 records the integrated proof
and layout audit, including the conditional deductions from Note 33.

The octonionic simplification and literature integration is documented
in Reviews 25--26. The standard-library checker
\path{numerics/check_octonion_reflected_lift.py}, run with
\texttt{--output}, verifies 325 exact polynomial assertions: the full
eight-dimensional reflected operator, physical trace, heat and jet
identities, and the planar bilinear reductions. The original 248-check
octonion driver was also replayed. The independent moment and multiplier
checks for the correlation formulas are supplied by
\path{numerics/check_theta_correlation_identities.py}; its scope and
the consolidated receipt are recorded in Review 26. These finite
controls support the general proofs above; they add no theta sign
certificate or interval-quadrature coverage.

The present integration incorporates the subsequent Poisson and signed-form
continuations (Notes 38--53 and Reviews 27--31). The fresh integration
replay executes the standard-library drivers
\path{numerics/check_convex_dyadic_weighted_remainder.py},
\path{numerics/check_poisson_mode_cutoff_obstructions.py}, and
\path{numerics/check_explicit_weighted_budget.py}, each with
\texttt{--output}. Their 129, 260, and 50 exact assertions pass.
These 439 assertions check finite algebra and rational budgets; the
analytic proofs of the infinite-domain claims are given above.
The six fixed-theta root certificates retain the complete-source
quadrature provenance of Note 40 and Review 27; they were not re-evaluated
in this manuscript integration. The newly integrated signed-form,
primitive, and smooth-control results are analytic, with independent
internal cross-review, and introduce no numerical sign coverage.
Review 32 and the current \path{manuscript/BUILD_RECORD.json} record the
source/export identity, compilation, reference checks, and layout review.

The main remaining obligation is the one-sided inner-region bound
\eqref{eq:final-target}, with losses locally bounded through candidate
positive-time collisions. Equations~\eqref{eq:corr-stationary-target}
and~\eqref{eq:corr-hybrid-transform} organize that same obligation
using complete source correlations. Improving the explicit height and
connecting the finite cover remain separate quantitative work.
Neither a fixed-time asymptotic nor a ratio divided by unknown physical
energy supplies the missing bound.

\par\endgroup

\section*{Acknowledgment}
This manuscript was drafted for Edward Baker with substantial assistance
from OpenAI language models. The present integration used Codex with the
recorded model \texttt{gpt-6-astra} and reasoning effort \texttt{ultra};
earlier research and numerical certificates also used
\texttt{gpt-6.1-sol}. Edward Baker's unpublished quaternion-form working
notes of January and February 2014 motivated the signed differential-form
construction. Those notes were circulated with mistakes and incompleteness
acknowledged; no general nonabelian boundary formula from them is assumed. LLM assistance included derivations, exposition,
code, and internal checks. The accompanying reviews distinguish
analytic arguments, validated interval calculations, exact algebra
checks, and their remaining limitations. These internal checks are not
external specialist refereeing or proof-assistant formalization.

\begin{thebibliography}{99}
\raggedright
\bibitem{debruijn}
N. G. de Bruijn,
\emph{The roots of trigonometric integrals},
Duke Mathematical Journal \textbf{17} (1950), 197--226.
\url{https://doi.org/10.1215/S0012-7094-50-01720-0}.

\bibitem{newman}
C. M. Newman,
\emph{Fourier transforms with only real zeros},
Proceedings of the American Mathematical Society \textbf{61} (1976),
245--251.
\url{https://doi.org/10.1090/S0002-9939-1976-0434982-5}.

\bibitem{polymath}
D. H. J. Polymath,
\emph{Effective approximation of heat flow evolution of the Riemann
$\xi$ function, and a new upper bound for the de Bruijn--Newman constant},
arXiv:1904.12438v2 (2019).
\url{https://arxiv.org/abs/1904.12438}.

\bibitem{robert}
O. Robert,
\emph{On van der Corput's $k$-th derivative test for exponential sums},
Indagationes Mathematicae \textbf{27} (2016), 559--589.
\url{https://doi.org/10.1016/j.indag.2015.11.009}.
Author manuscript:
\url{https://perso.univ-st-etienne.fr/rool6510/robert-2015-indag.pdf}.

\bibitem{sudbery}
A. Sudbery,
\emph{Quaternionic analysis},
Mathematical Proceedings of the Cambridge Philosophical Society
\textbf{85} (1979), 199--225.
\url{https://doi.org/10.1017/S0305004100055638}.
% Primary references checked 14 October 2026.
% Existing manuscript already has the sudbery key.
\bibitem{baez}
J. C. Baez,
\emph{The octonions},
Bulletin of the American Mathematical Society \textbf{39} (2002), 145--205.
\url{https://arxiv.org/abs/math/0105155}.
The Cayley--Dickson convention is in Section~2.2;
\url{https://math.ucr.edu/home/baez/octonions/node5.html}.

\bibitem{kauhanenorelma}
J. Kauhanen and H. Orelma,
\emph{Cauchy--Riemann operators in octonionic analysis},
Advances in Applied Clifford Algebras \textbf{28} (2018), article 1.
\url{https://doi.org/10.1007/s00006-018-0826-2};
\url{https://arxiv.org/abs/1701.08698}.

\bibitem{bourgain14}
J. Bourgain,
\emph{Decoupling, exponential sums and the Riemann zeta function},
arXiv:1408.5794v2, especially Theorem 6 and its subinterval discussion.
\url{https://arxiv.org/abs/1408.5794v2}.

\bibitem{trudgianyang23}
T. S. Trudgian and A. Yang,
\emph{Toward optimal exponent pairs},
arXiv:2306.05599v3, especially Section 1.1, Lemma 1.1, and equation (1.10).
\url{https://arxiv.org/abs/2306.05599v3}.

\bibitem{ruiz}
J. M. Ruiz,
\emph{Basic properties of real analytic and semianalytic germs},
Publications mathematiques et informatique de Rennes (1986), no. 4,
29--44, especially (7.2), p. 40.
\url{https://www.numdam.org/item/PSMIR_1986___4_29_0.pdf}.

\bibitem{taotrudgianyang25}
T. Tao, T. Trudgian, and A. Yang,
\emph{New exponent pairs, zero density estimates, and zero additive
energy estimates: a systematic approach}, arXiv:2501.16779v1 (2025),
Table 1. \url{https://arxiv.org/abs/2501.16779v1}.

\bibitem{harveylawson}
R. Harvey and H. B. Lawson, Jr.,
\emph{Calibrated geometries}, Acta Mathematica \textbf{148} (1982), 47--157,
Chapter IV, Theorem 1.6.
\url{https://doi.org/10.1007/BF02392726}.

\bibitem{karigiannis}
S. Karigiannis, \emph{Introduction to $G_2$ geometry},
Fields Institute Communications \textbf{84} (2020), 3--50,
especially Corollary 3.52.
\url{https://arxiv.org/abs/1909.09717}.

\bibitem{schraymanogue}
J. Schray and C. A. Manogue,
\emph{Octonionic representations of Clifford algebras and triality},
Foundations of Physics \textbf{26} (1996), 17--70.
\url{https://arxiv.org/abs/hep-th/9407179}.

\bibitem{kauhanenorelma19}
J. Kauhanen and H. Orelma, \emph{On the structure of octonion regular
functions}, Advances in Applied Clifford Algebras \textbf{29} (2019),
article 77. \url{https://doi.org/10.1007/s00006-019-0983-y}.

\bibitem{lipeng}
X. Li and L. Peng, \emph{The Cauchy integral formulas on the octonions},
Bulletin of the Belgian Mathematical Society \textbf{9} (2002), 47--64.
\url{https://doi.org/10.36045/bbms/1102715140}.

\bibitem{colombokraussharsabadini}
F. Colombo, R. S. Krau{\ss}har, and I. Sabadini,
\emph{Octonionic monogenic and slice monogenic Hardy and Bergman spaces},
Forum Mathematicum \textbf{36} (2024), no. 4.
\url{https://doi.org/10.1515/forum-2023-0039};
\url{https://arxiv.org/abs/2311.07691}.

\bibitem{constaleskrausshar}
D. Constales and R. S. Krau{\ss}har,
\emph{Octonionic Kerzman--Stein operators} (2021).
\url{https://doi.org/10.1007/s11785-021-01152-0};
\url{https://arxiv.org/abs/2012.11925}.

\bibitem{teodorescu24}
R. S. Krau{\ss}har, M. Ferreira, N. Vieira, and M. M. Rodrigues,
\emph{The Teodorescu and the $\Pi$-operator in octonionic analysis and
some applications}, Journal of Geometry and Physics \textbf{206}
(2024), 105328. \url{https://doi.org/10.1016/j.geomphys.2024.105328}.

\bibitem{gentilistruppa}
G. Gentili and D. C. Struppa,
\emph{Regular functions on the space of Cayley numbers},
Rocky Mountain Journal of Mathematics \textbf{40} (2010), 225--241.
\url{https://doi.org/10.1216/RMJ-2010-40-1-225}.

\bibitem{ghiloni}
R. Ghiloni, \emph{Slice Fueter-regular functions},
Journal of Geometric Analysis \textbf{31} (2021), 11988--12033.
\url{https://doi.org/10.1007/s12220-021-00709-x}.

\bibitem{xusabadini26}
Z. Xu and I. Sabadini,
\emph{Generalized partial-slice monogenic functions: the octonionic case},
Transactions of the American Mathematical Society \textbf{379} (2026),
3257--3291. \url{https://doi.org/10.1090/tran/9646};
\url{https://arxiv.org/abs/2503.12409v4}.

\bibitem{fuetermapping}
F. Colombo, I. Sabadini, and F. Sommen,
\emph{The Fueter mapping theorem in integral form and the F-functional
calculus}, Mathematical Methods in the Applied Sciences \textbf{33}
(2010), 2050--2066. \url{https://doi.org/10.1002/mma.1315}.

\bibitem{feehan}
P. M. N. Feehan, \emph{A Kato--Yau inequality for harmonic spinors and
decay estimate for eigenspinors}, Journal of Geometric Analysis
\textbf{11} (2001), 469--489.
\url{https://arxiv.org/abs/math/9903021}.

\bibitem{borrelliwu26}
W. Borrelli and R. Wu, \emph{On the nodal set of solutions to Dirac
equations}, Journal of Geometric Analysis \textbf{36} (2026), article 195.
\url{https://doi.org/10.1007/s12220-026-02443-8}.

\bibitem{csordas}
G. Csordas, \emph{Fourier transforms, positive definite kernels and the
Riemann $\xi$-function}, arXiv:1309.0055v2 (2014),
especially Theorems 3.5 and 3.7.
\url{https://arxiv.org/abs/1309.0055v2}.

\bibitem{dimitrovxu}
D. K. Dimitrov and Y. Xu, \emph{Wronskians of Fourier and Laplace
transforms}, Transactions of the American Mathematical Society
\textbf{372} (2019), 4107--4125.
\url{https://arxiv.org/abs/1606.05011}.

\bibitem{elkiesgross}
N. D. Elkies and B. H. Gross, \emph{Embeddings into the integral
octonions}, Pacific Journal of Mathematics \textbf{181} (1997),
special issue, 147--161.
\url{https://msp.org/pjm/1997/181-3/pjm-v181-n3-p06-s.pdf}.

\bibitem{kim93}
H. H. Kim, \emph{Exceptional modular form of weight 4 on an exceptional
domain contained in $\mathbb C^{27}$}, Revista Matem\'atica Iberoamericana
\textbf{9} (1993), 139--200. \url{https://doi.org/10.4171/RMI/134}.

\bibitem{pollackg2}
A. Pollack, \emph{Exceptional theta functions and arithmeticity of
modular forms on $G_2$}, arXiv:2211.05280v2 (2023).
\url{https://arxiv.org/abs/2211.05280v2}.

\bibitem{pollackspin8}
A. Pollack, \emph{Exceptional Siegel--Weil theorems for compact
$\operatorname{Spin}_8$}, arXiv:2308.09100 (2023).
\url{https://arxiv.org/abs/2308.09100}.

\bibitem{shan26}
Y. Shan, \emph{The global exceptional theta correspondence for
$F_4\times\operatorname{PGL}_2$}, arXiv:2501.19101v2,
revised 12 October 2026, especially Theorems B and E.
\url{https://arxiv.org/abs/2501.19101v2}.

\bibitem{viazovska}
M. S. Viazovska, \emph{The sphere packing problem in dimension 8},
Annals of Mathematics \textbf{185} (2017), 991--1015.
\url{https://doi.org/10.4007/annals.2017.185.3.7}.

\bibitem{romik}
D. Romik, \emph{On Viazovska's modular form inequalities},
Proceedings of the National Academy of Sciences \textbf{120} (2023),
e2304891120. \url{https://doi.org/10.1073/pnas.2304891120}.

\bibitem{lee24}
S. Lee, \emph{Algebraic proof of modular form inequalities for optimal
sphere packings}, arXiv:2406.14659 (2024).
\url{https://arxiv.org/abs/2406.14659}.

\bibitem{lee26polynomial}
S. Lee, \emph{Inequalities involving polynomials and quasimodular forms},
arXiv:2602.10536 (2026), preprint.
\url{https://arxiv.org/abs/2602.10536}.

\bibitem{lee26positive}
S. Lee, \emph{Positive quasimodular forms and the sign uncertainty
principle}, arXiv:2608.15415v2 (2026), preprint.
\url{https://arxiv.org/abs/2608.15415v2}.

\bibitem{conreyli}
J. B. Conrey and X.-J. Li, \emph{A note on some positivity conditions
related to zeta- and $L$-functions}, arXiv:math/9812166 (1998).
\url{https://arxiv.org/abs/math/9812166}.

\end{thebibliography}
\end{document}
